% concatenated from main.tex
\documentclass[reqno,11pt]{amsart}
\usepackage{presets}

\selectlanguage{english}
\begin{document}

%
    \title[Spectral rigidity of Sinai billiards]{A CAT(0)-approach to the marked length spectral rigidity of Sinai billiards}
    
    \author{Douglas Finamore}
    \address{Instituto de Matem\'atica, Estat\'istica e Computa\c{c}\~ao Cient\'ifica - UNICAMP\\
    	Rua S\'ergio Buarque de Holanda, 651 | 13083-859, Campinas, SP, Brasil.}
    \email{douglas.finamore@alumni.usp.br}

     
    \author{Martin Leguil}
    \address{Centre de Math\'ematiques Laurent Schwartz - \'Ecole Polytechnique\\
    	Route de Saclay, 91128 Palaiseau Cedex, France.}
    \email{martin.leguil@polytechnique.edu}
%

\begin{abstract} 
We study the spectral rigidity problem for Sinai billiards with finite horizon, specifically asking whether the geometry of the billiard table can be recovered from the lengths of its (marked) periodic trajectories.
To address this, we introduce an enriched marked length spectrum $\mathcal{EL}$ and prove that two Sinai billiards sharing the same $\mathcal{EL}$ must be isometric.
Our approach involves approximating the billiard flow using geodesic flows on smooth Riemannian surfaces. In the limit, these flows converge to $\mathrm{CAT}(0)$ spaces, which encode both the lengths of periodic orbits and the geometry of the boundary. We adapt Otal’s original method\textemdash developed for marked length spectrum rigidity in negatively curved surfaces\textemdash to this new setting. Here, the lack of curvature control is offset by metric comparison estimates.
By integrating the analysis of geodesic flows with perturbative techniques for periodic orbits, we establish a rigidity theorem for Sinai billiards with finite horizon.
These results extend the classical theory of marked length spectrum rigidity beyond the Riemannian setting, demonstrating that\textemdash even in discontinuous dynamical systems\textemdash geometric information is rigidly encoded in spectral data. 
\end{abstract}
\maketitle

\tableofcontents

\section{Introduction}

We consider dynamical systems with hyperbolic behaviour arising from geometry, such as geodesic flows on negatively curved Riemannian manifolds or dispersing billiards. A central question is how to parameterise the moduli space of such systems. Specifically, given a Riemannian manifold equipped with a negatively curved metric, one seeks to describe the set of all such metrics up to the action of the diffeomorphism group. Similarly, for billiards, the goal is to understand the set of isometries within a fixed diffeomorphism class of dispersing planar billiards.
These systems typically exhibit a dense set of periodic orbits (restricted to the non-wandering set in the non-compact case), making periodic data a natural choice for moduli. The information recorded for periodic orbits can be categorized as follows:
\begin{itemize}
    \item topological or symbolic data: on negatively curved manifolds, and more generally, for metrics with Anosov geodesic flow, it is well-known (see e.g.~\cite{klingenberg_riemannian_1974}) that there is a unique closed geodesic in each free homotopy class. For dispersing billiards, such as open dispersing billiards or Sinai billiards, the symbolic data naturally assigned to each periodic orbit is the sequence of scatterers encountered by the trajectory;
    \item length data: the \emph{length spectrum} is the collection of all lengths (or periods) of closed orbits;  
    \item Lyapunov exponents: each periodic orbit in a hyperbolic system is associated with exponents measuring the exponential divergence rate of nearby orbits. In two dimensions, this reduces to a single exponent.
\end{itemize}
A natural question that arises is: To what extent do periodic orbits characterise the system?
A dynamical analogue of Mark Kac’s famous question\textemdash ``Can one hear the shape of a drum?''\textemdash is the following inverse problem: Does the length spectrum encode the system’s geometry?
Unfortunately, this problem is typically intractable due to the lack of structure in the length spectrum. Even if two systems share the same length spectrum, one cannot \textit{a priori} rule out the possibility that a given length corresponds to periodic orbits with vastly different behaviour in each system. Matching data between such systems is therefore generally a hard task. 
To impose more structure on the spectrum, one usually augments it with symbolic data, yielding a \emph{marked length spectrum}. Informally, this can be viewed as a (set-valued) function
$$
\text{symbolic data } \omega \quad \mapsto\quad \{\text{lengths of all periodic orbits with symbol }\omega\}. 
$$
In the hyperbolic setting, the expansiveness of the dynamics typically reduces the set on the right-hand side to a single positive number. For certain systems\textemdash such as geodesic flows on negatively curved Riemannian manifolds or open dispersing billiards satisfying the no-eclipse condition\textemdash the dynamics is structurally stable. Consequently, the symbolic data does not depend on the choice of negatively curved metric for a given Riemannian manifold, nor on the specific geometry of the open dispersing billiard.
This defines a natural marked length spectrum, allowing us to rephrase the previous question as follows:
Does the marked length spectrum fully determine the system up to isometry?

The question of marked length spectral determination for negatively curved metrics on surfaces was answered affirmatively and simultaneously by Otal~\cite{otal_spectre_1990} and Croke~\cite{croke1990rigidity} in their celebrated work. Significant progress has recently been made in generalizing these results in several directions.
In higher dimensions, Guillarmou and Lefeuvre~\cite{guillarmou_marked_2019} proved local marked length spectral determination for metrics of non-positive curvature with Anosov geodesic flow. In the same work, they also provided new stability estimates, quantifying the extent to which the marked length spectrum controls the distance between metrics.
Even more recently, for surfaces of genus at least $2$, Guillarmou, Lefeuvre, and Paternain~\cite{guillarmou2025marked} were able to drop the negative curvature assumption and establish global marked length spectral determination. Specifically, they showed that two metrics with Anosov geodesic flows on a given surface are isometric via an isometry isotopic to the identity if they share the same marked length spectrum.

A related framework is that of flat metrics, more precisely, non-positively curved Euclidean cone metrics, where curvature is concentrated at conical singularities. Banković and Leininger~\cite{bankovic_marked-length-spectral_2017} showed that two flat metrics on a closed, oriented surface that assign the same lengths to all closed curves differ by an isometry isotopic to the identity.
It is also worth noting that the question of marked length spectral rigidity has been investigated in very different contexts. For instance, recent work by Abbondandolo and Mazzucchelli~\cite{Abb_Mazz} addresses the marked length spectral determination of metrics on the $2$-sphere $\mathbb{S}^2$ with an $\mathbb{S}^1$-action and a unique equator. Here, the dynamics is integrable, which sets it apart from the hyperbolic geodesic flows discussed earlier.

Spectral rigidity of billiards has been studied for various classes of domains. De Simoi, Kaloshin and Wei~\cite{desimoikal} showed that convex domains with $\mathbb{Z}_2$-symmetry and that are sufficiently close to a circle exhibit marked length spectral rigidity. In~\cite{de2023marked}, it is shown that open dispersing billiards with analytic boundaries satisfying the no-eclipse condition and with $\mathbb{Z}_2 \times \mathbb{Z}_2$-symmetry can be recovered from their marked length spectral data. The dynamics of such billiards is hyperbolic, aligning them more closely with Anosov geodesic flows. These billiards are structurally stable and admit natural symbolic dynamics. However, the non-wandering dynamics occurs only on a Cantor set, and analyticity is required to globally determine the domain.\\

In this paper, we investigate the case of Sinai billiards, obtained by removing from the $2$-torus $\mathbb{T}^2$ a finite number of strongly convex scatterers. While the dynamics of Sinai billiards shares many features with Anosov geodesic flows on surfaces\textemdash such as hyperbolicity, ergodicity, and exponential mixing, see e.g.~\cite{chernov_chaotic_2006, baladi2018exponential}\textemdash it also introduces several additional challenges:
\begin{itemize}
\item the billiard map has a dense set of singularities, and the billiard flow is not Anosov;
\item the system is not structurally stable: perturbing the geometry of the boundary can create or destroy periodic orbits.
\end{itemize}
A symbolic coding for Sinai billiards does arise naturally: by considering a $\mathbb{Z}^2$-periodic lift of the billiard table to $\mathbb{R}^2$ and labelling each scatterer, each orbit can be assigned a symbolic sequence consisting of the labels of the obstacles it successively encounters.
However, due to the lack of structural stability, the question of marked length spectral determination for Sinai billiards may seem less canonical than for Anosov metrics on surfaces. The marking is not \textit{a priori} stable under perturbation, and the symbolic data itself must encode substantial geometric information. In particular, if we discard the length data and collect all symbolic data corresponding to sequences realized by (periodic) billiard orbits into a \emph{symbolic spectrum}, we can even pose the following stronger question:
\begin{question}
    Does the symbolic spectrum determine a Sinai billiard up to isometry? 
\end{question}
While we are unable to answer this question for Sinai billiards, we note that several works address similar questions for flat billiards (see, e.g.,~\cite{duchin2021you, calderon2018hear}).
In the present work, rather than focusing on the marked length spectrum corresponding to periodic orbits of a Sinai billiard, we take a different approach. Specifically, we build on a construction\textemdash suggested by Birkhoff and Arnold and detailed in the work of Kourganoff~\cite{kourganoff_anosov_2016}\textemdash which demonstrates that the billiard flow on a Sinai billiard can, in a certain sense, be approximated by a one-parameter family of Anosov geodesic flows on a surface lifting the table. This allows us to define an \emph{enriched length spectrum}, obtained as the ``limit'' of the marked length spectra of the approximating geodesic flows. Although this construction introduces periodic data that do not correspond to actual billiard orbits, it renders the marking more natural by leveraging the symbolic data shared across the family of approximating geodesic flows.

Let us quickly recall the construction in~\cite{kourganoff_anosov_2016}. Fix a finite-horizon Sinai billiard $\mathcal{D}\subset \mathbb{T}^2$. Let \( E= \mathbb{T}^2\times \R \), and denote by \( p\colon E \to \mathbb{T}^2 \) the canonical projection onto the first two coordinates. We consider a smooth closed immersed surface \( K \) in \( E \) such that \( p(K) = \mathcal{D} \), and such that  $K \setminus p^{-1}(\partial\mathcal{D})$ has two connected components each of which is diffeomorphically mapped to \( \mathrm{int}(\mathcal{D}) \) under \( p \). 
We interpret \( K \) as a surface obtained from \( \mathcal{D} \) by gluing different copies of \( \mathcal{D} \) along their respective boundary components.

\begin{figure}[!h]
    \centering
    \includegraphics[width=.8\textwidth]{liftsurface.png}
    \caption{A surface $K$ lifting the table $\mathcal{D}$ and a deformation $K_\epsilon$ \\ (Credits: Mickaël Kourganoff, Jos Leys).\label{liftsurface_png}} 
\end{figure}

\par For \( \epsilon \in (0,1] \), we let \( \alpha_\epsilon \) be the contraction given by \( (x,y,z) \mapsto (x, y, \epsilon z) \) and write \( g_\epsilon \) to denote the metric induced on \( K_\epsilon := \alpha_\epsilon(K) \) by the Euclidean metric.
Pulling back via \( \alpha_\epsilon \), we will consider \( \{g_\epsilon\} \) as a family of metrics on the surface \( K \). 
We denote by \( \mathbf{g}_\epsilon \) the geodesic flow of the surface \( (K, g_\epsilon) \).
Then \( p \) induces a smooth mapping
\begin{equation*} 
\begin{split}
    \pi_K\colon T^1K^\circ &\to T^1 \mathcal{D} \\
    (x,v) &\mapsto \left(p(x), \frac{dp_x(v)}{\|dp_x(v)\|}\right).
\end{split}
\end{equation*} 
Let us denote by \( \mathcal{C}:= \mathrm{Conj}(\pi_1(K)) \) the set of free homotopy classes on \( K \). Note that $\mathcal{C}$ is independent of the choice of surface $K$ lifting $\mathcal{D}$ that satisfies our requirements. By Kourganoff~\cite{kourganoff_anosov_2016}, for $\epsilon_0>0$ small enough, the geodesic flows $\{\mathbf{g}_\epsilon\}_{0<\epsilon<\epsilon_0}$ associated to the family of metrics $\{g_\epsilon\}_{0<\epsilon<\epsilon_0}$ are all Anosov. Therefore, for $0<\epsilon<\epsilon_0$, there exists a unique closed geodesic \( \gamma_\epsilon^\alpha\) of \( (K, g_\epsilon) \) in each free homotopy class \( \alpha \in \mathcal{C}(K) \). We will see in Lemma~\ref{lem:unique_minimising_geodesic} that the 
family $\{p(\gamma_\epsilon^\alpha)\}_{0<\epsilon<\epsilon_0}$ converges in the Hausdorff topology to some closed curve \( \gamma^\alpha_0 \) in $\mathcal{D}$ as $\epsilon\to 0$, and their respective lengths $\{l_\epsilon(\gamma_\epsilon^\alpha)\}_{0<\epsilon<\epsilon_0}$ also converge to a limiting length $l_0(\gamma_0^\alpha)$. We can then define the~\emph{enriched marked length spectrum} of \( \mathcal{D} \) as the function
    \begin{align*}
        \mathcal{EL}_{\mathcal{D}}\colon \mathcal{C} &\to \R \\
        \alpha &\mapsto l_0(\gamma_0^\alpha). 
    \end{align*} 
\begin{remark}
For a ``generic'' Sinai billiard with finite horizon (see Proposition~\ref{prop_good_tables}), every periodic orbit of the billiard is realized as $\gamma_0^\alpha$ for some $\alpha \in \mathcal{C}$. Consequently, $\mathcal{EL}_\mathcal{D}$ can be interpreted as a natural completion of the marked length spectrum of $\mathcal{D}$, which would be the restriction of $\mathcal{EL}_\mathcal{D}$ to those $\alpha \in \mathcal{C}$ associated with genuine billiard orbits.
\end{remark}
\begin{remark}
As we shall see (e.g. in Section~\ref{sec:geodesic_approximations}), there are multiple ways to rigorously define the ``limit'' of the family of Anosov geodesic flows $\{\mathbf{g}_\epsilon\}_{0 < \epsilon < \epsilon_0}$ as $\epsilon \to 0$, all of which are equivalent. Although the resulting limit object is no longer an Anosov flow, it retains several key features inherited from the limiting $\mathrm{CAT(0)}$ spaces.
We also emphasise that the enriched spectrum $\mathcal{EL}_\mathcal{D}$ is independent of the specific choice of the surface $K$ lifting $\mathcal{D}$; in particular, it is intrinsic to the Sinai billiard table $\mathcal{D}$ (see e.g. Appendix~\ref{sec_billiards_cycles_lf}). 
\end{remark}
Let $\mathcal{D}_1,\mathcal{D}_2\subset \mathbb{T}^2$ be two Sinai billiards with finite horizon that are diffeomorphic to each other. We can choose surfaces $K_1,K_2$ lifting $\mathcal{D}_1,\mathcal{D}_2$ which are diffeomorphic to each other; in particular, the sets  of free homotopic classes on $K_1,K_2$ are naturally identified.  
Our main result is: 
\begin{theoalph} 
Let $\mathcal{D}_1,\mathcal{D}_2\subset \mathbb{T}^2$ be two Sinai billiards with finite horizon that are diffeomorphic to each other. If the billiards $\mathcal{D}_1,\mathcal{D}_2$ share the same enriched marked length spectrum, i.e., $\mathcal{EL}_{\mathcal{D}_1}=\mathcal{EL}_{\mathcal{D}_2}$, then $\mathcal{D}_1$ and $\mathcal{D}_2$ are actually isometric. 
\end{theoalph}
Natural questions then arise.
\begin{question}
Does partial information from the enriched marked length spectrum suffice to recover this conclusion? In particular, to what extent does the set of periodic billiard orbits encode the geometry of the billiard $\mathcal{D}$?
\end{question}

\begin{question}
For convex billiards, the Poisson relation\textemdash due to Andersson and Melrose~\cite{andersson1977propagation}\textemdash connects the singular support of the wave trace (defined via the Laplace spectrum) to the length spectrum. By analogy, in the case of Sinai billiards, is it possible to demonstrate that the additional length data, which do not originate from genuine billiard orbits, are also encoded in some Laplace spectral data?
\end{question}


\subsection*{Acknowledgements}
We are grateful to the LESET/MATH-AMSUD project, coordinated by Jérôme Buzzi, Pablo D. Carrasco, and Radu Saghin, for making this collaboration possible. 
We thank the Centre de Mathématiques Laurent Schwartz (CMLS) at École Polytechnique for its hospitality during the development of this project. 
We also thank Andrey Gogolev for his valuable insights and discussions related to this work, and Mickaël Kourganoff for permission to reproduce the figures from~\cite{kourganoff_anosov_2016}. 

D.F. was supported by CAPES Grant No. 88887.898617/2023-00 as part of the LESET/MATH-AMSUD project. 
M.L. was supported by the ANR AAPG 2021 PRC CoSyDy (Conformally symplectic dynamics, beyond symplectic dynamics; Grant No. ANR-CE40-0014) and by the ANR JCJC PADAWAN (Parabolic dynamics, bifurcations, and wandering domains; Grant No. ANR-21-CE40-0012). 

\section{Preliminaries}\label{sec_symbolics_prelim}
\par We begin by introducing some established facts and notation.
All the statements in this section can be found in \cite{chernov_chaotic_2006}, if only with slightly different notation. 
Let \( \mathcal{D}\subset \mathbb{T}^2 \) be a Sinai billiard table with finite horizon whose boundary \( \partial \mathcal{D} \) is of class \( C^r \),  \( r \in \mathbb{N}_{\geq 3}\).
More specifically, we have 
\[
    \partial\mathcal{D} = \Gamma_1 \cup \Gamma_2 \cup \cdots \cup \Gamma_m,
\]    
where \( \Gamma_i \) is a $C^r$ curve bounding a strongly convex open set \( \mathcal{O}_i \subset \mathbb{T}^2 \), called an obstacle, and \( \mathcal{D} \coloneqq \mathbb{T}^2\setminus \cup_{i=1}^m \mathcal{O}_i \).

\par For each \( i \in \{1,\cdots,m\} \), we denote by \( \ell_i>0 \) the perimeter of \( \mathcal{O}_i \), and let \( \gamma_i\colon \mathbb{T}_i\to \T^2 \) be an anticlockwise parameterisation of \( \Gamma_i \) by arc-length, where \( \mathbb{T}_i\coloneqq\R/\ell_i \Z \).
We set \( \mathbb{T}_\mathcal{D}\coloneqq\cup_{i=1}^m \mathbb{T}_i \), and let \( \gamma\colon \mathbb{T}_\mathcal{D}\to  \T^2 \), \( \T_i \ni s\mapsto \gamma_i(s) \). 
For each \( s \in \T_\mathcal{D} \), we denote by \( n_s \) the unit outward normal vector at the point \( \gamma(s)\in \partial \mathcal{D} \), and by \( \mathcal{K}(s) \) the curvature at this point; following the convention in \cite{chernov_chaotic_2006}, we assume that \( \mathcal{K}(s)> 0 \). 
\par We let \( \mathcal{M}\coloneqq\T_\mathcal{D} \times [-\frac{\pi}{2},\frac{\pi}{2}] \), and denote by \( f\colon\mathcal{M} \to\mathcal{M} \), \( (s,\varphi)\mapsto (s',\varphi') \) the billiard map within \( \mathcal{D} \), where \( s \) corresponds to the position on the boundary, and \( \varphi \) is the oriented angle from the inward normal vector \( -n_s \) to the velocity vector at \( \gamma(s) \).
\par Moreover, \( \mathcal{S}_0 \coloneqq \partial \mathcal{M} = \{(s,\varphi); |\varphi| = \pi/2\} \) is the set of \emph{grazing} collisions of \( f \). 
Set
\[
    \mathcal{S}_{n+1} \coloneqq \mathcal{S}_n \cup f^{-1}(\mathcal{S}_n) \text{ and }  \mathcal{S}_{-(n+1)} \coloneqq \mathcal{S}_{-n} \cup f(\mathcal{S}_{-n}).
\] 
Then, for every \( n \in \N \), \( \mathcal{S}_{n+1}\), resp. \( \mathcal{S}_{-(n+1)}\) is the singularity set of the map \( f^{n+1} \), resp. \( f^{-(n+1)} \), while these maps are  of class \( C^{r-1} \) at \( \mathcal{M}\setminus(\mathcal{S}_{n+1} \cup \mathcal{S}_{-(n+1)}) \).
The points in \( \mathcal{S} \coloneqq \cup_{i \in \Z} \mathcal{S}_i \) are called the \emph{singular points} of the billiard map, and \( \mathcal{M}^\mathrm{reg} \coloneqq \mathcal{M}\setminus\mathcal{S} \) is the \emph{regular set} of \( f \).
The singular set has empty interior, but is dense in \(\mathcal{M}\), while the regular set has full Lebesgue measure, but is badly disconnected. 
\par Finally, for each \( (s, \varphi) \in \mathcal{M} \), its free-flight time is the function 
\[
    \tau(s, \varphi) = \|\gamma(\mathrm{Pr}_1(f(s, \varphi))) - \gamma(s)\|,
\]
where we denote by \( \| \cdot \| \) the Euclidean distance in \( \mathbb{T}^2 \) and by \(\mathrm{Pr}_1\colon \mathcal{M} \to \partial\mathcal{D}\) the projection onto the first coordinate.
The finite horizon hypothesis means that there are positive constants \( \tau_\mathrm{min}>0 \) and \( \tau_\mathrm{max}>0 \) such that
\( \tau_\mathrm{min} \leq \tau(x) \leq \tau_\mathrm{max} \) for every \( x \in \mathcal{M} \).
The suspension flow of \( f \) with roof function \( \tau \) is called the billiard flow \( \Psi \) of \( \mathcal{D} \). 
Its phase space is denoted by \( \Omega \), and the usual coordinates are \((x, y, \omega)\), where \((x,y)\) are coordinates on \(\mathcal{D}\) and \(\omega\) is the angle between the velocity of the trajectory and the \(x\)-axis.
The saturations of \(\cS\) and \(\mathcal{M}^{\mathrm{reg}}\) under the flow form the singular and regular set of \(\Omega\), denoted by \(\cS\) and \(\Omega^{\mathrm{reg}}\).
It is known that the billiard flow preserves the Lebesgue measure \(dx\wedge dy\wedge d\omega\) and the contact form \(\lambda = \cos\omega dx + \sin\omega dy\) \cite{chernov_chaotic_2006}.
Indeed, the flow is generated by the Reeb field of such form.
\par We will adopt, instead, the Jacobi coordinates \((\eta, \xi, \omega)\), introduced in \cite{wojtkowski_two_1994}, where
\begin{align*}
    \eta &= x\cos\omega + y\sin\omega, \\
    \xi &= x\sin\omega - y\cos\omega.
\end{align*}
In particular, while in the usual coordinates the action of the flow between collisions is \(\Psi^t(x, y, \omega) = (x + t\cos\omega, y+t\sin\omega, \omega)\), in Jacobi coordinates this action is simply translation on the first coordinate: 
\[
    \Psi^t(\eta, \xi, \omega) = (\eta + t, \xi, \omega).
\]
Inverting the change of coordinates gives \(x = \eta\cos\omega  + \xi\sin\omega \) and \( y = \eta\sin\omega  - \xi\cos\omega \).
Hence the forms
\[
    d\eta + \xi d\omega = \cos\omega dx + \sin\omega dy = \lambda
\]
and
\[
   \lambda \wedge d\lambda =  d\eta\wedge d\xi\wedge d\omega = - dx\wedge dy\wedge d\omega
\]
are both preserved by the flow.
In Jacobi coordinates, the Reeb field of \(\lambda\) is exactly the coordinate field \(\partial_\eta\), which correspond to the displacement of a particle in the direction of the flow.
\par The map \(f\) and the flow \(\Psi\) are hyperbolic: at every regular point there are two distinct Lyapunov exponents (and a zero exponent in the direction \(\partial_\eta\), in the case of the flow). 
This gives a decomposition 
\[
    T_x\Omega = T^c_x\Omega \oplus T^\perp_x\Omega = \R\partial_\eta \oplus E^u(x) \oplus E^s(x),
\]
of the tangent space of every regular point, where \(E^s(x), E^u(x)\) are the stable and unstable bundles, respectively.
As usual, \(E^s(x)\), resp. $E^u(x)$, contracts, resp. expands exponentially under the flow as \(t \to \infty\), and is locally integrable to \(C^{r-2}\) curves, called the strong stable (resp. unstable) local manifolds \(\Ws{x}\) (resp. \(\Wu{x}\)).
The end points of the manifolds \(\cW_{\mathrm{loc}}^*(x)\) (not included in the manifolds) are always singular.
These manifolds are not uniform in length, and get smaller as one approaches the singular set.
In particular, every open set of \(\Omega\) contains points whose unstable manifolds are arbitrarily small, and therefore points in \(\Omega\) do not have product neighbourhoods.
The saturations of \(\Ws{x}\) and \(\Wu{x}\) with respect to the flow are the weak, or central, stable  and unstable manifolds, respectively, denoted by \(\Wcs{x}\) and \(\Wcu{x}\).

\subsection{Liftings and the symbolic space}
We wish to associate with each orbit of $f$ a sequence of symbols that represents the scatterers it encounters along the way.
However, this approach is problematic on the torus because an orbit could hit the same scatterer several times in succession.
To avoid this issue, let \( \widetilde{\mathcal{D}}\subset \R^2 \) be the \( \Z\times \Z \)-periodic billiard table such that \( \pi(\widetilde{\mathcal{D}})=\mathcal{D} \), where \( \pi\colon [0,1)^2\to \T^2 \) denotes the canonical projection. 
Abusing notation, we denote the billiard map and flow on \( \widetilde{\mathcal{D}} \) by \( f \) and \( \Psi \) as well.
\par Let \( \{B_1, \cdots, B_k\} \) be set of connected components of the lift \( \pi^{-1}(\partial\mathcal{D}) \subset [0,1)^2. \)
We index all the scatterers in the lifted table on \( \R^2 \) using the alphabet \( \mathcal{A} \coloneqq \mathbb{Z}^2\times \{1, \cdots, k\} \), as follows: 
\[
    B_{(i,j;l)} \coloneqq B_l + (i,j),
\]
that is, \( B_{(i,j;l)} \) is the unique scatterer obtained by translating \( B_l \) by \( (i,j) \in \Z^2 \).
Let \( \Sigma \) denote the space \( \mathcal{A}^\mathbb{Z} \) of two-sided sequences of elements in \( \mathcal{A} \).
This is a complete metric space with the usual distance
\[
    d_0(\alpha, \beta) = \sum_{i \in \Z} 2^{-|i|}(1 - \delta_{a_ib_i})
\]
between two sequences \( \alpha = \{a_i\}_{i \in \Z} \) and \( \beta = \{b_i\}_{i \in \Z} \).
In what follows, we will work with the equivalent metric \( \rho \) given by
\[
    \rho(\alpha, \beta) \coloneqq 2^{-\max\{n \in \N; a_i = b_i \text{ for all } |i| \leq n\}}.
\]
We associate to the points in \( \mathcal{M} \) the sequence of indices of scatterers where its lifted trajectory bounces. 
\begin{defi}
The map \( \imath \colon \mathcal{M} \to \Sigma \) is given by
\[
    \imath(s,\varphi) = \{a_i\}_{i \in \Z}, 
\]
where $B_{a_0}\subset \pi^{-1}(\partial \mathcal{D})$, and for all $i \in \mathbb{Z}$, \( a_i \) is the unique element of \( \mathcal{A} \) such that \( \mathrm{Pr}_1(f^i(\pi^{-1}(s), \varphi)) \in B_{a_i} \).
We call the set \( \Sigma_\mathcal{M} \coloneqq \imath(\mathcal{M}) \) the \emph{space of admissible sequences}.
\end{defi}
The space $\Sigma$ admits a natural $\mathbb{Z}^2$-action: for any $(i,j)\in \mathbb{Z}^2$, the action is defined by $\{(i_p,j_p;l_p)\}_{p \in \mathbb{Z}}\mapsto \{(i_p+i,j_p+j;l_p)\}_{p \in \mathbb{Z}}$. We denote by $\hat \Sigma$ and $\hat {\Sigma}_\mathcal{M}$ the respective quotients of $\Sigma$ and $\Sigma_\mathcal{M}$ under this action; the associated map is denoted by $\hat \imath\colon \mathcal{M} \to \hat \Sigma$. 
    \begin{figure}[!h]
        \centering
        \includegraphics[width=.8\textwidth]{Sinai2.png}
        \caption{Lifted table.}
        \label{fig:lifted_table}
    \end{figure}
\begin{remark}
    By definition, the space \( \hat{\Sigma}_\mathcal{M} \) is invariant under the left shift \( S\colon \{a_i\}_{i \in \Z} \mapsto \{a_{i+1}\}_{i \in \Z} \), and \( \hat\imath \circ f = S \circ \hat\imath \).
\end{remark}
\begin{prop}\label{pps:representation_map_continuity} 
    The mapping \( \imath\colon \mathcal{M} \to \Sigma_\mathcal{M} \) is a bijection satisfying the following properties:
    \begin{itemize}
        \item[(i)] the restriction of \( \imath \) to the regular set 
        $\mathcal{M}^\mathrm{reg}$ is continuous;
        \item[(ii)] the inverse mapping \( (\imath\rvert_{\mathcal{M}^\mathrm{reg}})^{-1} \) is Hölder continuous. In particular, \( \imath\rvert_{\mathcal{M}^\mathrm{reg}}\colon \mathcal{M}^\mathrm{reg} \to \Sigma_{\mathcal{M}^\mathrm{reg}} \) is a homeomorphism;
        \item[(iii)] the discontinuity set of \( \imath \) is precisely the singular set \( \cS \).
    \end{itemize}
\end{prop}

\begin{proof}
    For \( \alpha \in \Sigma \) and \( n \in \N \), consider the preimage \( C(\alpha; n) \coloneqq \imath^{-1}(B(\alpha, 2^{-n})) \) of an open ball of \( \rho \).
    It consists of points \( x \in \mathcal{M} \) whose first \( n+1 \) iterations under \( f \),  both past and future, coincide with the truncation of \( \alpha \), or, in other words, 
    \[
        C(\alpha; n) = \{x \in \mathcal{M} \,\rvert\, f^{i}(x) \in B_{a_i} \text{ for all } |i| \leq n+1\}.
    \]
    Let \( x, y \) be two distinct points in \( C(\alpha; n) \) lying in the same connected component of \( \mathcal{M} \), and \( \gamma \) a geodesic in \( \mathcal{M} \) joining them. 
    Considering an involution of the flow, if necessary, we may assume that \( \gamma \) is the trace left on the scatterer containing \( x \) and \( y \) by a dispersing wave front \( W \) of the billiard flow (the suspension of \( f \)). 
    Thus, the trajectories of the endpoints of \( W \) on \( \widetilde{\mathcal{D}} \) are exactly the trajectories defined by \( x \) and \( y \), and since \( x, y \in C(\alpha; n) \), these trajectories hit the same \( n \) first obstacles.
    In particular, any obstacle \( B \) on the path of \( W \) that is not in the sequence \( B_{a_1}, \cdots, B_{a_n} \) would be completely ``engulfed'' by \( W \), and can safely be ignored in the analysis which follows (see Figure \ref{fig:eclipsed_obstacle} below).
    \begin{figure}[!ht]
        \centering
        \includegraphics[width=.8\textwidth]{Eclipse.pdf}
        \caption{A smaller obstacle eclipsing a section of the interior of a wave front.}
        \label{fig:eclipsed_obstacle}
    \end{figure}
    \par As the image \( W_t \) of \( W \) under the time \( t \) map of the billiard flow grows exponentially (cf. \cite{chernov_chaotic_2006}), up to constants \( K, \cB > 0, \) depending on the table (i.e. on the curvature \( \mathcal{K} \) and the maximum free-flight time \( \tau_\mathrm{max} \)), we have
    \[
        \|f^n(x) - f^n(y)\| \geq Ke^{\cB n}\|x-y\|,
    \]
    where \( \|\cdot\| \) denotes the Euclidean metric on \( \mathcal{M} \).
    In particular, the coefficient in the exponent is
    \begin{equation}\label{eq:anosov_constats_billiard_table}
        \cB = \cB_{\mathrm{min}} := \left(\tau_{\mathrm{max}} + \frac{1}{2 \cK_{\mathrm{min}}}\right)^{-1},
    \end{equation}
    i.e., the minimum curvature of dispersing wave fronts of the table \( \mathcal{D} \).
    \par On the other hand, \( \|f^n(x) - f^n(y)\| \leq D \coloneqq \max_{1 \leq i \leq k} \mathrm{diam}(B_i) \), and therefore
    \begin{equation}\label{eq:lips_estimate}
        \|x -y\| \leq \frac{D}{K}e^{-\cB n} = c(\mathcal{D})e^{-\cB n},
    \end{equation}
    for some constant \( c(\mathcal{D}) \) depending on the table. 
    In particular, if \( \imath(x) = \imath(y) = \omega \), then \( x, y \in \cap_{n}\imath^{-1}(B(\alpha, 2^{-n})) \), and therefore 
     \( \|x -y\| \leq c(\mathcal{D})e^{-\cB n} \) for all \( n \in \N \), hence \( x = y \), and \( \imath \) is indeed a bijection.     
    \par Now, let us denote by \( \jmath \colon \mathcal{M}^\mathrm{reg} \to \Sigma \) the restriction of \( \imath \) to the regular set.
    We claim that \( \jmath^{-1}(B(\alpha, 2^{-n})) = C(\alpha; n)\cap\mathcal{M}^\mathrm{reg} \) is open.
    Indeed, since each \( f^i \) is continuous on \( \mathcal{M}^\mathrm{reg} \), there is a neighbourhood of any \( x \) in such preimage where every point bounces in exactly the same obstacles as \( x \) for \( |i| < n+1 \).
    Hence every point in \( C(\alpha; n)\cap\mathcal{M}^\mathrm{reg} \) is interior and the map \( \jmath \) is continuous, proving (i).
    \par As for (ii), note that if \( \rho(\alpha, \beta) = 2^{-n} \), then it follows from the inequality in \eqref{eq:lips_estimate} that, 
    \[
        \|\imath^{-1}(\alpha)-\imath^{-1}(\beta)\|  \leq c(\mathcal{D})e^{-\cB n} \leq c(\mathcal{D})\rho(\alpha, \beta)^{\cB \log 2^{-1}},
    \]
    hence \( \imath^{-1} \) is \( \cB \log 2^{-1}\)-Hölder continuous. 
    \par Finally, we prove (iii).
    Property (i) can be rephrased as saying that the discontinuity set of \( \imath \) is contained in the singular set \( \cS \), so all that remains to show is the converse inclusion.
    Let \( x_0 \in \cS \) and assume, without loss of generality, that \( f(x_0) \in \cS_0 \). 
    Let \( \pi_\mathcal{D}(x) \in \Gamma_i \) and \( \pi_\mathcal{D}(f(x_0)) \in \Gamma_j \).
    Then, due to the discontinuity of \( f \) at \( x_0 \), we can find points \( x \) arbitrarily close to \( x_0 \) and such that \( \pi_\mathcal{D}(f(x)) \notin \Gamma_j \), as well as points arbitrarily close and such that \( \pi_\mathcal{D}(f(x)) \in \Gamma_j \). 
    Thus, there are sequences \( x_n \to x_0 \) in \( \mathcal{M} \) such that \( \rho(\imath(x_n), \imath(x_{n+1})) = 1 \) for every \( n\in \N \), and \( x_0 \) is therefore a discontinuity point for \( \imath \).\qedhere
\end{proof}



\subsection{Periodic orbits.}\label{ssec:periodic_orbits} 
\par Although it is not needed for the main part of the argument, we collect in this section general properties of periodic orbits and their lengths. 
As a corollary of Proposition~\ref{pps:representation_map_continuity}, we have:
\begin{corollary}\label{cor:fin_card_peri}
	For each \( T>0 \), the number of periodic orbits of period at most \( T \) is finite.
\end{corollary}
\begin{proof}
    Fix \( T>0 \). Fix a periodic point $(s,\varphi)$ whose orbit under the billiard flow has length at most $T$. Let $q>0$ be its period under the billiard map, $f^q(s,\varphi)=(s,\varphi)$. In particular, we have $q \leq q_0$, with \( q_0:=[\frac{T}{\tau_\mathrm{min}}]+1 \). Let \(\imath(s,\varphi)=:(\rho_k)_{k}\in \mathcal{A}^\mathbb{Z} \). 
    Since $f^q(s,\varphi)=(s,\varphi)$, we have $\rho_{q}=\rho_0+(\tilde i,\tilde j;0)$, for some $(\tilde i,\tilde j)\in \mathbb{Z}^2$, and 
    $$
    \rho_p=\rho_{p \text{ mod }q}+\left\lfloor \frac{p}{q}\right\rfloor(i_0,j_0;0),\quad \forall\, p \in \mathbb{Z}.
    $$
    Moreover, by the finite horizon assumption, and $\mathbb{Z}^2$-periodicity of the unfolding, for each step $p\in \mathbb{Z}$, there exists a uniform upper bound $K\geq 1$ on the difference $|i_{p+1}-i_p|+|j_{p+1}-j_p|\leq K$ between the labels $(i_p,j_p)$ and $(i_{p+1},j_{p+1})$ of the cells at steps $p$ and $p+1$. In particular, we have $|\tilde i|+|\tilde j|\leq K q$. We conclude that periodic points of period at most $T$ are completely encoded by finite words $(\rho_1,\cdots,\rho_q)$, $q\leq q_0$, with at most $Kk$ choices for each symbol $\rho_p$. Therefore, the number of (infinite) words $\imath(s,\varphi)$ associated to periodic points with period at most $T$ is finite. By the injectivity of $\imath$ shown in Proposition~\ref{pps:representation_map_continuity}, we deduce that the set of periodic points for the billiard map $f$ (and consequently, the set of periodic orbits) of period less than $T$ is finite. \qedhere
   
\end{proof}

Given an integer $m\geq 2$, we denote by \( \mathfrak{D}^m \) the space of all Sinai billiard tables in \( \T^2 \) with \( m \) scatterers and $C^r$ boundary, $r \geq 3$, that have finite horizon. 
Each table \( \mathcal{D} \) is defined by its boundary map \( \T_\mathcal{D} \to \T^2 \), allowing us to identify \( \mathfrak{D}^m \) with the space of embeddings \( C^r_{\text{emb}}(\T_\mathcal{D}, \T^2) \), endowed with the $C^r$ topology; in particular it is a Baire space. 
\par We want to show that for a generic choice of table in \( \mathfrak{D}^m \), the set of periodic orbits has simple length spectrum and contains no singular orbits. 

\begin{defi}
For each integer $n\geq 1$, we define \( \mathcal{R}_n \) as the subset of \( \mathfrak{D}^m \) consisting of tables without any singular periodic orbits of period less than or equal to \( n \).
\end{defi}

\begin{defi}
For each integer $n\geq 1$, we define \( \mathcal{S}_n \) as the subset of \( \mathfrak{D}^m \) consisting of tables such that no distinct periodic orbits of period less or equal to \( n \) have the same length.
\end{defi}

\begin{prop}\label{prop_good_tables}
For each integer $n \geq 1$, the set \( \mathcal{R}_n \cap \mathcal{S}_n \)  is open and dense in $\mathfrak{D}^m$. 
    Let us then define the set 
    \begin{equation*}
        \mathfrak{G} := \bigcap_n \mathcal{R}_n \cap \mathcal{S}_n,
    \end{equation*}
   which represents the set of ``good'' Sinai billiard tables of class $C^r$ with finite horizon, specifically, a table $\mathcal{D}$ is in $\mathfrak{G}$ if it has no grazing periodic orbits and has a simple length spectrum. Then, the set $\mathfrak{G}$ is residual in the $C^r$ topology. 
\end{prop}
We defer the proof of Proposition~\ref{prop_good_tables} to Appendix~\ref{sec_periodicorbitsbilliard}, where  further properties of periodic orbits are investigated. 


\section{Geodesic approximations}\label{sec:geodesic_approximations}
\subsection{Kourganoff surfaces}
Recall that \( \Psi \) is the billiard flow of the table \( \mathcal{D} \), and let \( \Omega \) denote its phase space.
It is the suspension space of the billiard map \( f \) with roof function \( \tau \).
We consider \( E= \mathbb{T}^2\times \R \) with coordinates \( (x,y,z) \) and a canonical global frame \( (\partial_x, \partial_y, \partial_z) \).
We write \( p\colon E \to \mathbb{T}^2 \) to denote the canonical projection on the first two coordinates.
\par Consider a smooth closed immersed surface \( K \) in \( E \) such that \( p(K) = \mathcal{D} \). 
In what follows, we abuse a bit the notation and write
\begin{align*}
    K^\circ &:= K \cap p^{-1}(\mathrm{int}(\mathcal{D})), \\
    B(K)&:= K \cap p^{-1}(\partial\mathcal{D}).
\end{align*}
Moreover, we assume \( K \) satisfies
\begin{enumerate}
    \item[(K1)]\label{cond_k_un} \emph{
        The vector \( \partial_z \) is not tangent to \( K \) at any point \( x \in K^\circ \), i.e. \( \partial_z \notin T_xK \)
        };
    \item[(K2)] \emph{
        For every point \( x \in B(K) \) the curve defined by the intersection of \( K \) and the affine plane \( \{x\} + (\mathbb{R}\partial_z \oplus T^\perp_xK) \) has non-vanishing curvature around \( x \)
        }.
\end{enumerate}
\par Note that each connected component of \( K^\circ \) is diffeomorphically mapped to \( \mathrm{int}(\mathcal{D}) \) under \( p \), since \( p(K^\circ) = \mathrm{int}(\mathcal{D}) \) and \( \mathrm{int}(\mathcal{D}) \) are connected.
Similarly, each connected component of \( B(K) \) is projected diffeomorphically onto the boundary of a single scatterer \( \Gamma_i \subset \partial\mathcal{D} \).
In particular, if \( p(x) = p(x') \in \partial \mathcal{D} \), then we can join them by a path \( \eta\colon (-\delta, \delta) \to S \) that lies entirely in a single connected component \( S \approx \mathcal{D} \) of \( K^\circ \).
Since \( \overline{S} \approx \mathcal{D} \) and the projection of \( \eta \) has a single endpoint, we conclude that \( x = x' \) (otherwise \( \overline{S} \) would have more boundary components than \( \mathcal{D} \)).
Thus, \( \#p^{-1}(x) = 1 \) for all \( x \in \partial \mathcal{D} \). 
We should therefore think of \( K \) as a surface obtained from \( \mathcal{D} \) by gluing different copies of \( \mathcal{D} \) at their respective boundary components. 
The number of copies is exactly the number of connected components of \( K^\circ \).
In what follows, it might be helpful to think of \( K \) as two copies of \( \mathcal{D} \) identified by the corresponding obstacles. We refer the reader to Figure~\ref{liftsurface_png} for a picture of the construction. 

\par For \( \epsilon \in (0,1] \), we let \( \alpha_\epsilon \) be the contraction given by \( (x,y,z) \mapsto (x, y, \epsilon z) \) and write \( g_\epsilon \) to denote the metric induced on \( K_\epsilon := \alpha_\epsilon(K) \) by the Euclidean metric.
Pulling back via \( \alpha_\epsilon \), we will consider \( \{g_\epsilon\} \) as a family of metrics on the surface \( K \). 
We denote by \( \mathbf{g}_\epsilon \) the geodesic flow of the surface \( (K, g_\epsilon) \).
Then \( p \) induces a smooth mapping
\begin{equation}\label{eq:projection_flow_equivalence}
\begin{split}
    \pi_K\colon T^1K^\circ &\to\Omega \\
    (x,v) &\mapsto \left(p(x), \frac{dp_x(v)}{\|dp_x(v)\|}\right)
\end{split}
\end{equation}
(note that~(K1) ensures we are not dividing by \( 0 \) with such a definition).
Finally, we denote by \( A_0 \) the subset of \( \R \times T^1K^\circ \) consisting of points \( (t; x,v) \) such that the orbit segment \( \{\Psi^s(\pi_K(x,v)); s \in [0,t]\} \) has endpoints in \( \mathrm{int}(\mathcal{D}) \) and does not contain grazing collisions. 
\(A_0\) is easily seen to be open, and since it contains the subset \(\R \times \pi_K^{-1}(\Omega^\mathrm{reg})\), it is also dense in \(\R \times T^1K^\circ\).
Define mappings \( F_\epsilon\colon A_0 \to \Omega \), \( \epsilon \in [0,1] \) given by
\begin{align*}
    F_0\colon (t; x,v) &\mapsto \Psi^t(\pi_K(x,v)), \text{ and} \\
    F_\epsilon\colon (t; x,v) &\mapsto \pi_K(\mathbf{g}^t_\epsilon(x,v)), \text{ for } \epsilon > 0.
\end{align*}
\begin{theorem}[Theorem 5 in \cite{kourganoff_anosov_2016}]\label{thm:kourganoff_approximation}
    The family of maps \( F_\epsilon \) converges uniformly \( F_0 \), as \( \epsilon \to 0 \), on every compact subset of \( A_0 \).
    Moreover, there is \( \epsilon_0 > 0 \) such that \( \mathbf{g}_\epsilon \) is an Anosov flow for every \( \epsilon < \epsilon_0 \).
\end{theorem}
In particular, up to rescaling, there is no loss of generality in assuming
\begin{itemize}
    \item[(K3)] \emph{
        the family \( \{\mathbf{g}_\epsilon\}_{\epsilon \in (0,1]} \) is a one-parameter family of Anosov flows in \( K \).
        }
\end{itemize}
We will denote by \( \cK(\mathcal{D}) \) the collection of all such surfaces:
\[
    \cK(\mathcal{D}) := \{K \subset E; p(K) = \mathcal{D} \text{ and } K \text{ satisfies (K1), (K2) and (K3)} \}.
\]

\begin{figure}[!h]
    \centering
    \includegraphics[width=.8\textwidth]{Kourganoff.pdf}
    \caption{\textbf{Left:} A geodesic arc on some surface $K_\epsilon$ and its projection on the billiard table $\mathcal{D}$. \textbf{Right:} for $\epsilon>0$ small and $(t;x,v)\in A_0$, the projection $\{F_\epsilon(s;x,v)\}_{s\in [0,t]}$ closely shadows the billiard orbit segment $\{F_0(s;x,v)\}_{s\in [0,t]}$ (Credits for pictures: Mickaël Kourganoff, Jos Leys).} 
\end{figure}

\subsubsection{ Metric convergence in \( \cK(\mathcal{D}) \) } 
Fix a Kourganoff surface \( K \in \cK(\mathcal{D}) \).
Given a patch \(  \phi(u,v) = (x(u,v), y(u,v), z(u,v))  \) for \(  K  \), we consider the associated patch \(  \alpha_\epsilon \circ \phi = (x(u,v), y(u,v), \epsilon z(u,v))  \) for \(  K_\epsilon  \). 
Comparing the First Fundamental Forms \( I = E \, du^2 + 2F \, du \, dv + G \, dv^2 \) and \( I_\epsilon = E_\epsilon \, du^2 + 2F_\epsilon \, du \, dv + G_\epsilon \, dv^2 \) for \( K \) and \( K_\epsilon \) in these coordinates, we obtain 
\begin{equation}\label{eq:first_fund_forms_comparison}
\begin{split}
    E_\epsilon &= E + (\epsilon^2 - 1) (\partial_u z)^2, \\
    F_\epsilon &= F + (\epsilon^2 - 1) \partial_u z\partial_v z, \\
    G_\epsilon &= G + (\epsilon^2 - 1) (\partial_v z)^2. 
\end{split}
\end{equation}
Hence, for any \( \epsilon, \delta \in (0,1] \), the difference between the coefficients of the metrics \(g_\epsilon\) and \(g_\delta\), as well as the differences between their derivatives, are of order \( |\epsilon^2 - \delta^2| \).
It follows that the \(C^r\) norm of their difference, as a tensor field on \( K \), can be bounded by a term 
\begin{equation}\label{eq:bound_geom}
    \|g_\epsilon - g_{\delta}\|_{C^r(K)}     
    \leq  C|\epsilon^2 - \delta^2| \leq C,  
\end{equation}
where \( C \) is a constant depending only on \( K \) (it comes from integrating of the coordinate function \(z\) and its derivatives over the entire surface).
Hence, given a curve \( \gamma \), denoting by $l_\epsilon(\gamma)$ its length with respect to the metric $g_\epsilon$, the magnitude of length variation of $\gamma$ for two different metrics \( g_\epsilon \) and \( g_{\delta} \) is  
\begin{equation}\label{eq:length_estimate}
    \left| l_\epsilon(\gamma) - l_{\delta}(\gamma) \right| \leq \int_{\gamma} \left| \sqrt{I_\epsilon} - \sqrt{I_{\delta}} \right| 
    = O(|\epsilon -\delta|).
\end{equation}

A first consequence of the equation above is that a curve \( \gamma \)\footnote{
    Please remark that through this section we frequently abuse the terminology and refer to both the maps \( [a,b] \to K \) and their images as curves, letting the distinction be dictated by context. 
} is rectifiable (i.e., has finite length) with respect to \( g_\epsilon \) if and only if it is rectifiable with respect to the reference metric \( g \), thus the class of rectifiable curves is the same for all the metrics.
Given \( x,y \in K \), we will denote by \(  [x,y]_\epsilon  \) the minimising geodesic for \(  g_\epsilon  \) joining them.
Note that \( \{[x,y]_\epsilon\} \) is a collection of curves in \( K \) with endpoints \( x,y \), all of which are rectifiable with respect to every metric in the family \( \{g_\epsilon\} \).
Moreover, the distance functions \(  d_\epsilon  \) and \(  d_{\delta}  \) satisfy
\[
    d_\epsilon(x,y) := l_\epsilon([x,y]_\epsilon) \leq l_\epsilon([x,y]_{\delta}) = l_{\delta}([x,y]_{\delta}) + O(|\epsilon - {\delta}|) = d_{\delta}(x,y) + O(|\epsilon - {\delta}|),
\]
and therefore
\begin{equation}\label{eq:unif_conver_dist}
    |d_\epsilon(x,y) - d_{\delta}(x,y)| = O(|\epsilon - {\delta}|),
\end{equation}
for every \( x,y \in K \).
\par More generally, Equation \eqref{eq:bound_geom} implies that the Christoffel symbols of \( g_{\delta} \) converge uniformly to those of \( g_\epsilon \), and the corresponding geodesic vector fields on \( T(TK) \) can be made arbitrarily close in the reference metric \( g_1 \).
Rewriting the statement of the continuous dependence of ODE's on the parameters (c.f. \cite[Chapter 3.12]{walter2013ordinary}, for instance) in terms of the geodesic equations, we obtain the following
\begin{lemma}\label{lem:uniform_conv_geodesic_flows}
    Let \( \{X_n\} \) in \( T^1K \) be a Cauchy sequence with respect to the Sasaki metric \( G \) induced on \( T(TK) \) by \( g \).  
    Then, for every compact set \( J \subset \R \) and every \( \eta > 0 \), there are constants \( N = N(J, \eta) \in \N \) and \( \epsilon_0 = \epsilon_0(J, \eta) > 0 \) such that, for any \( \epsilon, \delta \) with \( |\epsilon-\delta| < \epsilon_0 \) and any \( n > N \), it holds
    \[
        \forall t \in J, \ \ d_G(\mathbf{g}_\epsilon^t(X_n), \mathbf{g}_{\delta}^t(X_n)) < \eta,
    \]
    where \( d_G \) is the metric induced by \( G \) in \( TK \).
\end{lemma}

\begin{lemma}\label{lem:uniform_lower_bound_inj_radius}
    There is a uniform lower bound \( R_0 > 0 \) on the injective radius \( R_\epsilon \) of the metrics \( g_\epsilon \). 
\end{lemma}
\begin{proof}
    As the geodesic flows \( \mathbf{g}_\epsilon \) are all Anosov, none of the surfaces \( K_\epsilon \) have conjugated points (c.f. \cite{klingenberg_riemannian_1974}).
    Therefore, according to Klingenberg's lemma (cf. \cite[Lemma 1.8]{AbreschMeyer1997}) the injective radius of \( K_\epsilon \) is
    \[
        \mathrm{inj}(g_\epsilon) = \frac{1}{2}l_\epsilon(\sigma_\epsilon)
    \]
    where \( \sigma_\epsilon \) is the closed geodesic of minimal \( g_\epsilon \)-length in \( K \). 
    We claim that 
    \[
        \inf \{l_\epsilon(\sigma_\epsilon)\} \geq \min\{l_i, 2\tau_\mathrm{min}\} > 0,
    \] 
    where \( l_i \) is the length of the scatterer \( \Gamma_i \subset \mathcal{D} \).
    \par Since the length functional \( l_\epsilon \) is additive with respect to the juxtaposition of paths, the free homotopy class \( [\sigma_\epsilon] \in \mathcal{C}(K) := \mathrm{Conj}(\pi_1(K)) \) must be one of the generator loops, due to minimality.
    Now, the generators of \( \pi_1(K) \) are divided in the following three categories.
    \begin{itemize}
        \item[(i)] loops freely homotopic to great circles in one of the connected components \( S \approx \mathrm{int}(\mathcal{D}) \) of \( K^\circ \);
        \item[(ii)] loops freely homotopic to a connected component \(\Gamma_i\) of \( B(K) \);
        \item[(iii)] loops with a base point in one connected component of \( B(K)\) and meeting two or more distinct connected components of \( K^\circ\). 
    \end{itemize}
    Thus, the image \( \gamma_\epsilon := p(\sigma_\epsilon) \) must be either:
    \begin{itemize}
        \item[(i)] a closed curve transverse to a great circle;
        \item[(ii)] a curve homotopic to one of the obstacles \( \Gamma_i \subset \mathcal{D} \);
        \item[(iii)] a closed curve tangent to at least two distinct boundary components \( \Gamma_i \) and \( \Gamma_j \).
    \end{itemize}  
    In the first case, the length \( l_E(\gamma_\epsilon) \) (with respect to the flat metric on torus) is bounded below by the minimum length of the great circles.
    Remark that these lengths must be strictly larger then the minimum flight time \( \tau_\mathrm{min} \) of the table due to the finite horizon. 
    In case (ii), the length is bounded below by the minimum length\( \{l_i\} \) of the scatterers \( \Gamma_i \). 
    In the last case, \( l_E(\gamma_\epsilon) \geq 2\tau_\mathrm{min} \)
    (see Figure \ref{fig:projected_limiting_curve} below).
    \begin{figure}[!ht]
        \centering
        \includegraphics[width=.7\textwidth]{lift3.pdf}
        \caption{Some loops $\{\sigma_i\}_{i=1}^3$ on the surface $K$, and their projections $\{\gamma_i=p(\sigma_i)\}_{i=1}^3$ on the billiard table $\mathcal{D}$.}
        \label{fig:projected_limiting_curve}
    \end{figure}    
    \par Finally, by construction \( \partial_z \) is not tangent to \( K \) at any point, hence \( \dot\sigma_\epsilon \notin \R \partial_z = \ker p. \) 
    It follows that, for any \( \epsilon > 0 \),
    \[
        l_\epsilon(\sigma_\epsilon) \geq l_E(\gamma_\epsilon) \geq \min\{l_i, 2\tau_\mathrm{min}\} > 0,
    \]
    which concludes.    
\end{proof}
\begin{lemma}\label{lem:convergence_geodesics}
    For any \( x,y \in K^\circ \) sufficiently close (\( d_1(x,y) < R_0 \)), the geodesics \( [x,y]_\epsilon \) converge uniformly to a \( C^1 \) curve \( [x,y]_0 \) in \( K \).
\end{lemma}
\begin{proof}
    Let \( \widetilde x \in K^\circ, x = p(\widetilde x) \in \mathrm{int}(\mathcal{D}) \), and consider a neighbourhood \( U \) of \( \widetilde x \), which we choose sufficiently small so that \( \overline U \) is a normal neighbourhood for all of the metrics \( g_\epsilon \), and is moreover mapped under \( p \) into an open ball \( B := B_{\mathcal{D}}(x, \delta) \) of radius \( \delta < \inf\{\tau(x, p); p \in \mathbb{S}^1\}. \)
    In other words, the billiard trajectories starting at \( x \), in any possible direction, do not experience any collision at all inside \( B \).
    \par Moreover, for a fixed \( \widetilde y \in U \), distinct from \( \widetilde x \), we write \( l_\epsilon = d_\epsilon(\widetilde x, \widetilde y) \) and let
    \( \widetilde{v}_\epsilon \) be the velocity of the curve \( [\widetilde x,\widetilde y]_\epsilon \) at the starting point \( \widetilde x \).
    Finally, we set \( y = p(\widetilde y) \), let \( T = |x-y|, v = T^{-1}(x-y) \), and consider the geodesics \( \gamma_\epsilon \) in \( K_\epsilon \) with initial condition \( (\widetilde x,\widetilde v) \), where \( \pi_K(\widetilde x, \widetilde v) = (x, v) \).
    \par By construction, the point \( (T; \widetilde x,\widetilde v) \) lies in \( A_0 \), and Theorem \ref{thm:kourganoff_approximation} guarantees that 
    \( \pi_K(\mathbf{g}_\epsilon^t(\widetilde x,\widetilde v)) \) converges uniformly to \( \Psi_\epsilon^t(x, v) \) in the interval \( [0, T] \) as \( \epsilon \to 0 \). 
    Thus, for \( \epsilon \) sufficiently small we have \( \gamma_\epsilon(T) \) arbitrarily close to \( \widetilde y \). 
    Since the exponential map at \( \widetilde x \) is a diffeomorphism in \( U \), we can estimate the distance between \( \widetilde v_\epsilon \) and \( \widetilde v \) by
    \[
        \|\widetilde v_\epsilon - \widetilde v\|_\epsilon \leq \left\|d(\exp_{\widetilde x}^{-1})\rvert_{\overline U}\right\|_\epsilon d_\epsilon(\widetilde y, \gamma'_\epsilon(T)).
    \]
    Around \( \widetilde x \) the norm of the derivative is close to \( 1 \).
    Moreover, as pointed out in \cite{kourganoff_anosov_2016}, as the surface \( K \) flattens out, the curvature away from \( B(K) \) tends to \( 0 \), while in a neighbourhood of \( B(K) \) it becomes strictly negative.
    Therefore, away from \( \widetilde x \), for \( \epsilon \) sufficiently small, all the sectional curvatures of \( K_\epsilon \) on \( \overline U \) can be made arbitrarily close to \( 0 \), if need be. 
    Hence there is a constant \( 1/R^2 > 0 \) which bounds all the sectional curvatures, and we can estimate the derivative of \( \exp_x \)  in terms of the norms of Jacobi fields (using \cite[Theorem 11.2]{lee2019introduction}, for instance,), which are of order \( O(R\sin(R^{-1})) \). 
    We may then find a constant \( L \), independent of \( \epsilon \), and such that
    \[
        \|\widetilde v_\epsilon - \widetilde v\|_\epsilon \leq Ld_\epsilon(\widetilde y, \gamma_\epsilon(T)).
    \]
    Thus the family \( \{\widetilde v_\epsilon \} \) has a unique accumulation point at \( \widetilde{v} \).
    In particular, it now follows from Lemma \ref{lem:uniform_conv_geodesic_flows} that \( \{ [x,y]_\epsilon\} \) and their derivatives form an equicontinuous family of curves in \( T^1K \), and therefore converge to a well-defined curve \( [x,y]_0 \).
    What is more, \( [x,y]_0 \) is differentiable at \(x\), and its tangent space is spanned by the limit vector \(v\).
\end{proof}
Essentially the same argument shows that we may fix the tangent vectors \(v_\epsilon\) and let the endpoint \(y\) vary instead, yielding the following.
\begin{corollary}
    Given \((x,v) \in T^1K^\circ\), let \(\gamma_{\epsilon;v}\) denote the unique geodesic of \((K, g_\epsilon)\) with initial condition \((x,v)\).
    If \(U\) is any connected neighbourhood of \(x\) not intersecting \(B(K)\), then there is a unique differentiable curve
    \(\gamma_{0;v}\), starting at \(x\), and whose tangent space at \(x\) is spanned by \(v\).
\end{corollary}
\begin{theorem}\label{thm:d_0_is_a_metric}
    For any distinct points \( x, y \in K \), 
    \[
        d_0(x,y) := \lim_{\epsilon\to 0}l_\epsilon([x,y]_\epsilon) > 0.
    \] 
    Moreover, \( K_0 := (K, d_0)  \) is a complete geodesic length space. 
    The admissible paths are precisely the rectifiable curves with respect to the reference metric \( g \) and the length function is 
\begin{align*}
    l_0\colon \mathrm{Rec}(K, g) &\to \R \cup \{+\infty\} \\
    \gamma &\mapsto \lim_{\epsilon \to 0}l_\epsilon(\gamma).
\end{align*}
The geodesics\footnote{
    We mean geodesic in the sense of length spaces, i.e., a locally minimising path with respect to \( d_0 \).
    In general, the tensors \( g_\epsilon \) degenerate as \( \epsilon \to 0 \) and the metric \( d_0 \) is not associated to any Riemannian structure.
    }
of \( K_0 \) are exactly the uniform limits of geodesics of \( K_\epsilon \) in the \( C^1 \)-topology.
In particular, they are differentiable curves on \(K\).
\end{theorem}
\begin{proof}
    Remark that \( d_0(x,y) \) is well defined for any points \(x, y \) due to Equation \eqref{eq:unif_conver_dist}. 
    The symmetry of \( d_0 \) and the triangle inequality follow straightforwardly from the definition, as well as the fact that \( d_0(x, x) = 0 \). 
    \par We begin by considering \( x \) in the interior of \( \mathcal{D} \).
    Fix \( R_0 \leq \inf_\epsilon\{\mathrm{inj}(g_\epsilon)\} \).
    Given \( y \neq x \), if \( l_\epsilon([x,y]_\epsilon) > R_0 \) for every \( \epsilon \) then \( d_0(x,y) > R_0 \).
    If not, then Lemma \ref{lem:convergence_geodesics} applies, and we get that \( [x,y]_\epsilon \) converges uniformly in the \( C^1 \) topology to a curve \( [x,y]_0 \) whose derivative at \( x \) is a vector \( v \) satisfying 
    \[
        dp_x(v) = \frac{p(x) - p(y)}{|p(x) - p(y)|}.
    \]
    Moreover, by considering \( \epsilon \) sufficiently small, the geodesics \( \gamma_\epsilon \) and \( [x, y]_\epsilon \) can be made arbitrarily close in the uniform \( C^1 \) topology on the interval \( [0, |p(x) - p(y)|]\). 
    Thus
    \[
        d_0(x,y) := \lim_{\epsilon\to 0}l_\epsilon([x,y]_\epsilon) = \lim_{\epsilon \to 0}l_\epsilon(\gamma_\epsilon) = |p(x) - p(y)| > 0.
    \]
    Finally, since every geodesic \(g_\epsilon\) is complete, we conclude that every minimising curve \( [x,y]_0 \) is uniquely extensible to small neighbourhoods of \(x\) and \(y\), as long as such neighbourhoods do not intersect \(B(K) \).
    \par In the general case, the intersection \( [x,y]_\epsilon \cap B(K) \) may be non-empty for infinitely many values of \( \epsilon \), but the curves \( ([x,y]_\epsilon) \) can be arbitrarily well approximated by polygonal curves lying entirely in \( K^\circ \), where the previous arguments apply to each sufficiently small segment.
    It follows that \( d_0(x, y) > 0 \) as well, which concludes the demonstration that \( K_0 = (K, d_0) \) is a metric space.
    It is also straightforward that any segment \( [x,y]_0 \) can be extended to the entire line.
    For endpoint points in \(B(K)\), the extension is not unique.
    \par Now, note that \( d_\epsilon \) converges uniformly to \( d_0 \) as a consequence to Equation \eqref{eq:unif_conver_dist}. 
    In particular, up to isometry, \( K_0 \) is the unique Gromov-Hausdorff limit of the complete length spaces \( K_\epsilon \), and therefore it is a complete length space on its own \cite[Theorem 7.5.1]{burago2022course}.
    Additionally, the length function \( l_0 \) of \( d_0 \) may be reconstructed from the middle points of distances (cf. \cite[Theorem 2.4.16.]{burago2022course}), but since those are the limits of middle points of the distances \( d_\epsilon \), it follows that
    \[
        l_0(\gamma) = \lim_{\epsilon \to 0}l_\epsilon(\gamma).
    \]
    Finally, let \( \gamma_\epsilon\) be a geodesic of \( K_\epsilon \) with initial condition \( (x, v) \in T^1K \).
    By Lemma \ref{lem:uniform_conv_geodesic_flows} there is a unique curve \( \gamma \) on \( K \) to which the family \( \gamma_\epsilon \) converges on every compact subset of \( \R. \)
    Given \(x_0 = \gamma(a_0), y_0 = \gamma(b_0)\) sufficiently close, there are sequences of points \( x_\epsilon = \gamma_\epsilon(a_\epsilon) \) and \( y_\epsilon = \gamma_\epsilon(b_\epsilon) \) with \( d_\epsilon(x_\epsilon, y_\epsilon) = l_\epsilon(\gamma_\epsilon\rvert_{[a_\epsilon, b_\epsilon]}) \) and such that
    \( x_\epsilon \to x_0, y_\epsilon \to y_0 \) as \( \epsilon \to 0 \).
    Hence
    \[
        |d_\epsilon(x_\epsilon, y_\epsilon) - d_\epsilon(x_0,y_0)| \leq d_\epsilon(x_0, x_\epsilon) + d_\epsilon(y_0, y_\epsilon) \to 0 \text{ as } \epsilon \to 0,
    \]
    and consequently 
    \[
        \lim_{\epsilon \to 0}l_\epsilon(\gamma_\epsilon\rvert_{[a_\epsilon, b_\epsilon]}) = \lim_{\epsilon \to 0}d_\epsilon(x_\epsilon, y_\epsilon) = d_0(x_0, y_0).
    \]
    But \(l_0(\gamma\rvert_{[a_0, b_0]}) = \lim_{\epsilon \to 0}l_\epsilon(\gamma_\epsilon\rvert_{[a_\epsilon, b_\epsilon]})\) by construction.
    In other words, the curve \( \gamma =\lim_{\epsilon \to 0}\gamma_\epsilon \) is a minimising curve between any two points in its image, which concludes.
\end{proof}

\begin{remark}
    The results above, in particular Theorem~\ref{thm:d_0_is_a_metric}, do not guarantee uniqueness of the geodesic  when the initial condition lies in \(T^1B(K)\).
    In such a situation, if \(v\) is transverse to \(B(K)\) (that is, if \(\pi_K(x,v)\) corresponds to a regular collision on \(\Omega\)), then Kourganoff's approximation theorem guarantees \(\gamma_{0,v}\) is unique.
    Otherwise, if \(v\) is tangent to \(B(K)\), then there are at least two geodesics in \(K_0\) starting at \(x\) with velocity \(v\). 
    One projects, on \(\Omega\), onto the billiard orbit whose trajectory in the table \(\mathcal{D}\) is the line \(t \mapsto x + tv\), and the other is the component of \(B(K)\) containing \(x\).
\end{remark}

\subsubsection{\( K_0 \) as a metric gluing}
\par There is yet another way one can construct \( K_0 \).
First, examining the proof of Lemma \ref{lem:convergence_geodesics}, we note that it was showed that, given a connected component \( S \) of \( K^\circ \), every \( x \in S \) has a neighbourhood \( U \) such that, for every \( y \in U \), the curve \( [x, y]_0 \) satisfies
    \[
        l_0([x, y]_0) = \lim_{\epsilon \to 0}l_\epsilon([x,y]_\epsilon) = |p(x) - p(y)|.
    \]
We can then restate the main ingredient of the proof as
\begin{lemma}\label{lem:d_0_is_distance_reestated}
    Let \( S \) be a connected component of \( K^\circ \) and \( |\cdot| \) the Euclidean distance in \( \mathcal{D} \). Then \( p\rvert_S\colon (S, d_0) \to (\mathcal{D}, |\cdot|) \) is a local isometry. 
\end{lemma}
Naturally, the argument is highly dependent on the approximation by geodesic flows, which fails if the points whose distance we are measuring are eclipsed by a scatterer, or after the billiard trajectory experiences a grazing collision.
To remedy that, we consider a different metric on \( \mathcal{D} \), making it into a length space.
\begin{defi}[Length metric of \( \mathcal{D} \)]
    The table \( \mathcal{D} \) has a length structure \( \mathcal{L} \) consisting of all Euclidean rectifiable paths \( \gamma\colon I \to \R^2 \) whose image is completely contained in \( \mathcal{D} \) together with the usual Euclidean length \( l_E \) induced by \( |\cdot| \).
    We call the \emph{length metric} of \( \mathcal{D} \), denoted by \( d_\mathcal{L} \), the metric induced by \( \mathcal{L} \), i.e.,
    \[
        d_\mathcal{L}(x, y) = \inf\{l_E(\gamma); \gamma \in \mathcal{L}, \gamma(0) = x, \gamma(1) = y\}.  
    \]    
\end{defi}
The geodesics of \( d_\mathcal{L} \) are straight lines in \( \mathcal{D} \), segments of boundary, and paths consisting of concatenations of these two types: such curves are straight lines that meet a scatterer at a tangential angle, follow along the boundary for a while, and leave into a straight trajectory again, with another grazing angle (see Figure \ref{fig:types_geodesics}).
\begin{figure}[!h]
    \centering
    \includegraphics[width=.8\textwidth]{Geodesics.pdf}
    \caption{The longer curve following the obstacles is a \(d_\mathcal{L}\)-geodesic, while the billiard trajectory segment is merely piecewise geodesic.}
    \label{fig:types_geodesics}
\end{figure}
\begin{remark}
    Just as \(K_0\), the space \((\mathcal{D}, d_\mathcal{L})\) is not uniquely geodesic.
    If \(x \in \Gamma_i \subset \partial\mathcal{D}\) is a boundary point, and \(v\) is the tangent vector to \(\Gamma_i\) at \(x\), then there are two geodesics starting at \(x\) in the direction of \(v\): the straight line \(x + tv\) into the interior of the table, and the curve \(\Gamma_i\) itself.
    Similarly, the points of \(B(K)\) are exactly the ones where \(d_0\) fails to be uniquely geodesic.
\end{remark}
\begin{lemma}\label{lem:sheet_isometry}
    For every connected component \( S \subset K^\circ \), \( p\rvert_{\overline{S}}\colon (\overline S, d_0) \to (\mathcal{D}, d_\mathcal{L}) \) is an isometry.
\end{lemma}
\begin{proof}
    The projection \( p \) is Lipschitz with respect to metric \( g\) on \(K\) and the Euclidean norm \(|\cdot| \) on \(\mathbb{T}^2\), hence the image \( p(\gamma) \) of every rectifiable curve in \( K \) is rectifiable with respect to the Euclidean metric, and therefore \( p(\mathrm{Rec}(K)) \subset \mathcal{L} \).
    \par According to Lemma \ref{lem:d_0_is_distance_reestated} the map \( p \) is local isometry on \( S \), hence it preserves the lengths of every rectifiable curve \( \gamma \) mapped into \( \mathrm{int}(\mathcal{D}) \), that is, 
    \[ 
        l_0(\gamma) = l_E(p(\gamma)).
    \]
    But every curve on \( \mathcal{D} \) can be arbitrarily well approximated, in the \(C^1 \) topology, by a polygonal curve on \( \mathrm{int}(\mathcal{D}) \), and length preservation follows.
    \par From length preservation we now get that \( p\rvert_{\overline S} \) is a non-expanding map, that is
    \[
        d_0(x, y) \geq d_\mathcal{L}(p(x), p(y)), 
    \]
    but since \( \overline{S} \) and \( \mathcal{D} \) are compact, we conclude that \( p\rvert_{\overline S} \) is a global isometry, as wanted.     
\end{proof}
\begin{theorem}\label{thm:K_0_as_gluing_metric_space}
    Suppose \( K^\circ \) has \( k \) distinct connected components. Then the space \( K_0 \) is isometric to the metric quotient 
    \[
        ^{\textstyle \left(\displaystyle{\coprod_{i=1}^k}(\mathcal{D}, d_\mathcal{L})\right)}\!\!\!\Big/_{\!\!\textstyle\sim \ }
    \]
    where \( \sim \) is the equivalence relation identifying corresponding boundary components.
\end{theorem}
\begin{proof}
    We denote by \( \widehat d \) the metric on the disjoint union \( X :=\coprod_{i=1}^k(\mathcal{D}, d_\mathcal{L}) \).
    It is defined as
    \[
        \widehat d(x,y) := \begin{cases}
            d_\mathcal{L}(x,y), \text{ if } x, y \text{ belong to the same copy of } \mathcal{D}; \\
            \infty, \text{ otherwise}.
        \end{cases}
    \]
    In general, for an equivalence relation \( R \) on the disjoint union, the semi-metric \(d_R\) ``identifying'' the \( R \)-related points (by assigning distance zero to them) is given by
    \[
        d_R(x,y) := \inf\left\{\sum_{i=0}^q\widehat d(x_i, y_i);\  x_0 = x, y_q = y \text{ and } y_i R x_{i+1} \text{ for } i = 1, \cdots, q-1\right\},
    \]
    but it is clear that in our situation such a semi-metric has the equivalent, simpler description
    \[
        d_\sim(x,y) = \inf\{d_\mathcal{L}(x, m_1) + d_\mathcal{L}(m_2, y); \  m_1 \sim m_2\},
    \]
    since it is possible to change from one copy of the table to the other by ``passing'' through the boundary of \( \mathcal{D} \) only once. 
    We note moreover that \(d_\sim (x,y) = 0 \) if and only if \( x \sim y \), so the only pairs of distinct points \( (x, x') \) whose distance is zero are  those where \(x\) and \(x'\) correspond to the same boundary point. 
    In particular, since \(X\) is compact the quotients \(K \approx X \big/\!\!\sim\) and \( X\big/d_\sim \) are homeomorphic\footnote{
       For a semi-metric \(d\) on \(X\), we denote by \( X\big/d\) the set of equivalence classes of the relation \(d(x,y) = 0\). For a general relation \(R\) on \(X\), it is not always true that the quotients \(X \big/R\) and \(X\big/d_R\) are the same. Even when they coincide as sets, in general one can only guarantee that the topology on \(X\big/d_R\) is coarser, i.e., \(\mathrm{id}\colon X\big/R \to X\big/d_R\) is continuous \cite{burago2022course}.
    }. 
    \par We need to show that \( \mathrm{id}\colon (K,d_\sim) \to (K, d_0 ) \) is an isometry.
    Due to Lemma \ref{lem:sheet_isometry}, we already know that the metrics coincide for any two points in the same copy of \(\mathcal{D}\), so it is enough to consider \(\widetilde x,\widetilde y\) lying in distinct connected components of \(X\). 
    Let \(x, y\) be the corresponding projected points onto \(\mathcal{D}\), and \([x,y]_\mathcal{L}\) the geodesic of \( (\mathcal{D}, d_\mathcal{L}) \) joining them. 
    There are two cases to analyse.
    \par \emph{Case I: \([x,y]_\mathcal{L}\) meets the boundary \(\partial\mathcal{D}\).}
    In this situation we have \(d_\sim(\widetilde x,\widetilde y) = l_E([x,y]_\mathcal{L}) = d_\mathcal{L}(x,y)\).
    Indeed, given any point \( \widetilde m \) in the boundary \(\mathcal{D}\), we have
    \[
        d_\mathcal{L}(\widetilde x, \widetilde m) + d_\mathcal{L}(\widetilde m, \widetilde y) =  d_\mathcal{L}(x, m) + d_\mathcal{L}(m, y) \geq d_\mathcal{L}(x,y),   
    \]
    with equality holding exactly when \( m \) is a point in the segment \([x,y]_\mathcal{L}\).
    Thus the infimum \( d_\sim (\widetilde x, \widetilde y) = d_\mathcal{L}(x, y)\) is achieved thanks to the assumption that \([x,y]_\mathcal{L} \cap \partial\mathcal{D} \neq \emptyset\). 
    \par On the other hand, given any \( m \in x,y]_\mathcal{L} \cap \partial\mathcal{D} \), we may lift the geodesic segments \( [x, m]_\mathcal{L}\) and \( [m, y]_\mathcal{L}\) isometrically to geodesics \( [\widetilde x, m]_0\) and \( [m, \widetilde y]_0 \) of \(K_0\).
    Hence the path of minimal length \( [\widetilde x, \widetilde y]_0 \) has length at most \(l_0([\widetilde x, m]_0) + l_0([\widetilde x, m]_0) = d_\mathcal{L}(x, y) \).
    But the projection map is non-expanding, that is, \(d_\mathcal{L}(x,y) \leq d_0(\widetilde x, \widetilde y \), thus \( d_0(\widetilde x, \widetilde y) = d_\mathcal{L}(x, y) = d_\sim (\widetilde x, \widetilde y)\), as we wanted.
    \par \emph{Case II: \([x,y]_\mathcal{L} \cap \partial\mathcal{D} = \emptyset \).}
    If the intersection is empty, then the geodesic \([x,y]_\mathcal{L}\) is a straight line contained entirely in \( \mathrm{int}(\mathcal{D})\). 
    We adopt a functional approach. 
    As in Section \ref{ssec:periodic_orbits}, for each \( i = 1, 2, \cdots, m \), we consider functions \( \tau^i_x\colon \T_i \ni s \mapsto \|x - \gamma_i(s)\| \) and the length functional
    \begin{align*}
        L_{x,y}^i\colon \T_i &\to \R_+ \\
        s &\mapsto \tau^i_x(s) + \tau^i_y(s).
    \end{align*}
    Remark that, by construction, \(d_\sim(\widetilde x, \widetilde y) = \inf_s\{L_{x,y}^i(s)\}\).
    \par Now, the paths of minimal length between \(x\) and \(y\) with a single element in the scatterer \( \Gamma_i \) are the critical points of \( L_{x,y}^i \), which are exactly the elements \( m = \gamma_i \) where the entry and exit angles are the same.
    So the length minimisers are the billiard trajectories of the form \([x,m]_\mathcal{L} \ast [m, y]_\mathcal{L}\).
    Note that the collision is regular, as the grazing collisions have already been considered in Case I.
    Thus this segment of billiard trajectory is well approximated by the proper segments of geodesics \( \gamma_\epsilon \) of \( K_\epsilon \).
    Hence 
    \[
        d_0(\widetilde x, \widetilde y) = \lim_{\epsilon \to 0}l_\epsilon (\gamma_\epsilon) = [x,m]_\mathcal{L} + [m, y]_\mathcal{L} = d_\sim(\widetilde x, \widetilde y),
    \]
    which concludes.
\end{proof}
The ``seam'' of the above construction, i.e., the subset of \(K_0\) corresponding to the boundary points, where the identification happens, is exactly set \(B(K)\), as defined earlier. 
We end the section with two immediate corollaries.
\begin{corollary}\label{cor:isometry-preserves-boundary}
Let \(K_0 = \coprod_{i=1}^k (\mathcal{D}, d_\mathcal{L})/\!\!\sim\) and \(K'_0 = \coprod_{i=1}^k (\mathcal{D}', d'_{\mathcal{L}})/\!\!\sim\) be metric gluings of Sinai billiard tables. 
If \(F\colon (K_0, d_0) \to (K'_0, d'_0)\) is an isometry, then the tables \((\mathcal{D}, d_\mathcal{L})\) and \((\mathcal{D}', d'_\mathcal{L})\) are also isometric.
\end{corollary}
\begin{proof}
An isometry preserves distances and consequently must map points where geodesic bifurcation happens to points with the same property\footnote{
    Putting it other way, if for any direction there is a single geodesic emanating from \(x\), the same must be true for \(F(x)\).}.    
Since the points were geodesics bifurcate are exactly those in the seam, we conclude that \(F\)
must send the set \(B(K)\) exactly onto \(B(K')\).  
But \(K_0\setminus B(K)\) is the disjoint union of open ``sheets'', each isometric to \((\mathcal{D},d_\mathcal{L})\).
Since \(F\) carries sheets to sheets, it induces on each one a distance‐preserving bijection with the corresponding sheet in \(K'_0\), i.e.\ an isometry \(\mathcal{D}\to\mathcal{D}'\).  
\end{proof}
\par Recall that a metric space is called non-positively curved, or locally CAT(0), if every point has a neighbourhood on which every geodesic triangle satisfies the CAT(0) inequality.
\begin{corollary}
    \( (K, d_0) \) is a non-positively curved metric space.
    The universal covering  \( \tilde K \), equipped with the length metric induced by \(d_0\), is a contractible Hadamard space, that is, a contractible and complete (globally) CAT(0) geodesic space. 
\end{corollary}
\begin{proof}
    First, we observe that \( (\mathcal{D}, d_\mathcal{L}) \) is locally CAT(0).
    Indeed, if a point belong to the table's interior, then it has a neighbourhood isometric to an open subset of \( \R^2 \). 
    Otherwise, choose a neighbourhood sufficiently small so that it does not completely contain any scatterer.
    If a geodesic triangle \( \Delta \) in this neighbourhood does not contain any segment of \( B(K) \) as one of its sides, then the distance between any points is the same as the distance between the corresponding comparison points in \( \R^2 \).
    Otherwise, CAT(0) inequality is clearly satisfied (see Figure \ref{fig:comparison_triangles_table} below).
    \begin{figure}[!ht]
        \centering
        \includegraphics[width=.8\textwidth]{Triangle.pdf}
        \caption{A geodesic triangle with a side in a scatterer.}        \label{fig:comparison_triangles_table}
    \end{figure}
    
    \par It then follows from iterated applications of Reshetnyak's Gluing Theorem \cite[Theorem 9.1.21]{burago2022course} that \( (K, d_0) \) is a non-positively curved metric space.
    Finally, applying the Cartan-Hadamard Theorem for locally CAT(0) spaces (cf. \cite[Theorem II.4.2]{bridson2011metric} gets us the assertion about the universal covering.
\end{proof}

\subsection{Coarse geometry of \( (K, d_0) \)}
    \label{section:coarse_geometry}
    \textbf{ 
    For simplicity, from now on we will assume that \( (K, d_0) \) consists of two copies of \( (\mathcal{D}, d_\mathcal{L}) \) glued along the boundary. 
    }
    
    Let \(\wK_0\) be the universal covering of \(K\).
    For simplicity, we denote the lift of the metric \(d_\epsilon\) to \(\wK_0\) by \(d_\epsilon\) as well, and the same for lift of the Riemannian tensor \(g_\epsilon\), when \(\epsilon > 0\).
    We write \( \wK_\epsilon \) to denote the space \((\wK, d_\epsilon)\).
    The canonical projection \(\wK_\epsilon \to K_\epsilon\) is then a smooth covering and a local isometry. 
    \par The space of (parameterised) geodesics of \(\wK_\epsilon\), \(\epsilon \in [0,1]\), is the subset \( \widehat{\cG}_\epsilon \subset C(\R, \wK)\) consisting of unit speed parameterisations\footnote{
        In a general metric space like \(\wK_0\), a parameterisation \(\gamma\) is unit-speed if \(d(\gamma(t), \gamma(s)) = |t-s|\) for every \(t,s \in \R\). 
    } 
    of geodesics \(\gamma\colon \R \to \wK_\epsilon\).
    The space of unparameterized geodesics is the quotient of \(\widehat{\cG}_\epsilon\) by the relation of reparameterisation.
    More specifically, \(\gamma_1 \sim \gamma_2\) if and only if there is \(T \in \R\) such that \(\gamma_1(t) = \gamma_2(t + T)\) for every \(t \in \R\), and we set
    \[
        \cG_\epsilon := \widehat{\cG}_\epsilon / \sim = \{ [\gamma]; d_\epsilon(\gamma(t), \gamma(s)) = |t-s| \text{ for every } t,s \in \R \},
    \] 
    equipped with the compact-open topology inherited from \( C(\R, \wK) \). 
    We remark that, in the Riemannian cases (i.e., for \(\epsilon > 0\)), there is a natural identification between \(\widehat{\cG}_\epsilon\) and \(T^1\wK_\epsilon\) given by
    \[
        \gamma \mapsto (\gamma(0), \dot\gamma(0)),
    \]
    and in this case the space \(\cG_\epsilon\) is simply the space of orbits of the geodesic flow acting on \(T^1\wK_\epsilon\).
    While there is no unit tangent bundle of \(\wK_0\) and therefore no such identification, we can still define an abstract geodesic flow on \(\wK_0\) as the \(\R\)-action on \(\widehat{\cG}_0\) by precomposition of translations:
    \begin{align*}
        \mathbf{g}_0\colon \widehat{\cG}_0 \times \R &\to \widehat{\cG}_0 \\
        (\gamma, t) &\mapsto \gamma(t + \cdot).
    \end{align*}
    \par For the Riemannian case it follows from the Anosov property of the geodesic flow on \(K_\epsilon\) that \(\wK_\epsilon\) is a Gromov-hyperbolic space (cf. \cite[Appendix B]{guillarmou_geodesic_2022} and \cite{knieper_new_2012}).
    Thus, every \(\wK_\epsilon\) has a well-defined boundary at infinity \(\bK{\epsilon}\), consisting of equivalence classes of geodesic rays which remain from a bounded distance of each other.
    Such space can be equipped with the cone topology, making the union \( \wK_\epsilon \cup \bK{\epsilon}\) into a compact space (see \cite[Chapter II.8]{bridson2011metric}. for more details).
    Moreover, since \(\wK_0\) is a proper metric space, it follows that \(\wK_\epsilon\) is a visibility space \cite[Lemma III.3.2]{bridson2011metric}. 
    What this means is that for any two distinct points \(\alpha, \omega \in \bK{\epsilon}\) there is a unique geodesic on \(\wK_\epsilon\) joining them, or, more globally, that we have a homeomorphism
    \[
        \cG_\epsilon \approx \bbK{\epsilon} := (\bK{\epsilon} \times \bK\epsilon) \setminus \mathrm{Diag}.
    \]
    In what follows, we will see that in the limit case \(\bK{0}\) has much of the same properties.
    \subsubsection{The space \(\cG_0\) and the abstract geodesic flow} Based on the relation between geodesics of \(\wK_0, \ K_0\) and \((\mathcal{D}, d_\mathcal{L})\), we classify geodesics in terms of their projections on the table \(\mathcal{D}\).
    First, we associate to a geodesic a set of collision points, which are the points at which the corresponding segment of trajectory on \(\mathcal{D}\) bounces at a scatterer.
    \begin{defi}[Collision point]
        For a geodesic \( \gamma \) in \( K_0 \),  the set 
        \[
            B(\gamma) := \partial(\gamma(\R) \cap B(K)) 
        \]
        is the set of \emph{collision points} of \(\gamma\).
    \end{defi}
    Recall that \(B(K)\) consists of closed curves on \(K\) whose image under the projection \(p\colon E \to \mathbb{T}^2\) are the scatterers on \(\mathcal{D}\).
    The intersection of a segment of \(\gamma\) with \(B(K)\) could be a single point \textemdash corresponding under \(p\) to a bouncing of the associated trajectory on \(\mathcal{D}\)\textemdash or an entire segment of a scatterer on the table, in which case the geodesic does not correspond to any billiard trajectory.
    This motivates the following terminology.
    \begin{defi}[Geodesic trichotomy]
        We will say a collision point \( x \in K_0\) is \emph{degenerate} if the connected component of \( \gamma(\R) \cap B(K) \) to which it belongs has positive length, i.e., is not an isolated point.
        \par A non-degenerate point is said to be \emph{regular} or \emph{grazing} if the corresponding collision of \( p(\gamma) \) is regular or grazing, respectively.
        \par Similarly, given a geodesic \(\tilde\gamma\) on \(\wK_0\), let \(\gamma\) be its projection on \(K_0\).
        Then the set \(B(\tilde\gamma)\) is defined as the pre-image of \(B(\gamma)\) under the canonical projection, and \(\partial B(\tilde\gamma)\) is the set of \emph{collision points} of \(\tilde\gamma\).
        A collision point of \(\tilde\gamma\) will be called regular, grazing or degenerate when its projection on \(K_0\) has the corresponding property. 
        \par We will say a geodesic is \emph{degenerate} if it contains a degenerate collision point.
        A non-degenerate geodesic is \emph{grazing} if it contains a grazing point. 
        Otherwise it is \emph{regular}.
        Let
        \[
            \greg := \{ [\gamma] \in \cG_0; \gamma \text{ is regular} \},
        \]
        and denote by \(\cG(\mathcal{D})\) the closure of \(\greg\) in \(\cG_0\) with respect to the compact-open topology.
        We say \(\cG(\mathcal{D})\) is the space of \emph{admissible} geodesics.
    \end{defi}
    The non-degenerate geodesics are the ones which project to honest billiard trajectories on \(\mathcal{D}\), while the degenerate geodesics project to curves which are tangent to a scatterer \(\Gamma_i\) for some positive time interval \([t, t+\epsilon]\).
    In particular, degenerate geodesics can not be the limit of regular ones, and therefore \(\cG(\mathcal{D})\) is exactly the set of non-degenerate geodesics on \(\wK_0\). 
    The following is immediate from the definitions just exposed.
    \begin{lemma}
        There is an homeomorphic correspondence \(\cG(\mathcal{D}) \to \Omega/\Psi \) between the space of admissible geodesics and the orbit space of the billiard flow.
    \end{lemma}
    
    If we lift \(\cG(\mathcal{D})\) and consider the space \(\widehat{\cG}(\mathcal{D})\) of \emph{unparameterised admissible geodesics}, then the situation is as follows.
    In \(K_0\), branching of geodesics can occur only at grazing collisions: at points in \(T^1K^\circ\), the geodesics \(\gamma_{\epsilon;v}\) converge uniformly to a well-defined geodesic \(\gamma_{0;v}\). 
    \par Therefore, similar to the Anosov case, at least at the set \(\cR := \pi_K^{-1}(\Omega^\mathrm{reg}) \subset T^1K\) there is an identification with a subset of geodesics \(\widehat{\cG}_0\) given by
    \[
        \cR \ni (x,v) \mapsto \gamma_{0;v} \in \widehat{\cG}^{\mathrm{reg}} \subset \widehat{\cG}_0,
    \]
    where \(\gamma_{0;v}\) is the unique geodesic of \(\wK_0\) emanating from \(x\) ``in the direction \(v\)'', in the sense that it is uniquely defined by the limit
    \[
        \gamma_{0;v}(t) = \lim_{\epsilon \to 0}\gamma_{\epsilon;v}(t).
    \]
    What is more, using this representation, the restriction of the abstract geodesic flow on \(\widehat{\cG}_0\) to the set of regular geodesics has the more familiar form
    \begin{equation}\label{eq:geodesic_flow_K_0}
    \begin{split}
        \mathbf{g}_0^t\colon \cR &\to \cR \\
        (x,v) &\mapsto (\gamma_{0;v}(t), \dot\gamma_{0;v}(t)) =  \lim_{\epsilon \to 0}\left(\gamma_{\epsilon;v}(t), \dot\gamma_{\epsilon;v}(t)\right),    
    \end{split}
    \end{equation}
    which is well defined for \( t\in \R \).
    Finally, Kourganoff's theorem \cite{kourganoff_anosov_2016} implies that such flow maps are uniformly continuous on compact sets.
    In particular, the flow has a unique continuous extension to \(\overline{\cR} = T^1\wK\), which induces, in particular, a one-to-one identification \(\widehat{\cG}(\mathcal{D}) \approx T^1\wK\).
    We may then restate Kourganouff's result as 
    \begin{theorem}[Kourganoff]\label{thm:kourganoff_on_steroids}
        Let \( \mathcal{D} \) be a Sinai billiard table and \( K \in \cK(\mathcal{D}) \). 
        Then the abstract geodesic flow on \(\widehat{\cG}_0\) restricts to the space of admissible geodesics.
        Such a restriction induces a ``geodesic flow'' \( \mathbf{g}_0 \) on \( K \), which is conjugated to the billiard flow \( \Psi \) on \( \mathcal{D} \).
        More specifically,
        \[
            \pi_K \circ \mathbf{g}^t_0(x,v) = \Psi^t\circ\pi_K(x,v)
        \]
        for every \((x,v) \in T^1K\), where \(\pi_K\colon T^1K \to \Omega\) is a homeomorphism whose restriction to \(T^1K^\circ\) agrees with the smooth projection defined in \eqref{eq:projection_flow_equivalence}.
     \end{theorem}

    
               
\subsubsection{The boundary \(\bK{0}\)}
    \par Being a Hadamard space, \( \wK_0 \) has a well-defined circle (boundary) at infinity \(\bK{0}\), defined as the set of equivalence classes of asymptotic rays.
    More explicitly, \(\gamma, \gamma'\colon [0, \infty) \to \wK_0\) are in the same class if \(\sup\{d_0(\gamma(t), \gamma'(t)); t \in [0, \infty)\} < \infty\).
    We denote the class of the ray \(\gamma\) by \(\gamma(\infty)\).
    More generally, if \(\gamma\colon \mathbb{R} \to \wK_0\) is a geodesic, we set its endpoints at infinity to be
    \begin{align*}
        \gamma(\infty) &:= \gamma^+(\infty),\\
        \gamma(-\infty) &:= \gamma^-(\infty),
    \end{align*}
     where \(\gamma^\pm: [0, \infty)\) is the geodesic ray defined by \( \gamma^\pm(t) = \gamma(\pm t). \)
    \par With the appropriate cone topology on \(\bK{0}\), the union \(\wK_0  \cup \bK{0} \) becomes a compactification of \(\wK_0\), homeomorphic to a disc \cite{bridson2011metric}.
    The fundamental group \(\pi_1(K)\) acts on \(\wK_0\) via isometries, and such an action extends continuously to \(\bK{0}\).
    \par Let \(\bbK{0} := \{ (x,y) \in \partial_\infty \tilde{K} \times \partial_\infty \tilde{K}; x \neq y\}\) denote the space of distinct boundary points. 
    There is a well-defined \(\pi_1\)-equivariant map
    \begin{align*}
        \partial\colon \cG_0 &\to \bbK{0}, \\
                    [\gamma] &\mapsto (\gamma(-\infty), \gamma(\infty)).
    \end{align*}       
    \begin{lemma}\label{lem:wK_is_visibility}
        \(\wK_0\) is a visibility space, that is, \(\partial\colon \cG_0 \to \bbK{0}\) is a homeomorphism.   
        Moreover, any two geodesics in \(\greg\) intersect at most once, and transversely.
    \end{lemma}
    \begin{proof}
        \(\partial\) is a continuous surjective closed mapping \cite{bankovic_marked-length-spectral_2017}, so it suffices to show it is injective. 
        Suppose, by contradiction, that is not the case.
        Then, according to the Flat Strip Theorem \cite{bridson2011metric}, the space \(\wK_0\) contains an isometric copy of \(\R \times [0,\epsilon]\).        
        \par Consider the family \( \{\R \times \{s\}\}_{s \in [0, \epsilon]} \) of parallel geodesics on \(\wK_0\).
        These geodesics project isometrically onto a continuous family of billiard trajectories \(\{\eta_s\}\) on \(\mathcal{D}\) (the projection is simply the composition of the canonical projection \(\wK_0  \to K_0\) with the mapping \( p\colon E \to \mathbb{T}^2 \)).
        Now, the trajectories \(\eta_s\) can not consist solely of grazing collisions.
        Thus, if \(\eta_s\) collides regularly with a scatterer \(\Gamma\), for \(s'\) sufficiently close to \(s\), the trajectory
        \(\eta_s\) collides with \(\Gamma\) as well.
        Since the trajectories are parallel and \(\Gamma\) is strictly convex, after the collision the trajectories are dispersed, and therefore grow apart from another, a contradiction. 
        \par Branching of geodesics in \(\cG\) can only occur at grazing collisions.
        Hence, given two distinct geodesics in \(\greg\), a tangential intersection between them would force them to be the same. Therefore any intersection of geodesics in \(\greg\) must be transverse.
        \par Finally, supposing they intersect at a point \(p \in \wK_0\), then the corresponding trajectories define a billiard Jacobi field \cite{wojtkowski_two_1994} starting at the corresponding point on the phase space \(\Omega\).
        Zeros of Jacobi fields are equivalent to focal points of the corresponding wavefronts. 
        Since a wavefront along a regular orbit in a Sinai billiard changes signs at most once, i.e., has at most one focal point,
        it follows that there are no conjugated points for regular orbits for Sinai billiards.
        In particular, the intersecting point must be unique.  
    \end{proof}
    As a consequence, we obtain
    \begin{corollary}
        \(\wK_0\) is a hyperbolic space. 
        In other words, there is \(\delta > 0\) such that every geodesic triangle on \(\wK_0\) has the property that each side is contained on the \(\delta\)-neighbourhood of the other two sides' union.
    \end{corollary}
    \begin{proof}
        This is the content of the Flat Plane Theorem (cf. \cite[Proposition III.H 1.4 and Theorem III.H 1.5]{bridson2011metric}.
    \end{proof}
    \begin{lemma}\label{lem:homeomorphic_boundaries}
        There is a family \(\{\rho_\epsilon\}_{\epsilon \in [0,1]}\) of homeomorphisms 
        \[
            \rho_\epsilon\colon \bK{0} \to \bK{\epsilon}
        \]
        which is equivariant with respect to the isometric action of \(\pi_1(K)\) on \(\bK{\epsilon}\), i.e., 
        \[
            \rho_\epsilon(\alpha\cdot\gamma(\infty)) = \alpha\cdot\rho_\epsilon(\gamma(\infty)) \quad\forall \gamma(\infty) \in \bK{0}, \forall \alpha \in \pi_1(K),
        \]
        and satisfies
        \[
            \gamma(\infty) = \lim_{\epsilon \to 0}\rho_\epsilon(\gamma(\infty)),
        \]
        where the limit is taken on the space of curves \(C([0, \infty], K)\), equipped with the compact-open topology.
    \end{lemma}
    \begin{proof}
    First, take \(\epsilon > 0.\) 
    Then, there is a Hölder continuous homeomorphism \(H_\epsilon\colon T^1\wK_\epsilon \to T^1\wK_1\) which semiconjugates the geodesic flows \cite{ghys_flots_1984}, namely, an orbit-equivalence. 
    In other words, \(H_\epsilon\) is an orbit equivalence: given a geodesic \(\gamma \in \cG_\epsilon\), there is a unique geodesic \(H_\epsilon\gamma \in \cG_1\) corresponding to \(\gamma\), and such correspondence is equivariant with respect to the natural action of \(\pi_1(K)\) on the space of geodesics, induced by deck transformations.
    \par The geodesic ray \(\gamma^+ \colon [0,\infty) \to K_\epsilon\) defines a unique boundary point \(\gamma(\infty)\). 
    We define the homeomorphism $\rho_\epsilon^{-1}$ by mapping the forward point at infinity of \(\gamma\) to the forward point at infinity of $H_\epsilon\gamma$:
    \begin{align*}
        \rho_\epsilon^{-1}\colon \bK{\epsilon} &\to \bK{1},\\
                \gamma(\infty) &\mapsto (H_\epsilon\gamma)^+(\infty).
    \end{align*}    
    This extension is well-defined because \(\gamma'\) is another geodesic ray in the class \(\gamma(\infty)\) if and only if \((\gamma',\dot\gamma')\) is a section of strong stable bundle \(E^{ss}_\gamma\) of \(g_\epsilon\) along \(\gamma\).
    Since \(H_\epsilon\) is Hölder, it preserves the invariant foliations, hence \(H_\epsilon\gamma'(\infty) = H_\epsilon\gamma(\infty)\) (that is to say, points in the geodesic \(H_\epsilon\gamma'\) belong to the stable foliations of \(H_\epsilon\gamma\) as well).
    \par As for the continuity of the map \(\rho_\epsilon^{-1}\) on the boundary, we note that the cone topology on \(\bK{\epsilon}\) coincides with the compact-open topology on the space of geodesic rays~\cite[Chapter II.8]{bridson2011metric}. 
    Moreover, if \(\sigma\) is another geodesic sufficiently close to \(\gamma\) along a compact subset of \([0,\infty)\), then \(H_\epsilon\sigma\) will be close to \(H_\epsilon\gamma\) along a corresponding compact, implying that such application is continuous in the compact-open topology of the space of geodesics (and, in particular, geodesic rays) as we wanted.
    \par Reversing the roles of the indices \(\epsilon\) and \(1\) and applying the exact same construction provides a continuous inverse \(\rho_\epsilon\), which is the family we seek. 

    \par For the case \(\epsilon = 0\), we recall that the Equation \eqref{eq:bound_geom} gives a uniform bound for the tensor norms
    \[
        \|g_\epsilon - g_\delta\|_{C^r(K)} < C, 
    \]
    which means, in particular, that
    \begin{equation}\label{eq:inner_prod_constraints}
        \frac{1}{\lambda}g_\epsilon < g_\delta < \lambda g_\epsilon
    \end{equation}   
    for every \(\epsilon, \delta > 0\) and for \(\lambda:= 1+C\).
    As the projections \( \wK_\epsilon \to K_\epsilon\) are local isometries, it follows that the inequalities in \eqref{eq:inner_prod_constraints} are also true for the universal coverings \(\wK_\epsilon\).
    Consequently, for every \(\epsilon > 0\), any geodesic in \(\cG_\epsilon\) is a quasi-geodesic\footnote{
        More specifically, \(\gamma_\epsilon\) is a \((\lambda, 0)\)-quasi-geodesic shadowing \(\gamma_1\)
    } on \(\wK_1\). 
    More explicitly, what we have are inequalities of the form
    \[
        \frac{1}{\lambda}|t - s| \leq d_1(\gamma_\epsilon(t), \gamma_\epsilon(s)) \leq \lambda|t - s|
    \]
    for any geodesic \(\gamma_\epsilon\colon \mathbb{R} \to \wK_\epsilon\) in the space \(\cG_\epsilon\), and any \(t,s \in \R\).
    \par We may now apply Morse-Klingenberg lemma (cf. \cite[Theorem III.1.7]{bridson2011metric}) that we can find a constant \(R = R(g_1, \lambda) > 0\), such that any geodesic in \(\cG_\epsilon\), \(\epsilon \in (0, 1]\) (seen as a curve in \(\wK_1\)) is in a \(R\)-neighbourhood of a geodesic in \(\cG_1\).
    In other words, given \(\gamma_\epsilon \in \cG_\epsilon\), there is \(\gamma_1 \in \cG_1\) such that
    \[
        \gamma_\epsilon(\mathbb{R}) \subset B_1(\gamma_1; R) := \coprod_{t \in \mathbb{R}}\{x \in \wK_1; d_1(x, \gamma_1(t)) < R\}.
    \]
    What is more, since the geodesic flow is \(\wK_1\) hyperbolic, there is at most one \(g_1\)-geodesic which such property, since no two geodesics could stay at finite distance from one another for all time.
    \par In particular, any geodesic \(\gamma\) in \(\cG_0\), being a uniform limit of sequences \(\gamma_n \in \cG_{\epsilon_n}\) on compacta, is necessarily also a quasi-geodesic of \(\wK_1\) laying in the \(R\)-neighbourhood of some \(\gamma_1 \in \cG_1\).
    \par Now, a quasi-geodesic ray \(c\colon [0, \infty) \to \wK_1\) defines a unique point at infinity \(\xi(c) \in \bK{1}\) \cite[Lemma III.H.3.1]{bridson2011metric}.
    We can then define a continuous bijection
        \begin{align*}
            \rho_0\colon \bK{0} &\to \bK{1} \\
            \gamma(\infty) &\mapsto \xi(\gamma^+).
        \end{align*}
    To define the inverse of this map, we note that, given \(\gamma(\infty) \in \bK{1}\), there is a unique geodesic \(\gamma_0 \in \cG_0\) obtained as the limit 
    \[
        \gamma_0 = \lim_{\epsilon \to 0}H_\epsilon^{-1}\gamma.
    \]
    In fact, the family \(\overline{\{H_\epsilon^{-1}(\gamma)\}}\) lies in the compact set
    \[
        \coprod_{\epsilon \in [0,1]}\cG_\epsilon \approx \coprod_\epsilon \bK{\epsilon}\times\bK{\epsilon},
    \]
    so that it contains all its accumulation points.
    If \(\gamma_0, \gamma_0'\) are two such accumulation points, then by construction they must lie at a finite \(g_1\)-distance of at most \(2R\) from one another, since they are limits of sequences of curves with such properties. 
    Now, from \eqref{eq:bound_geom} we conclude that they are also at bounded \(d_0\) distance from one another, and, consequently, define the same points at the boundary \(\bK{0}\):
    \[
    \begin{split}
        \gamma_0(\infty) &= \gamma_0'(\infty), \\
        \gamma_0(-\infty) &= \gamma_0'(-\infty).
    \end{split}
    \]
    Since, by Lemma \ref{lem:wK_is_visibility}, \(\wK_0\) is a visibility space, it follows that \(\gamma_0 = \gamma_0'\).
    Then the mapping
    \[
        \rho_0^{-1}\colon \gamma(\infty) \mapsto  \gamma_0(\infty) \in \bK{0}
    \]
    is a continuous inverse to \(\rho_0\), which concludes.
    \end{proof}

    \begin{remark}\label{rem:anosov_top_equivalences}
        The topological (or orbit) equivalences \(H_{\epsilon}\) are essentially unique, due to the Anosov Shadowing Theorem (more accurately, they are \emph{transversally} unique, cf. \cite[Theorem 5.4.1]{fisher2019hyperbolic}).
        The choice of homeomorphism does not change the equivalence between orbits, just the reparameterisation function \(\tau\) appearing in the relation
        \[
            H_\epsilon \mathbf{g}_\epsilon^t(z) = \mathbf{g}_1^{\tau(z, t)}(H_\epsilon(z)).
        \]
        One possibility is to choose \(H_\epsilon\) is the following construction, from \cite[Appendix B, Theorem B.1]{guillarmou_geodesic_2022}.
        Given \(v \in T^1\wK_\epsilon\), we write 
        \[
            (v^-_\epsilon, v^+_\epsilon) = (\gamma_{\epsilon; v}(-\infty), \gamma_{\epsilon; v}(\infty)) \in \bbK{}. 
        \]
        The homeomorphism \(H_\epsilon\colon T^1\wK_\epsilon \to T^1\wK_1\) is then obtained by setting \(H_\epsilon(v) = w\), where \(w \in T^1\wK_1\) is the unique vector such that
        \[
        (v^-_\epsilon, v^+_\epsilon) = (w^-_1, w^+_1) \quad \text{and} \quad B^1_{v^+_\epsilon}(\pi(v), \pi(w)) = 0.        
        \]
        Here \(\pi\colon T^1\wK \to \wK\) is the canonical projection and 
        \[
            B^1_{v^+_\epsilon}(\pi(v), \pi(w)) = \lim_{t \to \infty} d_1(\gamma_{1;v}(0), \gamma_{\epsilon;v}(t))-d_1(\gamma_{1;w}(0), \gamma_{\epsilon;v}(t))
        \]
        is the Busemann function of the boundary point \(v^+_\epsilon\) with respect to the metric \(g_1\).
        \par With such choice, the family \(\{H_\epsilon^{-1}\}\) becomes an uniformly equicontinuous deformation of the \(g_1\)-geodesics on \(\wK\).
        Indeed, if that were not the case there would be \(\eta> 0\) and sequences \((v_n, u_n) \in T^1K_1, \epsilon_n \to 0\) such that
        \[
            d_1(v_n, u_n) < \frac{1}{n} \quad \text{and} \quad d_\epsilon(H_\epsilon^{-1}(v_n), H_\epsilon^{-1}(u_n)) > \epsilon_0.
        \]
        The common limit \(\lim_n v_n = \lim_nu_n = z \in T^1K_1\) defines a unique geodesic \(\gamma \in \widehat{\cG}_0\).
        Moreover, by construction, the geodesics defined by the initial conditions \(H_\epsilon^{-1}(v_n)\) and \(  H_\epsilon^{-1}(u_n)\) have the same endpoints at infinity than \(v_n\) and \(u_n\), and must therefore both converge to \(\gamma\) on every compact subset as \(\epsilon \to 0\), a contradiction.
    \end{remark}
    \begin{remark}
        A simpler--but considerably less informative--proof of Lemma \ref{lem:homeomorphic_boundaries} would be just noticing that \[\mathrm{id}\colon \wK_0 \to \wK_\epsilon
        \]
        is a family of quasi-isometries, and apply \cite[Theorem III.H.3.9]{bridson2011metric}.
    \end{remark}
    \begin{corollary}\label{cor:dense_periodic_liftings}
        The boundary points \(\gamma(\infty)\) corresponding to lifts \(\gamma\) of periodic geodesics on \(K_0\) for a dense set on \(\bK{0}\).
        Moreover, the set of periodic points\footnote{
            i.e., lifts of periodic geodesics from \(K_0\).
        } 
        in the space of geodesics \(\widehat{\cG}_0\) is dense, and the abstract geodesic flow \(\mathbf{g}_0\colon \widehat{\cG}_0 \to \widehat{\cG}_0\) is transitive \(\mod{\pi_1(K)}\).
        In other words, the flow \(\mathbf{g}\) is weakly hyperbolic, in the sense of Ballmann \cite{ballmann1995lectures}.
    \end{corollary}
    \begin{proof}
        Each \(H_\epsilon\) is a flow equivalence, hence it maps periodic orbit onto periodic orbit.
        Since the geodesic flow \(\mathbf{g}_1\) is Anosov, its set of lifted periodic orbits is dense in \(\bK{1}\).
        As every periodic orbit on \(\wK_0\) is obtained as a limit \(\lim_{\epsilon \to 0}H_\epsilon\gamma\) for some periodic orbit \(\gamma\) on \(\wK_1\), and density on \(\bK{0}\) follows.
        Consequently, the periodic points of \(\cG_0 \approx \bbK{0}\), which are simply the points \((\gamma(-\infty), \gamma(\infty))\) for periodic \(\gamma\), are also dense.
        \par Finally, each geodesic flow \(\mathbf{g}_\epsilon\) is Anosov, and therefore transitive\(\!\mod{\pi_1(K)}\) on \(K_\epsilon\): any two open non-empty subsets \(U, V\) in \(\cG_\epsilon\) there is \(t \in \mathbb{R}\) and \(\alpha \in \pi_1(K)\) such that 
        \[
            \mathbf{g}_\epsilon^t(U) \cap \alpha(V) \neq \emptyset.
        \]
        Equivariance guarantees such property passes to the limit case.    
        \end{proof}

    
    \subsubsection{The Liouville current on \(\bK{0}\)}
    We fix a family of homeomorphisms \(H_\epsilon\colon \cG_0 \to \cG_\epsilon\) as in the last section, and identify all the boundaries \(\bK\epsilon\) with \(\bK{0} \approx \mathbb{S}^1\).  
    Given \(\gamma_0 \in \cG_0\), we write \(\gamma_\epsilon := H_\epsilon \gamma_0\), and may assume, without loss of generality, that \(\gamma_\epsilon(\infty) = \gamma(\infty) \in \mathbb{S}^1\) for all \(\epsilon \in [0,1]\). 
    In particular, \(\gamma\) is a quasi-geodesic on every \(\wK_\epsilon\) and thus define a unique point at infinity in \(\bK{\epsilon}\), which, under our former identifications, is simply \(\gamma(\infty)\).
    
    \par For \(\epsilon \in [0,1]\) and \(x, y, z, w \in \wK_\epsilon\), write 
    \[
        E_{\epsilon}(x, y, z, w) := d_\epsilon(x, y) + d_\epsilon(z, w) - d_\epsilon(x, w) - d_\epsilon(z, y).
    \]
    Following \cite{otal_sur_1992, bourdon_sur_1996}, given four distinct points \(\zeta_1, \zeta_2, \zeta_3, \zeta_4 \in \bK{0}\), we define their cross-ratio as
    \[
        \crat{0}(\zeta_1, \zeta_2, \zeta_3, \zeta_4) := \lim_{(x,y,z,w) \to (\zeta_1, \zeta_2, \zeta_3, \zeta_4)} e^{\frac{1}{2}E_0(x, y, z, w)}.
    \]
    This quantity does not depend on the sequences approaching the given points at the boundary, hence we can rewrite it as
    \begin{equation}\label{eq:cr_rat_def}
        \crat{\epsilon}(\zeta_1, \zeta_2, \zeta_3, \zeta_4) = \lim_{t \to \infty} e^{\frac{1}{2}E_0(\sigma_1(t), \sigma_1(-t), \gamma_3(t), \gamma_4(t))}
    \end{equation}
    where \(\gamma_i \in \cG_0\) are geodesics such that \(\gamma_i(\infty) = \zeta_i\), for \(i = 1,2,3,4\).
    Moreover, since \(d_0(x,y) := \lim_{\epsilon \to 0}d_\epsilon(x,y)\), we get
    \begin{align*}
        \lim_{t \to \infty}E_0(\gamma_1(t), \gamma_2(t), \gamma_3(t), \gamma_4(t)) &= \lim_{t \to \infty}\lim_{\epsilon \to 0}E_\epsilon(\gamma_1(t), \gamma_2(t), \gamma_3(t), \gamma_4(t)) \\
        &= \lim_{\epsilon \to 0}\lim_{t \to \infty}E_\epsilon(\gamma_1(t), \gamma_2(t), \gamma_3(t), \gamma_4(t))        
    \end{align*}
    where the limit operators commute due to the uniform convergence of \(d_\epsilon \to d_0\) on compact sets.
    Consequently, we have
     \begin{equation}\label{eq:cr_rat_lim}
        \crat{0}(\zeta_1, \zeta_2, \zeta_3, \zeta_4)
        = \lim_{\epsilon \to 0} \crat{\epsilon}(\zeta_1, \zeta_2, \zeta_3, \zeta_4),
    \end{equation}
    where \(\crat{\epsilon}\) is defined as in Equation \eqref{eq:cr_rat_def}, using \(E_\epsilon\) instead.
    Remark that we are considering the same geodesics \(\gamma_i \in \cG_0\), seen as quasi-geodesics on \(\wK_\epsilon\) approaching the same point \(\zeta_i\) at infinity, under the identifications expressed at the beginning of this section.
    \par The following is straightforward to conclude from the definition.
    \begin{prop}\label{pps:crat_properties}
        The mappings \(\crat{\epsilon}\colon (\bK{\epsilon})^4 \to \mathbb{R}, \epsilon \in [0,1]\), are positive functions satisfying
        \begin{enumerate}
            \item \( \crat{\epsilon} \) is invariant under the action of \( \pi_1(K) \) on \( (\bK{\epsilon})^4 \);
            \item \( \crat{\epsilon}(\zeta_1, \zeta_2, \zeta_3, \zeta_4) = \crat{\epsilon}(\zeta_2, \zeta_1, \zeta_4, \zeta_3)^{-1} \);
            \item \( \crat{\epsilon}(\zeta_1, \zeta_2, \zeta_3, \zeta_4) = \crat{\epsilon}(\zeta_3, \zeta_4, \zeta_1, \zeta_2) \);
            \item \( \crat{\epsilon}(\zeta_1, \zeta_2, \zeta_3, \zeta_4)\crat{\epsilon}(\zeta_1, \zeta_3, \zeta_2, \zeta_4) = \crat{\epsilon}(\zeta_1, \zeta_2, \zeta_2, \zeta_4) \);
            \item \( \crat{\epsilon}(\zeta_1, \zeta_2, \zeta_3, \zeta_4)\crat{\epsilon}(\zeta_2, \zeta_3, \zeta_1, \zeta_4)\crat{\epsilon}(\zeta_3, \zeta_1, \zeta_2, \zeta_4) = 1 \).
        \end{enumerate}
    \end{prop}

    We briefly recall some facts about geodesic currents.
    For \( \epsilon > 0 \), the surface \( K_\epsilon \in \mathcal{K}(\mathcal{D}) \) carries a smooth Riemannian metric \( g_\epsilon \). 
    This structure induces a canonical contact form \( \lambda_\epsilon \) on the unit tangent bundle \( T^1 K_\epsilon \), obtained by pulling back the canonical symplectic form from the cotangent bundle.
    The Reeb vector field of \( \lambda_\epsilon \) coincides with the geodesic spray of \( g_\epsilon \), that is, the vector field on \( T^1 K_\epsilon \) generating the geodesic flow.
    The associated Liouville measure \( \int\lambda_\epsilon \wedge d\lambda_\epsilon \) on \( T^1 K_\epsilon \) is a finite measure, invariant under the geodesic flow, and consequently under its time reversal.
    \par All such structures are lifted to \(T^1\wK_\epsilon\) via the local isomorphism defined by the covering map \(\wK_\epsilon \to K_\epsilon\), so that the space \(T^1\wK_\epsilon\) has the corresponding contact and volume forms \(\widetilde{\lambda}_\epsilon\) and  \( \int\widetilde\lambda_\epsilon \wedge d\widetilde\lambda_\epsilon \). 
    Given a rectifiable curve \( \gamma \subset \wK_\epsilon \), its unit tangent lift \( \dot{\gamma} \subset T^1 \wK_\epsilon \) defines the \emph{Liouville current} \( m_\epsilon \) on the space of geodesics \( \mathcal{G}_\epsilon \) by integration against \( d\widetilde\lambda_\epsilon \), assigning mass to Borel sets of geodesics intersecting \( \gamma \) transversely~\cite{wilkinson_anosov}.
    \par More generally, a geodesic current for \(K_\epsilon\) is any Radon measure on the space \(\bbK{\epsilon}\) that is invariant under both the diagonal action of \(\pi_1(K_\epsilon)\) and the flip transformation \( (\zeta_1, \zeta_2) \mapsto (\zeta_2, \zeta_1)\).
    There is an affine one-to-one correspondence between the cones of finite measures on \(T^1K_\epsilon\) invariant under the geodesic flow and the space of geodesic currents with the weak-* topology (cf. \cite{bonahon_bouts_1986}, \cite{furman_coarse-geometric_2002}).
    \par In particular, the contact structure \( \lambda_\epsilon \), the cross-ratio \( \crat{\epsilon} \), and the Liouville current are realted in the following fashion. 
    For measurable rectangles in the space of geodesics determined by intervals \( [\zeta_1, \zeta_2] \) and \( [\zeta_3, \zeta_4] \) in the boundary \( \partial \widetilde{K}_\epsilon \), we have:
    \[
        m_\epsilon\big([\zeta_1, \zeta_2] \times [\zeta_3, \zeta_4]\big) 
        = \left| \ln \crat{\epsilon}(\zeta_1, \zeta_2, \zeta_3, \zeta_4) \right|
        = \int_{\Sigma} d\lambda_\epsilon,
    \]
    where \( \Sigma \subset T^1 \widetilde{K}_\epsilon \) is a transversal to the geodesic flow intersecting exactly the geodesics joining \( [\zeta_1, \zeta_2] \) to \( [\zeta_3, \zeta_4] \), oriented accordingly.

    Such a construction can not possibly be replicated on \(\wK_0\), since the Riemannian tensors degenerate as \(\epsilon \to 0\) and the metric on \(\wK_0\) is not induced by any Riemannian metric. 
    Instead, we define a Liouville current on \(\cG_0\) as a limit.
    
    \begin{defi}[Liouville pre-measure]
    We use the notation \(|\zeta_1, \zeta_2|\) to denote any of the intervals \([\zeta_1, \zeta_2], (\zeta_1, \zeta_2], [\zeta_1, \zeta_2)\) or \((\zeta_1, \zeta_2)\). 
    Consider the set map \(m_0\) defined on the space of rectangles \(|\zeta_1, \zeta_2| \times |\zeta_3, \zeta_4| \subset \bbK{0}\) as   
    \[
        m_0\big(|\zeta_1, \zeta_2| \times |\zeta_3, \zeta_4|\big) 
        := \left| \ln \crat{0}(\zeta_1, \zeta_2, \zeta_3, \zeta_4) \right|.
    \]
    \end{defi}
    
    \begin{prop}
        The mapping \(m_0\) extends uniquely to a current on the space of geodesics \(\cG_0\).
    \end{prop}
    \begin{proof}
        We begin by noting that the family
        \[
            \left\{|\zeta_1, \zeta_2| \times |\zeta_3, \zeta_4|; \ \zeta_i \in \bK{0}\right\} \cup \{\emptyset\}
        \]
        forms a semi-ring of sets.
        Let \(R\) be the ring of sets generated by this semi-ring.
        In other words, \(R\) consists of all finite disjoint unions of rectangles of the form \(|\zeta_1, \zeta_2| \times |\zeta_3, \zeta_4|\), together with \(\emptyset\).
        We extend \(m_0\) to this ring by setting the value of \(m_0\) at \(A = A_1 \cup \cdots \cup A_\ell\), with \(A_i \cap A_j = \emptyset\), to be
        \[
            m_0(A) := \sum_{i=1}^\ell m_0(A_i).
        \]
        This definition is well posed because any two finite disjoint decompositions of a set in \( R \) can be refined to a common disjoint partition into rectangles, and since \( m_0 \) is additive on such rectangles, both decompositions yield the same total as rearrangements of the same sum.        
        \par Moreover, \(m_0\) is a pre-measure on \(R\).
        Indeed, the \(\sigma\)-algebra generated by \(R\) is the Borel \(\sigma\)-algebra of \(\bK{0}^4 \approx \mathbb{T}^4\), where all the currents \(m_\epsilon\) are defined.
        If \(A \in R\) happens to be a countable disjoint union \(A = \bigsqcup_i A_i\), with each \(A_i \in R\), then we may assume, without loss of generality, that \(A\) and each of the \(A_i\) are rectangles.        
        From the approximation
        \begin{equation}\label{eq:m_0_as_limit}
            m_0\big(|\zeta_1, \zeta_2| \times |\zeta_3, \zeta_4|\big)
            = \lim_{\epsilon \to 0}\left| \ln \crat{\epsilon}(\zeta_1, \zeta_2, \zeta_3, \zeta_4) \right|
            = \lim_{\epsilon \to 0} m_\epsilon\big(|\zeta_1, \zeta_2| \times |\zeta_3, \zeta_4|\big),
        \end{equation}
        where the leftmost equality is justified by Equation~\eqref{eq:cr_rat_lim}, and the fact that each \(m_\epsilon\) is a measure, we obtain
        \[
            m_0(A) = \lim_{\epsilon \to 0} m_\epsilon(A)
                   = \lim_{\epsilon \to 0} \sum_i m_\epsilon(A_i)
                   = \sum_i \lim_{\epsilon \to 0} m_\epsilon(A_i)
                   = \sum_i m_0(A_i),
        \]
        where the interchange of limit and sum is justified by the uniform convergence of \(m_\epsilon \to m_0\) on fixed rectangles.
        \par This means \(m_0\) is in fact a pre-measure on the ring \(R\), and it follows from Caratheodory's Extension theorem that it extends to a measure \(m_0\).
        What is more, the pre-measure \(m_0\) is \(\sigma\)-finite, because every \(m_\epsilon\) enjoys such property and the convergence \(m_\epsilon \to m_0\) is uniform. 
        So we conclude \(m_0\) has a unique extension to the Borel \(\sigma\)-algebra of \(\bbK{0}\), which is a Radon measure by construction.
        
        Finally, we note that \(\pi_1\) and flip invariance of \(m_0\) on the sets of \(R\) are direct consequences of Proposition \ref{pps:crat_properties}, which concludes.
    \end{proof}
    In light of the last result, we may define
    \begin{defi}[Liouville current of \(K_0\)]
        The Liouville current \(m_0\) on \(\cG_0\) is the unique current defined on the Borel subsets of \(\bK{0}^4\) by the formula
        \[
            m_0(A) = \lim_{\epsilon \to 0}m_\epsilon(A).
        \] 
        Its pullback under the canonical projection \(\widehat{\cG}_0 \to \cG_0\) defines a Radon measure \(\widehat{m}_0\) on \(\widehat{\cG}_0\), which we call the Liouville current as well.
    \end{defi}
    Remark that the \(\cG_0\) is exactly the orbit space of the abstract geodesic flow \(\mathbf{g}_0\) on  \(\widehat{\cG}_0\), since the orbit \(\mathbf{g}^\R_0(\gamma) = \{\gamma( \cdot + s); s \in \R\}\) is exactly the set of all unit speed reparameterisations. 
    In particular, \(\widehat{m}_0\) is a ``transverse'' measure invariant under both the geodesic flow and the action of \(\pi(K)\) on geodesics.
    
\textemdash positive Radon-Nikodym derivative with respect to \(m_0\), every regular geodesic is a density point, I'm not sure \textemdash then the enriched spectrum would not be necessary, that is, we could forget the larger measure \(m_0\) on the space of geodesics of the surface \(K\), and work only with periodic data from the table \(\mathcal{D}\).

    
      




\section{Spectral Rigidity}\label{sec:spectral_rigidity}
In what follows, we continue with last section's convention that the Kourganoff surface \(K\) is a gluing of two copies of the table \(\mathcal{D})\).
Thus, the interior \( K^\circ \) has two connected components, the closures of which we will denote by \( \mathcal{D}_{\mathrm{up}} \) and \( \mathcal{D}_{\mathrm{down}} \), respectively.
\par Recall that \( \mathcal{C}(K) := \mathrm{Conj}(\pi_1(K)) \) is the set of free homotopy classes on \( K \).
\begin{lemma}\label{lem:unique_minimising_geodesic}
    Every element \( \alpha \in \mathcal{C}(K) \) contains a single closed geodesic curve \( \gamma^\alpha_0 \), which is a curve of minimal \(d_0\)-length among all those of \( \alpha \).
\end{lemma}
\begin{proof}
    We know that each surface \( K_\epsilon \) enjoys such property, because the geodesic flows \( \mathbf{g}_\epsilon \) are all Anosov \cite{klingenberg_riemannian_1974}.
    For a fixed \( \alpha \in \mathcal{C}(K) \) we write \( \gamma_\epsilon^\alpha\) to denote the unique closed geodesic of \( K_\epsilon \) whose free homotopy class is \( \alpha \).
    The family \( \overline{\{\gamma^\alpha_\epsilon\}} \) is a compact subset of \( C^1(\mathbb{S}^1, K) \), by the Arzelà–Ascoli theorem.
    We claim that \( \overline{\{\gamma^\alpha_\epsilon\}}\setminus\{\gamma^\alpha_\epsilon\} \) consists of a single curve \( \gamma^\alpha_0\), with the desired properties.
    \par First notice that for any \( \epsilon, \delta \) sufficiently close, the geodesic flows \( \mathbf{g}_\epsilon, \mathbf{g}_\delta \) are orbit equivalent via a homeomorphism \( H_{\epsilon, \delta}\colon T^1K \to T^1K \), which can be made arbitrarily close to the identity by choosing \( |\epsilon - \delta| \) small enough.
    In particular, the curves \( \gamma^\alpha \) and \( \gamma_\delta \) not only have similar lengths (due to estimate \eqref{eq:length_estimate}), but are also uniformly close in the \( C^1 \) topology, as a consequence to Lemma \ref{lem:uniform_conv_geodesic_flows}.
    Indeed, if \( (x_\epsilon, v_\epsilon) = (\gamma^\alpha_\epsilon(0), \dot\gamma^\alpha_\epsilon(0)) \), then the orbit of \( H_{\epsilon, \delta}(x_\epsilon, v_\epsilon) \) under \( \mathbf{g}_\delta \) must be close as well, and by uniqueness it coincides with \( \gamma^\alpha_\delta\).
    Thus the curves \( \gamma^\alpha_\epsilon \) become arbitrarily close as \( \epsilon \to 0\), and the accumulation point of \(\{\gamma^\alpha_\epsilon\}\) is unique.
    \par The curve \( \gamma^\alpha_0 \) is a closed geodesic for \( d_0 \), being the \(C^1\)-limit of closed geodesics. 
    Given any other curve \( \gamma \in \alpha \), we have \( l_\epsilon(\gamma) > l_\epsilon(\gamma^\alpha_\epsilon) \), therefore \( l_0(\gamma) \geq l_0(\gamma^\alpha_0)\).
    If equality holds, then \( \gamma \) must be locally minimising, and therefore a geodesic. 
    Hence \( \gamma \) must be the limit of a sequence of geodesics \( \gamma_{\epsilon_n} \) of \( K_{\epsilon_n} \).
    We claim that for sufficiently large \( n \) the curves \( \gamma \) and \( \gamma_{\epsilon_n} \) are homotopic, and therefore \( \gamma_{\epsilon_n} = \gamma^\alpha_{\epsilon_n} \), implying \( \gamma = \gamma^\alpha_0 \).
    Indeed, let \( \delta < R_0\), the common injective radius of the family \( K_\epsilon \), and fix \( N \) such that 
    \( \gamma_{\epsilon_n} \) is contained in an \( \delta \)-neighbourhood \( B(\gamma; \delta) \) of \( \gamma \) for every \( n > N \) (here the neighbourhood is measured with respect to the reference metric \(g = g_1 \)).
    By our choice of \( \delta \), observe that that the set
    \[
        D_\delta := \{(x,y) \in K; d_1(x,y) < \delta\}
    \]
    is the diffeomorphic image of the set 
    \[
       T^\delta K := \{(x,v) \in T^1K; g_1(v, v) < \delta^2\}
    \]
    under the diffeomorphism \( (x,v) \mapsto (x, \exp_x^1(v))\).
    In particular, the zero section in \( TK \) is a deformation retract of \( T^\delta K \), and therefore \( D_\delta \) is homotopic to the diagonal set \( \{(x,x)\} \).
    Now, since 
    \[
        \gamma_{\epsilon_n} \subset B(\gamma; \delta) \subset D_\delta,
    \]
    it follows that \( \gamma \) and \( \gamma_{\epsilon_n}\) are freely homotopic.
    \par Hence \( \gamma^\alpha_0 \) is indeed the unique geodesic in \( \alpha \), which concludes.
\end{proof}
\begin{defi}[Enriched marked length spectrum]
    Let \( \mathcal{D} \) be a Sinai billiard and \( K \) a surface in \( \cK(\mathcal{D}) \).
    The function 
    \begin{align*}
        \mathcal{EL}_{\mathcal{D}; K}\colon \mathcal{C}(K) &\to \R \\
        \alpha &\mapsto l_0(\gamma^\alpha)
    \end{align*}
    is called the enriched marked length spectrum of \( \mathcal{D} \).    
\end{defi}
\par Suppose now \( \mathcal{D} \) and \( \mathcal{D}' \) are two distinct billiard tables in \( \mathfrak{D}^m \approx C^r_{\text{emb}}(\T_\mathcal{D}, \T^2) \).  
We consider the latter as a metric space with the \( C^r \) metric.  
Remark that the tables \( \mathcal{D} \) and \( \mathcal{D}' \) are diffeomorphic.  
In fact, when we consider the embeddings \(\gamma\colon\mathbb{T}_\mathcal{D} \to \mathbb{T}^2 \), \( \gamma'\colon\mathbb{T}_\mathcal{D} \to \mathbb{T}^2 \) defining their boundaries, we see that  
\[
    F\colon \mathbb{T}_\mathcal{D} \times I \to \mathbb{T}^2, \quad (s,t) \mapsto \gamma(s) + t(\gamma'(s)-\gamma(s))
\]
is an isotopy of \( \mathbb{T}_\mathcal{D} \) on \( \mathbb{T}^2 \), and by the isotopy extension theorem (cf. Chapter 8, Theorem 1.3 in \cite{hirsch2012differential}), it extends to a diffeotopy of \( \mathbb{T}^2 \) with compact support.  
\par We will consider surfaces \( K \subset \cK(\mathcal{D}), \ K' \subset \cK(\mathcal{D}') \), obtained by gluing two copies of \( \mathcal{D} \), respectively \( \mathcal{D}' \), that is  
\[
    K \setminus B(K) = \mathcal{D}^\mathrm{up} \coprod \mathcal{D}^\mathrm{down}, \qquad K' \setminus B(K') = (\mathcal{D}')^\mathrm{up} \coprod (\mathcal{D}')^\mathrm{down}.
\]  
The surfaces \( K \) and \( K' \) are also diffeomorphic. 
Indeed, \( p(\mathcal{D}^\mathrm{up}) \approx \mathcal{D} \approx \mathcal{D}' = p((\mathcal{D}')^\mathrm{up}) \), and similarly for \( \mathcal{D}^\mathrm{down} \).  
As \( \partial \mathcal{D} \approx \partial \mathcal{D}' \), we obtain a diffeomorphism \( h\colon K \to K' \), as wanted.  
\par In what follows, we will fix \( K \) as our reference surface.
As in Section \ref{sec:geodesic_approximations}, we will consider families of metrics \( g_\epsilon \), \( g'_\epsilon \) on \( K \) converging to the gluing metric when \( \epsilon \) approaches \( 0 \).  
More specifically, \( g = g_1 \) is our reference metric on \( K \), induced by the Euclidean metric of \( E = \mathbb{T}^2 \times \R \).  
For each of the tables, the metrics \( g_\epsilon \), \( g'_\epsilon \) are defined as the pullback to \( K \), under the mapping \(\alpha_\epsilon\colon (x,y,z) \mapsto (x,y,\epsilon z) \), of the metric induced on \(\alpha_\epsilon(K) \), \(\alpha_\epsilon(K') \), respectively.  
The \( g'_\epsilon \) are seen as metrics in \(K\) by pulling back with the diffeomorphism between the surfaces.
\par We have metric convergence \( (K, g_\epsilon) =: K_\epsilon \xrightarrow{GH} K_0 := (K, g_0) \), and similarly \( (K, g'_\epsilon) =: K'_\epsilon \xrightarrow{GH} K_0' := (K, g'_0) \), in the Gromov-Hausdorff distance, as \( \epsilon \to 0 \).  
The metric space \( K_0 \) (respectively \( K_0' \)) consists of two copies of \( (\mathcal{D}, d_\mathcal{L}) \) (respectively \( (\mathcal{D}', d_\mathcal{L}) \)) glued at the boundary.  
Each family \(\{g_\epsilon\}_\epsilon\), \(\{g'_\epsilon\}_\epsilon\) satisfies uniform estimates as in \eqref{eq:bound_geom}  
\[
    \|g_\epsilon - g_\delta\|_{C^r(K)} \leq  C|\epsilon - \delta|, \qquad 
    \|g'_\epsilon - g'_\delta\|_{C^r(K)} \leq  C'|\epsilon - \delta|.
\]
From now on we set \( C:= \max\{C, C'\} \), in order to have a single constant.  
Note that we cannot, a priori, say anything about how close two metrics \( g_\epsilon, g'_\delta \) coming from \emph{different tables} are.  
\par In what follows \( l_\epsilon \), \( l'_\epsilon \) will denote the length functions with respect to the metrics \( g_\epsilon \), \( g'_\epsilon \), and we write \( \mathcal{EL} = \mathcal{EL}_{\mathcal{D}} \), \( \mathcal{EL}' = \mathcal{EL}_{\mathcal{D}'} \) to denote the enriched length spectrums of the given tables.

\begin{theorem}\label{thm:conj_abst_flow}
    Let \( \mathcal{D} \) and \( \mathcal{D}' \) be distinct billiard tables with the same enriched length spectrum.
    Then the surfaces \( \wK_0 \), \( \wK_0' \) have conjugated abstract geodesic flows, and the conjugacy
    \[
        h\colon \widehat{\cG}_0 \to \widehat{\cG}_0'
    \]
    preserves the Liouville currents, i.e, \(h_*\widehat{m}_0 =\widehat{m}_0'\).
    Moreover, \(h_0\) is \(\pi_1(K)\)-equivariant and descends to a conjugacy of the abstract geodesic flows of the compact surfaces \(K_0\) and \(K'_0\).
\end{theorem}
\begin{proof}
    The condition \( \mathcal{EL} = \mathcal{EL}' \) implies that for every conjugacy class \( \alpha \in \mathcal{C}(K) \) we have
    \[
       \lim_{\epsilon \to 0} \left(l_\epsilon(\gamma^\alpha)-l'_\epsilon(\gamma^\alpha)\right) = 0.
    \]
    In particular, the functions \( l_\epsilon - l'_\epsilon \) are converging uniformly to zero on compact subsets.  
    Indeed, given \(T > 0\), for every \(\alpha \in \mathcal{C}(K)\) such that \(l_\epsilon(\gamma^\alpha) < T\) we have 
    \[
    \begin{split}
        l_\epsilon(\gamma^\alpha) - l_\delta(\gamma^\alpha)| &< CT|\epsilon - \delta|; \\
        |l'_{\epsilon'}(\gamma^\alpha) - l'_{\delta'}(\gamma^\alpha)| &< CT|\epsilon' - \delta'|, \\
    \end{split}
    \]
    and thus, given \(\eta >0\), we may find \( \epsilon_0 >0\) such that
    \[
        |l_\epsilon(\gamma^\alpha) - l_0(\gamma^\alpha)| + |l'_{\epsilon'}(\gamma^\alpha) - l'_0(\gamma^\alpha)| < \frac{\eta}{2},\quad \forall \epsilon, \epsilon' <  \epsilon_0.
    \]
    As, by hypothesis, \(l_0(\gamma^\alpha) = l'_0(\gamma^\alpha)\), we have
    \[
        |l_\epsilon(\gamma^\alpha) - l'_\epsilon(\gamma^\alpha)| < \eta, \ \ \forall \alpha \in \mathcal{C}(K); l_\epsilon(\gamma^\alpha) < T.
    \]
    Consequently, 
    \[
        \frac{l_\epsilon(\gamma^\alpha)}{l_\epsilon'(\gamma^\alpha)} = 1 + O\left(\frac{|\epsilon-\epsilon'|}{l_\epsilon(\gamma^\alpha)}\right) \to 1.
    \]

\par At each \(\epsilon>0\), the metrics \(g_\epsilon, g_\epsilon'\) define topologically  equivalent (or orbit equivalent) Anosov geodesic flows.
We consider a family of homeomorphisms \(h_\epsilon\colon \widehat{\cG}_\epsilon \to \widehat{\cG}_\epsilon'\) between the phase spaces constructed in the following way. 
For \(\epsilon > 0\) we have canonical identifications \(\widehat{\cG}_\epsilon \approx T^1\wK_\epsilon\). 
Start with a Hölder homeomorphism \(h_1\colon T^1\wK_1 \to T^1\wK'_1\) semi-conjugating the geodesic flows.
In particular, such map is uniformly continuous.
Then, we consider families of equicontinuous equivalences \(H_\epsilon\colon T^1\wK_1 \to T^1\wK_\epsilon \) and \(H'_\epsilon\colon T^1\wK'_1 \to T^1\wK'_\epsilon \), such as the ones constructed in Remark \ref{rem:anosov_top_equivalences}, for instance.
Let
\[
    h_\epsilon := H'_\epsilon \circ h_1 \circ H^{-1}_\epsilon.
\]
As \(h_1\) is uniformly continuous, we get an equicontinuous family of semi-conjugacies.
Note that such family is equivariant with respect to the canonical \(\pi_1(K)\) actions via isomorphisms.
The family of inverses \(\{h_\epsilon^{-1}\}\) has the same properties.
\par We define \(h_0\) as the ``limiting homeomorphism'', in the following sense.
Let \(\gamma\) be a lifting of a dense geodesic in \(K_0\). 
Choose a reference point \(x_0 := \gamma(0)\).
For \(\epsilon > 0\), we let \(\gamma_\epsilon\) be the unique geodesic in \(\widehat{\cG}_\epsilon\) with same endpoints at infinity as \(\gamma\).
Then \(\gamma_\epsilon \to \gamma\) on compact sets as \(\epsilon \to 0\).
In particular, \(x_\epsilon := \gamma_\epsilon(0) \xrightarrow{\epsilon \to 0} x_0\).
We write 
\[
    v_\epsilon := (x_\epsilon, \dot\gamma_\epsilon(0)) \in T^1\wK_\epsilon \quad \text{and} \quad \gamma' := \lim_{\epsilon \to 0}h_\epsilon(v_\epsilon) \in \widehat{\cG}_0.
\]
Note that the projection of each \(\gamma_\epsilon\) is dense on \(K_\epsilon\), and therefore the same is true for each \(h_\epsilon(v_\epsilon)\) and for the limit \(\gamma'\).
Finally, set \(h_0\) on the \(\pi_1(K)\)-orbit \(\cO(\gamma) := \{\alpha\cdot \gamma; \alpha \in \pi_1(K)\}\) by letting
\[
    h_0(\alpha\cdot\gamma)(t) := \alpha\cdot\gamma' (t). 
\]
\(\cO(\gamma)\) is a saturated (hence invariant) subset under the geodesic flow, and it is clear by definition that \(h_0\) is a flow conjugacy on \(\cO(\gamma)\):
\[
    h_0(\mathbf{g}_0^s(\alpha\cdot\gamma)) = h_0(\alpha\cdot\gamma(\cdot + s)) = \alpha\cdot\gamma'(\cdot+ s) = \mathbf{g}'^s_0(h_0(\alpha\cdot\gamma)).
\]
We claim that \(h_0\) is uniformly continuous, and therefore has a unique extension \(h\colon \widehat{\cG}_0 \to \widehat{\cG}_0'\), which is the desired conjugacy.
To that end, it is sufficient to show that \(h_0\) preserves the ``returning times'' of \(\gamma\).
\par We argue by approximation using the geodesic flows on \(T^1K_\epsilon\), which is better suited to the use of recurrence arguments due its compactness.
Note that all flows and homeomorphisms defined above are \(\pi_1(K)\)-equivariant and descend well to the quotient.
\par Since the geodesic flows are Anosov, for a given \(\delta>0\) we may find positive \(T = T(\delta)\) and \(\epsilon_0\) such that the set of \(g_{\epsilon_0}\)-periodic points of period less than \(T\) is \(\delta\)-dense on \(T^1K_{\epsilon_0}\), i.e., it intersects every ball of radius \(\delta\).
In terms of the compact-open topology on the space of geodesics, this translates to saying that for every geodesic \(\sigma\) and every \(t < T\) there is a periodic orbit \(\rho\) such that
\[
   \sup_{0 \leq s \leq t} d_\epsilon(\sigma(s), \rho(s)) < \delta.
\]
By the equicontinuity of the orbit equivalences $\{h_\epsilon\}_{\epsilon < \epsilon_0}$ between the geodesic flows $\{\mathbf{g}_\epsilon\}_{\epsilon < \epsilon_0}$, enlarging \(\delta\) a bit, and choosing a smaller \(\epsilon_0\) if necessary, we may assume that \(g_{\epsilon}\)-periodic points of period less than \(T\) are \(\delta\)-dense on \(T^1K_{\epsilon}\) for every \(\epsilon < \epsilon_0\).
Moreover, we have \(T \xrightarrow{\delta\to 0} \infty\).
\par Now, equicontinuity of \(h_\epsilon^{-1}\) means there is an \(\eta = \eta(\delta)\) such that the set of \(g'_\epsilon\)-periodic points is \(\eta\)-dense on \(T^1K'_{\epsilon}\) for every \(\epsilon < \epsilon_0\).
We have \(\eta \xrightarrow{\delta \to 0} 0\).
Suppose, for some \(t < T\), and \(\epsilon\) small, we have 
\begin{equation}\label{eq:returning_times}
    d_\epsilon(x_\epsilon, \mathbf{g}^t_\epsilon(x_\epsilon)) < \delta. 
\end{equation}
Then, there is a \(g_\epsilon\)-periodic point \(p_\epsilon\) of length \(l_\epsilon(p_\epsilon) \approx t < T\) such that \(d_\epsilon(x_\epsilon, p_\epsilon) < \delta\), and consequently, for \(y_\epsilon := h_\epsilon(x_\epsilon)\) and \(q_\epsilon := h_\epsilon(p_\epsilon)\),
\[
    d'_\epsilon(y_\epsilon, \mathbf{g}^t_\epsilon(y_\epsilon)) \leq d'_\epsilon(y_\epsilon, q_\epsilon) + d'_\epsilon(q_\epsilon, \mathbf{g}^t_\epsilon(y_\epsilon)) \leq 2\eta + |l_\epsilon(p_\epsilon) - l'_\epsilon(q_\epsilon)|.
\]
Now, \(|l_\epsilon(p_\epsilon) - l'_\epsilon(q_\epsilon)| \xrightarrow{\epsilon \to 0} 0\), hence we have 
\[
    d_\epsilon(x_\epsilon, \mathbf{g}^t_\epsilon(x_\epsilon)) < \delta \implies d'_\epsilon(y_\epsilon, \mathbf{g}^t_\epsilon(y_\epsilon)) \leq \eta' = \eta'(\delta),
\]
for all \(t < T\) and for some constant \(\eta'= 2\eta + O(\epsilon)\).
In terms of lifts to \(\wK_\epsilon\), we have a unique point \(z_0\) from fibre of \(\mathbf{g}^t_\epsilon(x_\epsilon)\) in a \(\delta\)-neighbourhood of the reference point \(x_0\), which means \(z_0 = \sigma(t)\) for a geodesic \(\sigma\) obtained as the lifting of \(\pi(\gamma_\epsilon)\) with a different starting point.
But this means exactly that \(\sigma = \alpha\cdot\gamma_\epsilon\) for some \(\alpha \in \pi_1(K)\), and we recover uniform continuity at height \(\epsilon\) in the original setting.
\par Finally, if \(d_0(x_0, \mathbf{g}_0^t(x_0)) < \delta\), then the estimate \eqref{eq:returning_times} is satisfied for every \(\epsilon\) sufficiently small.
We may then let \(\epsilon \to 0\) to get uniform continuity of \(h_0\), as wanted.
Since \(\pi(\gamma)\) is a dense geodesic on \(K_0\), it follows that for every \(\sigma \in \widehat{\cG}_0\) there are sequences \(t_n \in \R, \alpha_n \in \pi_1(K)\) such that \(\alpha_n\cdot\mathbf{g}_0^{t_n}\gamma \to \sigma\), and therefore \(\overline{\cO(\gamma)} = \widehat{\cG}_0\).
Similarly for the image \(\cO(\gamma') = h_0(\cO(\gamma))\).
Hence \(h_0\) extends uniquely to a flow conjugacy \(h_0\colon \widehat{\cG}_0 \to \widehat{\cG'}_0\), as wanted.

\par We are left to show that \(h\) preserves the currents.
Since the currents are given in terms of cross-ratios, as in Equation \eqref{eq:cr_rat_def}, it is sufficient to show equality between the cross ratios at \(\cG_0\) of corresponding quadruples of points at infinity.

\par Corollary~\ref{cor:dense_periodic_liftings} guarantees the set of lifts of periodic geodesics of \(K_\epsilon\) is dense in \(\cG_\epsilon\).
Thus, any rectangle \(|\zeta_1, \zeta_2| \times |\zeta_3, \zeta_4| \subset \bbK{\epsilon}\) may be arbitrarily well approximated by rectangles whose boundary is the union of segments of periodic orbits.
We can use this to get uniform bounds on the deviation between the cross ratios \(\crat{\epsilon}\) and \(\crat{\epsilon}'\), in the following way.
\par Recall that we may write \(\crat{\epsilon}(\zeta_1, \zeta_2, \zeta_3, \zeta_4) = \lim_{t \to \infty} e^{\frac{1}{2}E_\epsilon\overline{\sigma}(t)}\), where
\[
    E_\epsilon\overline{\sigma}(t) = d_\epsilon(\sigma_1(t), \sigma_1(-t)) + d_\epsilon(\sigma_2(t), \sigma_2(-t)) - d_\epsilon(\sigma_3(t), \sigma_3(-t)) - d_\epsilon(\sigma_4(t), \sigma_4(-t)),
\]
and the \(\sigma_i\) are the geodesics representing pairs at infinity as follows:
\[
    \begin{split}
    \sigma_1 &= (\xi_1, \xi_2); \\
    \sigma_2 &= (\xi_3, \xi_4); \\
    \sigma_3 &= (\xi_1, \xi_4); \\
    \sigma_4 &= (\xi_3, \xi_2).
    \end{split}
\]
Now, because of the density of periodic orbits and the continuity of \(h_\epsilon\) in the compact-open topology, given \(\eta > 0\) we may choose periodic orbits \(\gamma^n_i \in \cG_\epsilon\) such that
\[
    \min\left\{\sup_{t \in [-n, n]}\{d_\epsilon(\sigma_i(t), \gamma^n_i(t))\}, \sup_{t \in [-n, n]}\{d'_\epsilon(h_\epsilon\sigma_i(t), h_\epsilon\gamma^n_i(t))\}\right\} < \frac{\eta}{4}.
\]
For a quadruple \(\overline\sigma = (\sigma_1, \sigma_2, \sigma_3, \sigma_4)\), we write \(E'_\epsilon\overline{\sigma}\) to denote the corresponding cross-ratio of the geodesics \(h_\epsilon\sigma_i \in \cG_\epsilon'\), with respect to \(d_\epsilon'\).
Thus, we have, for any \(t \in [-n, n]\)
\begin{align*}
    |E_\epsilon\overline{\sigma}(t) - E'_\epsilon\overline{\sigma}(t)| &\leq  |E_\epsilon\overline{\sigma}(t) - E_\epsilon\overline{\gamma}(t)| + |E_\epsilon\overline{\gamma}(t) - E'_\epsilon(\overline{\gamma}(t)| + |E'_\epsilon\overline{\gamma}(t) - E'_\epsilon\overline{\sigma}(t)|\\
    &\leq \eta + 4\delta + \eta,
\end{align*}
where we are using that the spectra are \(\delta\)-close at height \(\epsilon\) in order to bound the middle term in the right hand side.
As \(\eta\) is arbitrary, we conclude that 
\[
    \sup_{t \in [-n,n]}|E_\epsilon\overline{\sigma}(t) - E'_\epsilon\overline{\sigma}(t)| \leq 4\delta.
\]
Letting \(n \to \infty\) and using that both \(E_\epsilon\overline{\sigma}\) and \(E_\epsilon'\overline{\sigma}\) are convergent (in particular, bounded) sequences gives us 
\[
    |\crat{\epsilon}(\zeta_1, \zeta_2, \zeta_3, \zeta_4) - 
    \crat{\epsilon}'(h_\epsilon\zeta_1, h_\epsilon\zeta_2, h_\epsilon\zeta_3, h_\epsilon\zeta_4)| \leq L(\epsilon)\delta, 
\]
where \(L(\epsilon)\) is a constant obtained by applying the mean value theorem to the difference \(e^{\frac{1}{2}E_\epsilon\overline{\sigma}(t)} - e^{\frac{1}{2}E'_\epsilon\overline{\sigma}(t)}\), and depends, a priori, on \(\epsilon\).
Now, on the other hand, \(E_\epsilon\overline{\sigma} \to E_0\overline{\sigma}\), and similarly for \(E'_\epsilon\overline{\sigma}\). 
Hence, the sequences are also bounded in \(\epsilon\), which means we can find a uniform constant \(L\) such that
\[
    |\crat{\epsilon}(\zeta_1, \zeta_2, \zeta_3, \zeta_4) - 
    \crat{\epsilon}'(h_\epsilon\zeta_1, h_\epsilon\zeta_2, h_\epsilon\zeta_3, h_\epsilon\zeta_4)| \leq L\delta, \quad \forall \epsilon \in (0, \epsilon_0).
\]
Finally, letting \(\epsilon \to 0\) we may choose \(\delta\) arbitrarily small, and therefore
\[
    \crat{0}(\zeta_1, \zeta_2, \zeta_3, \zeta_4) - 
    \crat{\epsilon}'(h\zeta_1, h\zeta_2, h\zeta_3, h\zeta_4),
\]
which concludes.
\end{proof}

\subsection{Average angle dissipation}
    We recall that, in a CAT(0) space \((X, d)\), the angle \(\angle_x(\gamma, \gamma') \in [0, \pi)\) between two geodesics \(\gamma, \gamma'\colon [0,1] \to X\) intersecting at the point \(x = \gamma(0) = \gamma'(0)\) is given by
    \[
        \angle_x(\gamma, \gamma') = \lim_{t\to 0} 2\arcsin{\left(\frac{1}{2t}d(\gamma(t), \gamma'(t))\right)}.
    \]
    In particular, for any geodesic triangle \(\triangle(x,y,z)\) the following standard comparison estimate holds:
    \begin{equation}\label{eq:GB_estimate_subs}   
        \angle_x(y,z) + \angle_y(z,x) + \angle_z(x,y) \leq \pi.
    \end{equation}
    When \((X,d)\) is a Riemannian manifold with the Riemannian metric, this coincides with the Riemannian angle (cf. \cite[Corollary II.1A.7]{bridson2011metric}).
    Specialising to the case our case, we get that angles in \(\wK_0\) and \(K_0\) are limits of corresponding angles in \(K_\epsilon\):
    \begin{equation*}
    \begin{split}
        \angle_x(\gamma, \gamma') &:= \lim_{t\to 0} 2\arcsin{\left(\frac{1}{2t}d_0(\gamma(t), \gamma'(t))\right)} \\
                                    & = \lim_{t\to 0} \lim_{\epsilon \to 0}2\arcsin{\left(\frac{1}{2t}d_\epsilon(\gamma(t), \gamma'(t))\right)} \\
                                    & = \lim_{\epsilon \to 0} \lim_{t\to 0} 2\arcsin{\left(\frac{1}{2t}d_\epsilon(\gamma(t), \gamma'(t))\right)} \\
                                    & = \lim_{\epsilon \to 0}\arccos{g_\epsilon(\dot\gamma(0), \dot\gamma'(0))},
    \end{split}
    \end{equation*}
    where the limit operators commute due to uniform convergence on compacts.
    \begin{remark}
        Note that angles are always well defined on \(\wK_0\), even though they do not specify the geodesic, at least not at degenerate points.
        In fact, two geodesics bifurcate at a point \(x\) if and only if their angle at \(x\) is zero, and we may characterise the boundary set \(B(\wK)\) exactly as the locus of points where geodesic bifurcation happens, that it, points from which emanates a pair of geodesics with angle zero.
    \end{remark}
    \par Equipped with such angles, we may define coordinate charts on \(\widehat{\cG}_0\) as follows.
    Given a reference geodesic \(\sigma \in \widehat{\cG}_0\), let \(\mathcal{U}_\sigma\) denote the set of all \(\gamma \in \widehat{\cG}_0\) intersecting \(\sigma\) transversally.
    Such intersection point is necessarily unique, and we set
    \begin{align*} 
        \varphi_\sigma\colon \mathcal{U}_\sigma &\to \R \times [0, \pi) \times \R\\
        \gamma &\mapsto (t(\gamma), \theta(\gamma), s(\gamma)),
    \end{align*}
    where
    \begin{itemize}
        \item[(i)] \(t = t(\gamma)\) and \(s = s(\gamma)\) are the unique real numbers such that \(\sigma(-t) = \gamma(-s)\);
        \item[(ii)] \(\theta = \theta(\gamma):= \angle_{\sigma(-s)}(\sigma, \gamma)\) is the angle between \(\gamma\) and \(\sigma\) at the intersection point.
    \end{itemize}
    Note that \(\mathcal{U}_\sigma\) is an open set on \(\widehat{\cG}_0\), since transversality is an open condition. 
    If \(\sigma\) happens to be a degenerate geodesic on \(\wK_0\), then none the bifurcating geodesics are in \(\mathcal{U}_\sigma\), since the intersection at bifurcation points is necessarily tangent.
    Hence, \(\varphi_\sigma\) is a well-defined  continuous mapping in the compact-open topology of \(\widehat{\cG}_0\), though it need not be a surjection onto \(\R \times [0, \pi)\).

    \par The coordinates above can be given to any of the geodesic spaces \(\widehat{\cG}_\epsilon\). 
    It is well-known (cf.~\cite{otal_spectre_1990, wilkinson_anosov, bankovic_marked-length-spectral_2017}, among others), that in the Riemannian case the Liouville form \(\lambda\) is expressed in such coordinates as
    \[
        \lambda_{\epsilon}(t,\theta,s) = \frac{1}{2}\sin\theta\ d\theta \ dt \ ds.
    \]
    In particular, given a geodesic segment \([x,y]_\epsilon\) in \(\wK_{\epsilon}\), let \(A_\epsilon(x,y) = \{\gamma \in \widehat{\cG}_\epsilon; \gamma(\R) \pitchfork [x,y]_\epsilon\}\).
    Then Crofton's formula states:
    \[
        \widehat{m}_\epsilon(A_\epsilon(x,y)) = \int_{A_\epsilon(x,y)} d\lambda_\epsilon = l([x,y]_\epsilon) = d_\epsilon(x,y).
    \]
    In the limit there is not, \textit{a priori}, a Liouville form on \(\wK_0\), but using the correspondence between Liouville forms and Liouville currents for the Anosov flows, we get that \(\widehat{m}_0(A_0(x,y)) = d_0(x,y)\).
    In terms of geodesic segments \(\sigma\colon [a,b] \to \wK_0\), this maybe be stated as the following Crofton formula:
    \begin{equation}\label{eq:crofton_height_0}
        \widehat{m}_0(\{\gamma \in \widehat{\cG}_0; \gamma(\R) \pitchfork \sigma([a,b])\}) = b-a = l_0(\sigma).
    \end{equation}      
    

    
    \par There is an action \(\mathbb{S}^1 \curvearrowright \widehat{\cG}_0\) given as follows.
    A geodesic \(\gamma\colon \R \to \wK_0\) is a differentiable curve, and \(v := \dot\gamma(0)\) is a unit vector. 
    For any \(\theta \in \mathbb{S}^1\), we denote by \(\theta\cdot v \in T^1\wK \) the positive rotation of \(v\) by \(\theta\), and let \(\theta\) act on \(\gamma\) by setting
    \[
        \theta\cdot\gamma := \lim_{\epsilon \to 0}\gamma_{(\epsilon, \theta\cdot v)} \ ,
    \]
    where \(\gamma_{(\epsilon, \theta\cdot v)}\) denotes the unique geodesic on \(\wK_\epsilon\) with initial condition \(\theta\cdot v\).
    It is clear from this construction that 
    \[
        \angle_{\gamma(0)}(\gamma, \theta\cdot\gamma) = \theta.
    \]
    Now suppose \(\mathcal{D}, \mathcal{D}'\) have the same enriched spectra and let \(h\colon\widehat{\cG}_0 \to \widehat{\cG}_0'\) be the \(\widehat{m}_0\)-preserving conjugacy of geodesic flows. The map
    \(h\) takes intersecting geodesics into intersecting geodesics, and the point of intersection is unique since \(\wK_0'\) has no conjugated points.
    Thus, we have a well-defined angle function \(\vartheta\colon (\gamma,\theta) \mapsto \angle(h\gamma, h(\theta\cdot\gamma))\), where the angle is measured at the unique intersection point.  
    This is easily checked to be continuous in both coordinates.
    Following Otal, we define a mapping \(\Theta\colon [0, \pi] \to [0, \pi]\) which measures the average angle distortion imposed by the conjugacy \(h\), by setting
    \[
        \Theta(\theta) := \frac{1}{\widehat{m}_0(\widehat{\cG}_0)}\int_{\widehat{\cG}_0}\vartheta(\gamma, \theta) d\widehat{m}_0(\gamma) = \int_{\widehat{\cG}_0}\vartheta(\gamma, \theta) d\mu(\gamma)
    \]
    where \(\mu\) is the probability measure obtained by normalising \(\widehat{m}_0\)\footnote{
        We could define \(\Theta\) on the entire circle \(\mathbb{S}^1\), but the flip invariance of the map \(h\) means \(\Theta(-\theta) = \Theta(\theta)\), so we lose nothing by considering the restriction to the half circle.
    }.

    We have the following result.
    \begin{lemma}\label{lem:results_wilkinson_adapted}
        The function \(\Theta\) satisfies
        \begin{itemize}
            \item[(i)] \textbf{Symmetry}: \(\Theta(\pi - \theta) = \pi - \Theta(\theta)\) for all \(\theta \in [0,\pi]\), and consequently \(\Theta(\pi/2) = \pi/2\).
            \item[(ii)] \textbf{Superadditivity}: For all \(\theta_1, \theta_2 \in [0, \pi]\) with \(\theta_1 + \theta_2 \leq \pi\):
            \[
                \Theta(\theta_1 + \theta_2) \geq \Theta(\theta_1) + \Theta(\theta_2).
            \]
            \item[(iii)] For all \(a \in [0, \pi]\):
            \begin{equation}\label{equation_grand_theta}
                \int_0^a \frac{\sin \theta}{\sin \Theta(\theta)}  d\theta \leq \Theta(a).
            \end{equation}
        \end{itemize}
    \end{lemma}
    \begin{proof}
        Items (i) and (ii) can be deduced exactly as in \cite[Proposition 4.3]{wilkinson_anosov}.
        Remark that, while no Gauss-Bonnet theorem can be used in \(\wK_0'\), the crucial part of the argument is the actual estimate on the sum of the angles, for which Inequality \eqref{eq:GB_estimate_subs} can be used to the same effect.
        \par As for item (iii), we can follow the steps outlined in~\cite[Proposition 4.5]{wilkinson_anosov}, which are as follows:
        \begin{enumerate}[label=\alph*.]
            \item\label{pointt_un} apply Jensen’s Inequality when averaging against the Liouville current, leveraging the strict convexity of the function $\csc=\frac{1}{\sin}$;
            \item\label{pointt_deux} use the ergodicity of the geodesic flows with respect to the Liouville measure to replace spatial averages with temporal averages, as guaranteed by Birkhoff’s Ergodic Theorem;
            \item\label{pointt_trois} rely on the preservation of the Liouville current under a time-preserving conjugacy between Anosov geodesic flows on surfaces. 
        \end{enumerate} 
        While step~\ref{pointt_un} is quite general, steps~\ref{pointt_deux}-\ref{pointt_trois} can be adapted to our context. Specifically, we could use the ergodicity of billiard flows in Sinai billiards with horizon to establish the ergodicity of $\mathbf{g}_0$ and $\mathbf{g}_0'$ for the respective measures $\mu$ and $\mu'$ --viewing $K$, resp. $K'$ as two copies of $\mathcal{D}$, resp. $\mathcal{D}'$ glued along their boundaries-- and then apply Theorem~\ref{thm:conj_abst_flow} which shows that the conjugacy $h$ preserves the Liouville currents. 

        Alternatively, we can derive approximate results for the approximating geodesic flows \(\mathbf{g}_\epsilon\) and \(\mathbf{g}_\epsilon'\) and take the limit as $\epsilon\to 0$. Indeed, the ergodicity of these flows with respect to their Liouville measures $\widehat{m}_\epsilon$, $\widehat{m}_\epsilon'$ follows from the Anosov property. Although the orbit equivalence $h_\epsilon$ between \(\mathbf{g}_\epsilon\) and \(\mathbf{g}_\epsilon'\) does not exactly preserve the Liouville currents,  Theorem~\ref{thm:conj_abst_flow} ensures that $\lim_{\epsilon\to 0}(h_\epsilon)_* \widehat{m}_\epsilon-\widehat{m}_\epsilon'=0$. Thus, the change of variables in the proof of~\cite[Proposition 4.5]{wilkinson_anosov} introduces an error term that vanishes as $\epsilon\to 0$, which allows us to recover~\eqref{equation_grand_theta} in the limit. 
    \end{proof}

Using Lemma \ref{lem:results_wilkinson_adapted} above, we may conclude that \(h\) preserves the angles between geodesics. 
If \(\wK_0\) were negatively curved, we could use the fact that the sum of the angles on a geodesic triangle is always strictly smaller than \(\pi\), except if the triangle degenerates into a single point. 
This is exactly how Otal originally argued in \cite{otal_spectre_1990}.
In our case, however, the space is only CAT(0), so in principle there's nothing to prevent the existence of triangles with sum of intern angles exactly equal to \(\pi\).
In fact, if \(h\) is ``close to the identity'', a triple of geodesics intersecting at an interior point \(x \in \wK^\circ\) would be mapped into three geodesics of \(\wK_0'\) whose intersection determines a geodesic triangle \(\triangle\) also inside \((\wK')^\circ\).
Since the negative curvature of the limit space \(\wK_0\) is concentrated at the boundary set \(B(\wK)\), it would follow that \(\triangle\) is a flat triangle, and the sum of its intern angles would be equal to \(\pi\).
\par This indicates that, in order to use angle preservation to derive some sort of isometry, we must look at triangles on a neighbourhood of the boundary set.
This is how we proceed to conclude the proof.

\begin{lemma}
    The conjugacy \(h\colon\widehat{\cG}_0 \to \widehat{\cG}_0'\) is angle preserving. 
    In other words, $\vartheta(\gamma, \theta) = \theta$, for all $\gamma \in \widehat{\cG}_0$.
\end{lemma}
\begin{proof}
    Let
    \[
        a := \sup\{t \in (0, \pi]; \Theta(\theta) < \theta, \quad \forall \theta \in (0, t)\}.
    \]
    Note that \(a \leq \pi/2\) since the later is a fixed point (Lemma \ref{lem:results_wilkinson_adapted} (i)), and also that \(\Theta(a) = a\) necessarily, due to continuity.
    Lemma \ref{lem:results_wilkinson_adapted} (iii) implies
    \[
        \int_0^a\sin{\theta}\csc{(\Theta(\theta))}d\theta \leq \Theta(a) = a.
    \]
    Now, for \(\theta \in [0, a] \subset [0, \pi/2]\) the condition \(\Theta(\theta) < \theta\) means \(\sin{\theta}\csc{(\Theta(\theta))} > 1\), so that we have
    \[
        a < \int_0^a\sin{\theta}\csc{(\Theta(\theta))}d\theta \leq  a,
    \]
    which is a contradiction as soon as \(a \neq 0\).
    Thus \(\Theta(\theta) \leq \theta\) in its entire domain (cf. \cite[Proposition 4.4]{wilkinson_anosov}), and symmetry of \(\Theta\) implies equality. 
    Since \(\Theta(\theta) = \theta = \mathbb{E}_\mu[\vartheta(\cdot, \theta)]\) and \(\vartheta(\gamma, \theta) \geq \vartheta(\gamma, \theta_1) + \vartheta(\gamma, \theta_2)\) for \(\theta = \theta_1 + \theta_2\), equality holds pointwise almost everywhere. 
    Continuity extends this to all \(\gamma\), and we get the desired angle preservation.
\end{proof}
\begin{theorem}[Enriched spectrum rigidity]
    Suppose \( \mathcal{EL} = \mathcal{EL}' \).
    Then there is an isometry \( \phi\colon \T^2 \to \T^2 \) such that \( \mathcal{D}' = \phi \circ \mathcal{D} \).
    Consequently, the induced map
    \begin{align*}
        \Phi\colon \Omega &\to \Omega' \\
        (x, \omega) &\mapsto (\phi(x), \omega)
    \end{align*}
    is a flow conjugacy between billiard flows.
\end{theorem}

\begin{proof}
Let \(h\colon\widehat{\cG}_0 \to \widehat{\cG}_0'\) be as in Theorem \ref{thm:conj_abst_flow}.
It suffices to show that it is induced by an isometry of the torus.
\par By the angle preservation, we conclude immediately that \(h\) must preserve the set of admissible geodesics. 
In fact, if \(\gamma \notin \widehat{\cG}(\mathcal{D})\), then there is a degenerate point \(x \in \gamma(\R)\). 
Let \(\sigma\) be a geodesic though \(x\) such that \(\angle_x(\sigma, \gamma) = 0\). 
Then the angle between \(h\sigma\) and \(h\gamma\) at their respective degenerate point of intersection is also \(0\), by angle preservation, and therefore \(h\gamma \notin \widehat{\cG}(\mathcal{D}')\).
\par Now, let \(x \in B(\wK)\), and let \(\Gamma\) be the connected component of \(B(K) \subset K_0\) containing the projection of \(x\) under the canonical projection. Note that 
\(\Gamma\) is the gluing of two copies of a scatterer from \(\mathcal{D}\) and is itself a closed geodesic.
We consider its lift \(\gamma_0\) to \(\wK_0\) such that \(\gamma_0(0) = x\). 
Let \(\gamma_1\) and \(\gamma_2\) be two distinct geodesics emanating from \(x\), which we may choose to be regular and such that \(\gamma_1\) is between \(\gamma_0\) and \(\gamma_2\).
Then by construction
\[
    \angle_x(\gamma_0, \gamma_1) + \angle_x(\gamma_1, \gamma_2) + \angle_x(\gamma_2, -\gamma_0) = \pi.
\]
By construction, $h$ is a limit of orbit equivalences between the geodesic approximations $\mathbf{g}_\epsilon$ and $\mathbf{g}_\epsilon'$, hence it maps the geodesic $\gamma_0$ associated to the ``boundary component'' $\Gamma$ to a geodesic $h\gamma_0$ associated to the corresponding ``boundary component'' $\Gamma'$ of $\mathcal{D}'$.   
Now, suppose that the geodesics \(h\gamma_i\) intersect forming a non-degenerate triangle, and let \(p \in \wK_0'\) be the intersection point of \(h\gamma_0\) and \(h\gamma_1\).
Since the canonical projections \(\wK_0' \to K_0' \to \mathcal{D}'\) are local isometries, and the angle is a local construction, we see that the angle \(\angle_p(h\gamma_0, h\gamma_1)\) is the same as the angle between the billiard trajectory corresponding to \(h\gamma_1\) and the boundary component of \(\partial\mathcal{D}'\) containing the projection of \(p\).
But the billiard trajectory is a line segment while the scatterer is a convex curve, and it is clear that the angle they form is strictly less than the angle in the corresponding comparison triangle in \(\R^2\).
The same holds for the angle \(\angle(h\gamma_0, h\gamma_2)\), implying that
\[
    \angle(h\gamma_0, h\gamma_1) + \angle(h\gamma_1, h\gamma_2) + \angle(h\gamma_2, -h\gamma_0) < \pi,
\]
a contradiction.
Therefore \(h\) induces a bijection 
\[
    \widetilde h_0 \colon B(\wK) \to B(\wK').
\]
    \begin{figure}[!ht]
        \centering
        \includegraphics[width=1.1\textwidth]{conjug.pdf}
        \caption{Three admissible geodesics 
intersecting at a boundary point are mapped into a triangle of admissible geodesics,
whose sum of the intern angles is smaller than $\pi$, contradicting angle preservation.}
        \label{fig:geodesic_triangles_table}
    \end{figure}
We claim that it is, in fact, an isometry.
Indeed, the connected components of \(B(\wK)\) are geodesics so if \(x, y\) belong to the same connected component, by continuity so do \(\widetilde h_0(x), \widetilde h_0(y)\), and since \(h\) is a flow conjugacy, it follows that \(d_0(x, y) = d'_0(\widetilde h_0(x), \widetilde h_0(y))\).
\par On the other hand, if \(x, y \in B(\wK)\) are in distinct components, let \(\gamma\) be a geodesic such that
\[
    \gamma(0) = x \quad \text{ and } \quad \gamma(1) = y.
\]
Now, \(\gamma([0,1])\) projects to a finite segment of billiard trajectory on \(\mathcal{D}\), and as such it has only finitely many collision points.
This means there are at most finitely many times 
\[
    0 = t_0 < t_1 < \cdots < t_k = 1
\]
such that \(x_i = \gamma(t_i) \in B(\wK)\). 
Since \(h\) is a flow conjugacy preserving the boundary points, we conclude \(h\gamma(t_1) \in B(\wK')\), and moreover, that \(t_1\) is the first collision time of \(h\gamma\) after \(t_0\), as there are no other times \(t \in [t_0, t_1]\) such that \(\gamma(t) \in B(\wK).\)
Iterating this argument, we conclude that \(h\gamma\) has exactly \(k\) collision times, the same ones as \(\gamma\).
In fact, due to our construction we must have
\[
    h\gamma(t_i) = \tilde{h}_0(x_i), \ \ \text{for } i \in \{0, \cdots, k\},
\]
and 
\[
    d'_0(\widetilde h_0(x), \widetilde h_0(y)) = \sum_{i = 1}^kl_0(h\gamma\rvert_{[t_{i-1}, t_i]}) = \sum_{i = 1}^kl_0(\gamma\rvert_{[t_{i-1}, t_i]}) = d_0(x, y).
\]
\par All that is left to do now is to extend \(\widetilde{h}_0\) to an isometry \(\widetilde{h}\colon (\wK, d_0) \to (\wK', d'_0)\). 
In order to do that, consider the periodic geodesics of \((\wK, d_0)\), i.e., geodesics projecting to periodic orbits of the flow. 
Given such a periodic geodesic \(\gamma\), take \(x_0 = \gamma(t_0) \in \gamma(\R) \cap B(\wK)\), which exists because the table \(\mathcal{D}\) has finite horizon.
What is more, there is \(t_1 \in [\tau_{\mathrm{min}}, \tau_{\mathrm{max}}]\) such that \(x_1:= \gamma(t_1) \in B(\wK)\), and is minimal with respect to this property. 
In other words, \(t_1\) is the first collision time of the geodesic \(\gamma\) which is strictly larger than \(t_0\).
Just as we argued above, the first collision time of \(h\gamma\) after \(t_0\) must be \(t_1\) as well.
\par We extend \(\widetilde{h}_0\) to the segment \(\gamma([t_0, t_1]\) by setting 
\[
    \widetilde{h}_0(\gamma(t)) := h\gamma(t) \ \ \text{for all } t \in [t_0, t_1],
\]
and similarly for the subsequent (countably many) collision times \(t_2 < t_3 < \cdots \).
This extension is clearly isometric along the periodic rays, and an argument completely analogue to the one used to show \(\widetilde{h}_0\) preserves the distance between points in different components of \(B(\wK)\) shows that distance between points in distinct periodic geodesics is also preserved.
\par We have then an isometry
\[
    \widetilde{h}_0\colon B(\wK) \cup \mathrm{Per} \to B(\wK') \cup \mathrm{Per'}.
\]
Now, \(B(\wK) \cup \mathrm{Per}\) is left invariant under the action of \(\pi_1(K)\), and 
\(\mathrm{Per}\) is dense in the universal covering \(\wK\).
As \(\wK\) is complete, we conclude that there is a unique \(\pi_1(K)\)-invariant isometric extension
\[
    \widetilde h \colon \overline{B(\wK) \cup \mathrm{Per}} = \wK \to \wK'= \overline{B(\wK') \cup \mathrm{Per}'},
\]
which then descends to an isometry \(h\colon K_0 \to K_0'\). 
By Corollary \ref{cor:isometry-preserves-boundary} the tables \(\mathcal{D}\) and \(\mathcal{D}'\) are isometric, which concludes.\qedhere

\end{proof}





%%% Appendices
\appendix

\section{Generic properties of periodic orbits of Sinai billiards}\label{sec_periodicorbitsbilliard}

\subsection{Length functionals and periodic orbits}
We use the notation introduced in Section~\ref{sec_symbolics_prelim}. We let \( \gamma_{(0,0)}:=\pi^{-1}\circ \gamma \), and for each integer vector \( (i,j)\in \Z^2 \), we let \( \gamma_{(i,j)}:=\gamma_{(0,0)}+(i,j) \). 
Denote by
\begin{align*}
    \tau_{(i,j)} \colon \T_\mathcal{D}\times \T_\mathcal{D} &\to \R_+ \\ 
    (s,s')& \mapsto \tau_{(i,j)}(s,s'):=\|\gamma_{(0,0)}(s)-\gamma_{(i,j)}(s')\|
\end{align*}
the Euclidean distance between the points of \( \partial \widetilde{\mathcal{D}} \) with respective parameters \( s,s' \), in the respective cells \( (0,0) \) and \( (i,j) \). 
 Recall that for any \( (s,s')\in \mathbb{T}_\mathcal{D} \times \mathbb{T}_\mathcal{D} \), we have
\begin{equation}\label{eq:first_gen_fct}
    \partial_1 \tau_{(i,j)}(s,s')=-\sin \varphi_{(i,j)},\quad \partial_2 \tau_{(i,j)}(s,s')=\sin \varphi_{(i,j)}',
\end{equation}
where \( \varphi_{(i,j)} \) and \( \varphi_{(i,j)}' \) are the respective angles that the line segment between \( \gamma_{(0,0)}(s) \) and \( \gamma_{(i,j)}(s') \) makes with respect to the normal vector at these points. 
Moreover,
\begin{equation}\label{eq:der_seconde_gen} 
\begin{array}{ll}
    &\partial_{11} \tau_{(i,j)}(s,s')=\mathcal{K}(s)\cos \varphi_{(i,j)}+\frac{\cos^2 \varphi_{(i,j)}}{\tau_{(i,j)}(s,s')},\\
    &\partial_{12} \tau_{(i,j)}(s,s')=\frac{\cos \varphi_{(i,j)}\cos\varphi_{(i,j)}'}{\tau_{(i,j)}(s,s')},\\
    &\partial_{22} \tau_{(i,j)}(s,s')=\mathcal{K}(s')\cos \varphi_{(i,j)}'+\frac{\cos^2 \varphi_{(i,j)}'}{\tau_{(i,j)}(s,s')},
\end{array} 
\end{equation}
where $\mathcal{K}(s)$, $\mathcal{K}(s')$ are the respective curvatures at the points with parameters $s$ and $s'$. 
Let \( \sigma=(\bar s_{0},\cdots,\bar s_{q-1}) \) denote the collision points of a periodic orbit of (least) period \( q \geq 2 \), with \( \bar s_k\in \T_{\rho_k} \), \( \rho_k\in \{1,\cdots,m\} \), for \( k\in \{0,\cdots,q-1\} \). We let \( \rho(\sigma):=(\rho_k)_{k=0,\cdots,q-1} \). 
Let us consider the lift to \( \widetilde{\mathcal{D}} \) of this orbit starting in the cell \( (0,0) \), and denote by  \( \{(\tilde i_k,\tilde j_k)\}_{k=0,\cdots,q} \) the labels of the cells where the first \( (q+1) \) bounces happen, with \( (\tilde i_0,\tilde j_0)=(0,0) \).
For \( k\in \{1,\cdots,q\} \), set \( I_k:=(\tilde i_k,\tilde j_k)-(\tilde i_{k-1},\tilde j_{k-1}) \). We denote \( I(\sigma):=(I_k)_{k=1,\cdots,q} \).
As we assume that the table has finite horizon, there is a constant \( M > 0 \) such that \( I_k\in [-M,M]^{2} \), for any \( k\in \{1,\cdots,q\} \). 

The following result is essentially~\cite[Lemma A.2]{de2023marked}; since the context is slightly different, we reproduce its proof for the reader's convenience.
\begin{lemma}\label{lem:quadratic-form}
    Let \( \sigma=(\bar s_{0},\cdots,\bar s_{q-1}) \) denote the collision points of a periodic orbit of (least) period \( q \geq 2 \).
    With the above notations, let \( \rho:=\rho(\sigma)=(\rho_k)_{k=0,\cdots,q-1} \), and  \( I=I(\sigma):=(I_k)_{k=1,\cdots,q} \).
    We define the length functional
    \begin{equation*}
        L_\rho^I\colon
        \left\{
        \begin{array}{rcl}
            \T_{\rho_0}\times \cdots \times \T_{\rho_{q-1}} &\to& \R\\
            (s_{0},\cdots,s_{q-1}) &\mapsto& \sum_{k = 1}^{q-1}\tau_{I_k}(s_{k-1},s_{k})+\tau_{I_q}(s_{q-1},s_0)
        \end{array}
        \right..
    \end{equation*} 
    Then \( \sigma \) is a critical point of \( L_\rho^I \), and it holds:
    \begin{align*}
        L_{\rho}^I(s_{0},\cdots,s_{q}) - 
        L_{\rho}^I(\bar s_{0},\cdots,\bar s_{q}) &= 
        Q_{\rho}^I(s_{0}-\bar s_{0},\cdots, s_{q}-\bar s_{q}) + R_{\rho}^I,
    \end{align*}
    where \( Q_{\rho}^I \) is a positive definite quadratic form and \( R_{\rho}^I \) is a remainder term that satisfies the estimate:
    \begin{align*}
        R_{\rho}^I = O(\|(s_{0}-\bar s_{0},\cdots, s_{q}-\bar s_{q})\|^{3}).
    \end{align*}
\end{lemma}
\begin{proof} 
    Let us abbreviate \( L=L_{\rho}^I \). 
    The fact that \( \sigma=(\bar s_0,\cdots,\bar s_{p-1}) \) is a critical point for \( L \) among configurations \( (s_0,\cdots,s_{p-1}) \) follows from \eqref{eq:first_gen_fct}, as \( \sigma \) is a closed billiard orbit. 
    Hence the lemma follows from Taylor's formula,	provided that we show that the Hessian of \( L \) at \( (\bar s_{0},\cdots,\bar s_{q}) \) is positive definite.	
    It is immediate from the definition of \( L \) that \( \partial_{ij}L = 0 \) if \( |i-j| > 1 \); in other terms, the Hessian of \( L \) is a tridiagonal (symmetric) matrix
    \begin{equation*}
    D^2L = 
    \begin{bmatrix}
        \partial_{00}L & \partial_{01}L & 0 & \cdots & 0 \\
        \partial_{10}L & \partial_{11}L & \partial_{12}L & \cdots & 0 \\
        0 & \partial_{21}L & \partial_{22}L & \cdots & 0 \\
        \vdots & \vdots & \vdots & \ddots & \vdots \\
        0 & 0 & 0 & \cdots & \partial_{q-1,q-1}L 
    \end{bmatrix}.
    \end{equation*} 
    For notational convenience, let \( \tau_{k} = h(\bar{s}_{k},\bar s_{k+1}) \); let \( \bar\varphi_{k}\in[-\frac \pi 2,\frac \pi 2] \) denote the angle formed by the outgoing trajectory at the \( k \)-th collision point and the unit normal vector	to the domain at \( \bar s_{k} \); finally, let \( \mathcal{K}_{k} \) denote the reciprocal of the radius of curvature at the point \( \bar s_{k} \).
    Recall that, by convention, \( \mathcal{K}_{k} > 0 \) for any \( k \). 	
    By \eqref{eq:der_seconde_gen}, the diagonal terms are given by:
    \begin{align*}
        \partial_{00} L &= \frac1{\tau_{0}}\cos^{2}\bar\varphi_{0}+\mathcal{K}_{0}\cos\bar\varphi_{0}\\
        \partial_{jj} L &= \left[\frac1{\tau_{j-1}}+\frac1{\tau_{j}}\right]\cos^{2}\bar\varphi_{j}+2\mathcal{K}_{j}\cos\bar\varphi_{j}\text{ for } 0 < j < q\\
        \partial_{qq} L &= \frac1{\tau_{q-1}}\cos^{2}\bar\varphi_{q}+\mathcal{K}_{q}\cos\bar\varphi_{q}, \intertext{while the off-diagonal terms are given by:}
        \partial_{jj+1}L &= \frac1{\tau_{j}}\cos\bar\varphi_{j}\cos\bar\varphi_{j+1}.
    \end{align*}
    Using the above expressions it is simple to prove the following lemma.
    \begin{lemma}\label{lem:positive-definite}
        For \( 0 \leq k \leq q \), let \( f_{k} \) denote the determinant of the \( (k+1)\times (k+1) \) top-left minor of \( \partial_{ij}L \).
        Then \( f_{k} > 0 \) for all \( 0 \leq k\leq q \).
    \end{lemma}
    \begin{proof}[Proof of Lemma~\ref{lem:positive-definite}]
   See the proof of~\cite[Lemma A.3]{de2023marked}. 
\end{proof}
    By Sylvester's criterion, the above lemma implies that all eigenvalues of \( \partial_{ij}L \) are positive, which completes the proof of our result.
\end{proof}


\begin{defi}
    The spectrum of the table \( \mathcal{D} \) is the metric space \( (\Spec{\mathcal{D}}, d_H) \), where \( \Spec{\mathcal{D}} \) is the set of all periodic orbits of the billiard flow of \( \mathcal{D} \), and \( d_H \) is the Hausdorff metric on the phase space \( \Omega \).    
\end{defi} 
It follows from Corollary~\ref{cor:fin_card_peri} that \( \Spec{\mathcal{D}} \) is a discrete totally bounded space.

\par 
    Remark that a critical point of some length functional \( L \) as in Lemma~\ref{lem:quadratic-form} is not necessarily a periodic orbit, as it may consist of a non-physical trajectory, that is, one that passes through one the obstacle. 

Let us introduce the following sets of orbits:
\begin{defi}
    Given a word $\omega = (\omega_0, \cdots, \omega_n) \in \mathcal{A}^{n+1}$, the cylinder \( C(\omega; n) \) is the set of points \( x \in \mathcal{M} \) whose lifted orbits have the first \( n+1 \) bounces in the scatterers \( B_{\omega_1}, \cdots, B_{\omega_n} \), in this order.
\end{defi}

\begin{defi}
    A \emph{generalised periodic orbit} of \( \mathcal{D} \) is the projection onto \( \mathcal{M} \) of a critical point of one the functionals \( L \). 
    A generalised orbit that is not a periodic orbit of the billiard \( \mathcal{D} \) is called a \emph{ghost} or \emph{non-physical}. 
\end{defi}

\begin{theorem}[Theorem 2.2 in \cite{bunimovich_variational_1995}]\label{thm:bunimovich_periodic} 
    Let \(  \omega = (\omega_0, \cdots, \omega_n)\in \mathcal{A}^{n+1} \) be a word in \( \Sigma_{\mathcal{M}} \) such that 
    \begin{align*}
        \pi(B_{\omega_i}) &\neq \pi(B_{\omega_{i+1}}), \text{ for } i = 0, \cdots, n-1; \\
        \pi(B_{\omega_n}) &= \pi(B_{\omega_0}).
    \end{align*}
    Then there is at most one generalised periodic orbit in the cylinder \( C(\omega; n) \).
\end{theorem}

\par On the other hand, one periodic orbit in \( \Spec{\mathcal{D}} \) is associated to several different words \( \omega \), corresponding to distinct liftings.
Even when we restrict ourselves to liftings starting at the \( (0,0) \) cell, distinct starting points in the same orbit may lead to distinct words. 
We will therefore need a choice of point in each of the periodic orbits of \( \Spec{\mathcal{D}} \).
\par Given \( O \in \Spec{\mathcal{D}} \), we can see it as a union of line segments in the square \( [0,1)^2. \) 
The finite set \( (\cup_i B_i) \cap O \) correspond to all the collision points of the orbit \( O \).
It is an ordered set, with the lexicographical ordering inherited from \( [0,1)^2 \times [-\pi/2, \pi/2] \). 
Set
\[
    \zeta(O) \coloneqq \min \{O \cap (\cup_i B_i)\} \in \mathcal{M}.
\]
\begin{defi}
    The \emph{type} of the periodic orbit \( O \in \Spec{\mathcal{D}} \) is the unique sequence \( [O] \coloneqq \iota(\zeta(O)) \in \Sigma_M \).
\end{defi}
\subsection{Good tables}\label{sec:good_tables}
Recall that we denote by \( \mathfrak{D}^m \) the space  of all Sinai billiard tables in \( \T^2 \) with \( m \) scatterers and $C^r$ boundary, $r \geq 3$.  
\par We want to show that for a generic choice of table in \( \mathfrak{D}^m \), the set of periodic orbits has simple length spectrum and contains no singular orbits. 
\newline
\subsubsection{Table perturbations and periodic orbits}
We begin by analysing how perturbations on the scatterers affect periodic orbits.

\begin{defi}
    Given \( p = \gamma(s_0) \in \partial\mathcal{D} \), we say a table \( \mathcal{D}' \in \mathfrak{D}^m \) is a localised perturbation of \( \mathcal{D} \) at \( p \) if there is a neighbourhood \( U \) of \( s_0 \in \T_\mathcal{D} \) such that the defining maps of \( \mathcal{D} \) and \( \mathcal{D}' \) coincide outside \( U \). 
\end{defi} 
Intuitively, a small perturbation around a single collision point of a periodic orbit should generate a new periodic orbit close to the original one. 
In order to see that this is in fact what happens, we investigate how a localised perturbation at a point \( p=\gamma(s_0) \) alters the length functional of the table.
\par We begin by describing the perturbation in terms of a $C^r$ function \( \lambda \) defined on \( \mathbb{T}_{\mathcal{D}} \).
For \( \epsilon >0 \), consider the \( \epsilon \)-perturbation
\begin{equation}\label{eq:defor_descrip}
    \gamma^\lambda(s) = \gamma(s) + \epsilon\lambda(s)n_s,
\end{equation}
where \( n_s \) is the normal vector to \( \gamma(s) \) at \( s \).
For sufficiently small \( \epsilon \), the perturbation \( \gamma^\lambda \) remains strictly convex and the resulting billiard table is still a Sinai billiard.
The perturbation \( \gamma^\lambda \) is a localised perturbation at \( p = \gamma(s_0) \) exactly when \( \lambda \) is supported in a small neighbourhood around \( s_0 \).
\par As before, we look at liftings of such curves to the Euclidean plane and consider mappings
\begin{equation*}
    \tau^\lambda_{(i,j)}(s,s') := \|\gamma^\lambda_{(0,0)}(s)-\gamma^\lambda_{(i,j)}(s')\|,
\end{equation*}
whose first order approximation is 
\begin{equation*}
    \tau^\lambda_{(i,j)}(s,s') = \tau_{(i,j)}(s,s') + \epsilon(\lambda(s)\cos\varphi_{(i,j)}(s,s') +\lambda(s')\cos\varphi'_{(i,j)}(s,s')) + O(\epsilon^{2}).
\end{equation*}

Now, we suppose \( \sigma = (s_0, \cdots, s_{q-1}) \) corresponds to a sequence of bouncing points of the table \( \mathcal{D} \).
We want to study the relation between critical points for the length functionals of the original and perturbed tables.
For notational convenience, we drop the indices \( (i,j) \).
Then the length functional \( L^\lambda \) can be written in terms of the original length functional \( L \) as
\begin{align*}
    L^\lambda(\sigma) &= \sum_{i=0}^{q-1}\tau^\lambda(s_i, s_{i+1}) \\
    &= \sum_{i=0}^{q-1}h(s_i, s_{i+1}) + \epsilon\sum_{i=0}^{q-1}(\lambda(s_i)\cos\varphi(s_i, s_{i+1}) + \lambda(s_{i+1})\cos\varphi'(s_i,s_{i+1}) + O(\epsilon^2),
\end{align*}
or, more concisely, as
\begin{equation}\label{eq:pert_len_func}
    L^\lambda(\sigma) = L(\sigma) + \epsilon P^\lambda(\sigma) + O(\epsilon^2).   
\end{equation}
If \( \sigma \) is a periodic orbit of the unperturbed table \( \mathcal{D} \), then consecutive angles \( \varphi \) and \( \varphi' \) have the same cosine, and hence \( P^\lambda \) becomes
\begin{equation}\label{eq:descrip_Plamb_per_point}
    P^\lambda(\sigma) = 2\sum_{i=0}^{q-1}\lambda(s_i)\cos\varphi(s_i, s_{i+1}),
\end{equation}
with the usual convention that \( s_q = s_0 \). 
Writing \( \sigma^\lambda := (s_0^\lambda, \cdots, s^\lambda_{q-1}) \) to denote the generalised periodic orbit of the perturbed table with same type as \( \sigma \), we have 
\begin{equation*}
    s_i^\lambda = s_i + \epsilon\psi_i + O(\epsilon^2).
\end{equation*}
Thus, \( \sigma^\lambda = \sigma + \epsilon\psi \) satisfies
\begin{equation}\label{eq:per_orb_perturb}
\begin{split}
    0 &= DL^\lambda(\sigma^\lambda) \\
    &= DL(\sigma^\lambda) + \epsilon DP^\lambda(\sigma^\lambda) + O(\epsilon^2) \\
    &= DL(\sigma) + \epsilon D^2L(\sigma)\psi + \epsilon DP^\lambda(\sigma) + O(\epsilon^2) \\
    &= \epsilon D^2L(\sigma)\psi + \epsilon DP^\lambda(\sigma) + O(\epsilon^2),
\end{split}
\end{equation}
where the last equality comes from the fact that \( \sigma \) is a critical point for \( L \).
Thus, the first order change \( \psi \) we see in the collision points for a critical point \( \sigma^\lambda \) is the solution of the linear system
\begin{equation}\label{eq:1st_order_posi}
    D^2L(\sigma)\psi = -DP^\lambda(\sigma).
\end{equation}
We already know that the Hessian of \( L \) is a non-degenerate positive definite matrix, so for a non-positive solution \( \psi \) to \eqref{eq:1st_order_posi} to exist it is sufficient that the term \( -DP^\lambda(\sigma) \) is non-zero.
We conclude that, unless \( \sigma \) is a critical point of \( P^\lambda \), the perturbed flow will have a new generalised periodic orbit arbitrarily close to the original one.
In particular, for a generic choice of perturbations, such perturbed orbits always exist. 
Naturally, if \( \sigma \) is a critical point of \( P^\lambda \), then \( \psi = 0 \) is the unique solution of \eqref{eq:1st_order_posi} and \( \sigma \) remains a generalised periodic orbit after the perturbation.
\begin{defi}
    Given a point \( p = \gamma(s_0) \) in the obstacle \( \Gamma_i \subset \partial\mathcal{D} \), we say a localised perturbation is a \emph{retraction of \( \mathcal{O}_i \) at \( p \)} if it is of the form
    \begin{equation*}
        \widetilde\gamma(s) = \gamma(s) - \epsilon\Psi(s)n_s,
    \end{equation*}
    where \( \epsilon >0 \) and \( \Psi \) is a positive bump function supported at a small interval around \( s_0 \), with \( \Psi(s_0) = 1. \)
\end{defi}
\begin{remark}\label{rem:pertubed_periodic_orbits}
    If a periodic orbit \( \sigma \) has a grazing collision at \( p \), then any retraction of the obstacle tangent to the orbit at \( p \), not matter how small, removes the tangency without changing the original orbit, i.e. the perturbed orbit \( \sigma^\lambda \) has the same trace, but without the grazing collision at \( p \).
    \par On the other hand, a generic perturbation of \( \sigma \) at a regular collision point will generate a non-singular generalised orbit. 
    Indeed, the law of reflection means that a change in the angle of one of the collisions will propagate through all the orbit, so that, generically, a change of the angle and position of one the vertices will either kill the periodic orbit \( \sigma \) by pushing it into an obstacle (i.e., the perturbed orbit \( \sigma^\lambda \) is a ghost orbit) or it will move the orbit away from the obstacles it previously touched tangentially (that is, \( \sigma^\lambda \) is regular). 
\end{remark}
\subsubsection{Killing off singular orbits} As suggested in Remark \ref{rem:pertubed_periodic_orbits}, localised perturbations can ``destroy'' singular periodic orbits by turning them into regular or ghost ones.
We claim this procedure can be applied, in a controlled way, to remove any finite number of singular orbits without producing new ones.
\begin{lemma}\label{lem:fin_hor_persists}
    If \( \mathcal{D} \) has finite horizon, so does every sufficiently close localised perturbation of \( \mathcal{D}. \)   
\end{lemma}
\begin{proof}
    Let \( p \) be a collision point in \( \partial\mathcal{D} \), and let \( L \) be the image of the orbit of \( p \) in the torus.
    If \( p \) is regular, then, for sufficiently small perturbations localised at \( p \), the corresponding perturbed orbit \( L' \) is still regular and bounces at the perturbed obstacle \( \partial\mathcal{D}'_i \).
    Thus, if a perturbation around \( p \) were to create a periodic orbit with no collisions at all, then the collision at \( p \) must be grazing.
    It must also be the only collision on the orbit \( L \); otherwise, a localised perturbation would not eliminate all collisions.
    Hence, the line \( L \) is closed, tangent to the obstacle \( \partial\mathcal{D}_i \) at \( p \), and a positive distance away from all the other curves in \( \partial\mathcal{D} \).
    In particular, this means the obstacle \( \partial\mathcal{D}_i \) lies entirely on one side of \( L \).
    Now, on the side of \( L \) opposite the obstacle, any line parallel to \( L \) and sufficiently close to it would not meet any of the obstacles \( \partial\mathcal{D}_j \), which contradicts our finite horizon hypothesis.
    It follows that localised perturbations of finite horizon tables can not create tables of unbounded horizon, as we wanted.    
\end{proof}
\par Given \( n \in \N \), there exists a positive constant \( M \), depending on \( n \) and on the table \( \mathcal{D} \) (more specifically, on the maximum free-flight time \( \tau_\mathrm{max}) \) with the following property: for every periodic point \( p \in \mathcal{M} \) of order less or equal to \( n \), the fundamental trace of its lift to \( \widetilde{\mathcal{D}} \) is completely contained in \( [-M, M]^2 \).
    \begin{figure}[!ht]
        \centering
        \includegraphics[width=.8\textwidth]{Sinai.pdf}
        \caption{Physical and non-physical trajectories on the lifted table.}
        \label{fig:lifted_table_per}
    \end{figure}
\par We consider the set of all possible symbolic sequences of period \( n \) orbits:
\[
    \Sigma_n \coloneqq \left\lbrace (a_0, \cdots, a_{n-1}) \;\middle|\;
    \begin{tabular}{@{}l@{}}
        \( a_i \in [-M, M] \times \{1, \cdots, k\} \subset \mathbb{A}; \) \\ 
        \( \pi(B_{a_i}) \neq \pi(B_{a_{i+1}}) \text{ for } i = 0, \cdots, n-2 \); \\
        \( \pi(B_{a_{n-1}}) = \pi(B_{a_0}) \).
    \end{tabular}
    \right\rbrace.
\]
\par Now, according to Theorem \ref{thm:bunimovich_periodic}, to each element of the finite set \( \Sigma_n \) corresponds at most one periodic orbit, which can be non-physical.
Whether or not a sequence has an associated periodic trajectory depends on the length functional of the table.
The important point here is that localised perturbations may change the length functional slightly, but they do not create or destroy scatterers, and hence do not change the set \( \Sigma_n \).
Moreover, as a consequence of Equation \eqref{eq:1st_order_posi}, ghost and regular trajectories are stable under perturbations, whereas singular orbits can become either regular or non-physical.
Either way, the sequence associated to a generalised orbit will still be associated to the corresponding perturbed orbit.
 \begin{lemma}\label{lem:remov_tangency}
     Let \( p \in \partial\mathcal{D} \) and \( n \in \N \). 
     Then, for any sufficiently small neighbourhood \( U \subset \partial\mathcal{D} \) of \( p \), there exists localised perturbations of \( \mathcal{D} \) at \( p \) with no singular points of order \( n \) or less inside \( U \). 
 \end{lemma}
 \begin{proof}
    As there are finitely many periodic orbits of order \( n \), the result is straightforward for any point not lying on a periodic orbit.
    Assume \( p \in \Gamma_a, a \in \mathbb{A} \), is a vertex of a periodic trajectory.
    Let 
    \[
        \mathcal{A} \subset \bigcup_{i=2}^n\Sigma_i
    \]
    be the set of all sequences of \( n \) or less elements containing \( a \).
    We further partition \( \mathcal{A} \) into the sets \( \mathcal{A}_{\mathrm{emp}} \), of sequences corresponding to no periodic trajectory; and its complement \( \mathcal{A}_\mathrm{orb}:= \mathcal{A}\setminus \mathcal{A}_{\mathrm{emp}} \), of sequences corresponding to periodic trajectories (physical or not).
    \par As there are finitely many closed orbits of order less than \( n \), there exists a neighbourhood \( U \) of \( p \) which does not contain the vertices of any periodic orbits of order \( n \) or less other than ones reflecting exactly at \( p \).
    Now, the generalised orbits reflecting at \( p \) will all persist under perturbation localised at \( U \), provided it is sufficiently small.    
    In other words, small localised perturbations do not decrease the size of the set \( \mathcal{A}_\mathrm{orb} \).
    They could however, create a new periodic orbit \( O \) associated to one the sequences previously in \( \mathcal{A}_\mathrm{emp} \).
    Since, however, such set is finite, there are only finitely many perturbations which could possible create such news orbits.
    \par Choose the perturbation to be a retraction \( \widetilde\gamma \) of \( \Gamma_a \) at \( p \) with support \( J \) such that \( \widetilde\gamma(J) \subset U \). 
    If any of the orbits containing \( p \) were grazing at \( p \), the grazing collision does not happen any more after the perturbation.
    The regular and ghost orbits reflecting at \( p \) will persist after the perturbation, provided it is small enough.
    Moreover, we can choose our perturbation in such a way that any singular orbit having a regular collision at \( p \) is either regular or non-physical after perturbation.
    Finally, if the perturbation creates any new periodic orbits, that is, if it increases the size of \( \mathcal{A}_\mathrm{orb} \), then we repeat the procedure considering the new orbits and a suitable new neighbourhood \( \widetilde U \subset U \).
    Since \( \#\mathcal{A}_\mathrm{emp} < \infty \), the process terminates after finitely many steps.
\end{proof}
Recall that \( \mathcal{R}_n \) is the subset of \( \mathfrak{D}^m \) consisting of tables without any singular periodic orbits of period less or equal to \( n \). 
\begin{lemma}\label{lem:dense_nonsingular_n}
For every \( n\geq 2 \), the set \( \mathcal{R}_n \) is open and dense in \( \mathfrak{D}^m \). 
\end{lemma}
\begin{proof}
    As there are only finitely many periodic orbits of order less than \( n \), and regular orbits are persistent under sufficiently small perturbations, it follows that any table close enough to a table in \( \mathcal{R}_n \) is also in \( \mathcal{R}_n \).
    We need only show that \( \mathcal{R}_n \) is dense.
    Let \( \mathcal{D} \in \mathfrak{D}^m \).
    Using induction on the order of periodic orbits, we construct an element of \( \mathcal{R}_n \) arbitrarily close to \( \mathcal{D} \). \\
    \par\noindent\emph{The base case \( n=2 \)}. In this somewhat degenerate situation where there are only two obstacles \( \Gamma_1 \) and \( \Gamma_2 \), the only possible singular periodic orbit have to two regular collisions at one scatterer, say \( \Gamma_1 \), and one grazing collision at the point \( p \in \Gamma_2 \).
    Note that the two regular collisions at points \( x, y \in \Gamma_1 \) must be normal collisions, and that convexity of \( \Gamma_1 \) does not allow for the existence of sequences \( x_i, y_i \in \Gamma_1 \) such that the line segments \( [x_i,y_i] \) are realised as billiard trajectories while simultaneously converging to \( [x,y] \) as \( i \to \infty \). 
    Thus, any sufficiently small retraction of \( \mathcal{O}_2 \) at \( p \) would immediately remove the tangency while preserving the orbit. \\
    \par\noindent\emph{The general case}. Assume, by induction, that the result holds for \( \mathcal{R}_{n-1} \).
    We lift our initial table \( \mathcal{D} \) and consider the set \( \Sigma_n \) of all possible encodings of periodic orbits of order \( n \).
    Since there only finitely many periodic orbits, the collection of all reflection points of singular orbits is also finite, and we can choose a minimal finite sub-collection \( \{p_1, \cdots, p_k\} \) containing at least one reflection point of each singular orbit. 
    \par Starting with \( p_1 \), we can use Lemma \ref{lem:remov_tangency} to find an arbitrarily small perturbation \( \mathcal{D}' \) of \( \mathcal{D} \) around \( p_1 \) which removes all the singular orbits of order \( n \) reflecting at \( p_1 \), without creating any new ones.
    Moreover, since the perturbation is localised away from singular orbits not reflecting at \( p_1 \), such orbits comprise the full set of singular orbits of order \( n \) of \( \mathcal{D}' \).
    \par Now, according to Lemma \ref{lem:fin_hor_persists}, \( \mathcal{D}' \) still has finite horizon, so that we can once again repeat the construction done to \( \mathcal{D} \) to encode its periodic orbits using the set \( \mathfrak{D}_n \). 
    Furthermore, the set \( \{p_2, \cdots, p_k\} \) is a minimal sub-collection of reflection points representing all the singular orbits of \( \mathcal{D}' \), and we can repeat the procedure for \( p_2 \).
    \par After finitely many steps we obtain a perturbed table, which we still call \( \mathcal{D}' \), all of whose periodic orbits of order \( n \) are regular.
    In particular, since regular orbits are persistent under perturbation, any table sufficiently close to \( \mathcal{D}' \) also has this property.
    It follows from the induction hypothesis that \( \mathcal{D}' \) can be approximated by a table with no singular orbits of order less than or equal to \( n \).
    As such approximations can be made arbitrarily close, it follows that there are tables in \( \mathcal{R}_n \) arbitrarily close to \( \mathcal{D} \), as we wanted.
\end{proof}

Recall that \( \mathcal{S}_n \) is the subset of \( \mathfrak{D}^m \) consisting of tables such that no distinct periodic orbits of period less or equal to \( n \) have the same length. 


\begin{lemma}\label{lem:dense_simple_n}
For every \( n \geq 2 \), the set \( \mathcal{S}_n \) is open and dense in \( \mathfrak{D}^m \). 
\end{lemma}

\begin{proof}
    From the equality in \eqref{eq:per_orb_perturb} we can see that the first order variation of perimeter from a perturbation of the form \( \gamma^\lambda(s) = \gamma(s) + \epsilon\lambda(s)n_s \) is
    \begin{equation}\label{eq:1st_ord_length}
        \Delta_\sigma(\psi) = P^\lambda(\sigma) + DP^\lambda(\sigma)(\psi),
    \end{equation}
    where \( \sigma = (s_1, \cdots, s_n) \) are the reflection points of the original periodic orbit and the corresponding periodic orbit for the perturbation reflects at points \( s_i^\lambda = s_i + \epsilon\psi_i \).
    In particular, arbitrarily small perturbations induce arbitrarily small changes in length.
    \par Given a table \( \mathcal{D} \) in \( \mathcal{S}_n \), since there only finitely many periodic orbits of order \( n \) or less, we have \( \lvert L(\sigma) - L(\sigma')\rvert > \eta > 0 \) for any two periodic orbits \( \sigma, \sigma' \).
    Any perturbation of \( \mathcal{D} \) small enough not to alter the length of such orbits by more than \( \eta/2 \) would still be an element of \( \mathcal{S}_n \), so \( \mathcal{S}_n \) is open. 
    \par As for density, let \( U_0 \) be an open set in \( \mathfrak{D}^m \), and \( \mathcal{D} \) a table in \( U \).
    Due to the existence of the minimum free-flight time for any table in \( \mathfrak{D}^m \), it is sufficient to show that arbitrarily close to \( \mathcal{D} \) there is a table where no two periodic orbits of the same order have the same length.
    If \( \mathcal{D} \) itself has such a property there is nothing to show, so assume \( O, O' \) are two distinct closed trajectories in \( \mathcal{D} \) with same length.
    Suppose first that \( O \) has a reflection point \( p \) which is not a collision point of \( O' \).
    We choose the perturbation \( \lambda \) localised around \( p \) so that it changes the tangent vector of the obstacle at \( p \), but does not move \( p \) itself. 
    More explicitly, \( \lambda \) satisfies
    \begin{equation}
    \begin{cases}\label{eq:cond_lamb}
        \lambda(s_0) = 0, \\
        \lambda'(s_0) \neq 0.
    \end{cases}
    \end{equation}
    Applying such conditions to Equations \eqref{eq:descrip_Plamb_per_point} and \eqref{eq:1st_ord_length}, and using Equation \eqref{eq:1st_order_posi}, we obtain that
    \begin{equation*}
        \Delta_\sigma(\psi) = D^2L(\sigma)(\psi,\psi) > 0.
    \end{equation*}
    Thus, the increase in length of the perturbed orbit \( \widetilde O \) is non-zero, and now the two orbits have distinct lengths. 
    As this process is localised at \( p \), it does not change the perimeter of any other orbit not reflecting at \( p \).
    Thus, the only situation where such a procedure could fail is when \( p \) is a collision point of more than one orbit (in particular, if the orbits \( O, O' \) share all their reflection points).
    If this happens, the orbits cannot have the same reflection angles, so the corresponding vectors \( DP^\lambda(\sigma) \) and \( DP^\lambda(\sigma') \) are distinct.
    In particular, after a preliminary perturbation of the type \eqref{eq:cond_lamb}, the corresponding orbits now have different reflection points, and the procedure above can be applied.
    \par Finally, as all the perturbations involved can be made arbitrarily small, it follows that \( \mathcal{D} \) can be approximated arbitrarily well by a table in \( \mathcal{S}_n \), and this set is therefore dense.    
\end{proof}

    
\section{Billiard cycles and the enriched length functional}\label{sec_billiards_cycles_lf}
\par As seen in Section \ref{sec:geodesic_approximations}, every Kourganouff surface \(K\), when compressed on its third coordinate, converges to a metric space \(K_0\), whose geometry is determined by the table \(\mathcal{D}\).
This suggests that the enriched marked length spectrum, while depending heavily on the limiting surface \(K_0\), is in fact an object intrinsic to the table \( \mathcal{D} \).
\par In what follows, we describe \( \mathcal{EL}_{\mathcal{D}; K} \) in terms of lengths on the table only.
We begin with the following definition.
\begin{defi}[Billiard cycle]
    A closed curve \( \sigma\colon I \to \mathcal{D} \) is a \emph{billiard cycle} if there is a finite set \( t_0 < t_1 < \cdots < t_{k-1} \) in \( I \) such that
    \begin{itemize}
        \item[(i)] \( \sigma(t_i) \in \partial\mathcal{D} \) for \( 0 \leq i \leq k \);
        \item[(ii)] the segment \( \sigma\rvert_{[t_i, t_{i+1}]} \) is a geodesic of \( d_\mathcal{L} \) for every \( i = 0, \cdots, k, \) where \( t_k = t_0 \);
        \item[(iii)] at each point \( \sigma(t_i) \) the curve undergoes a specular reflection, that is:
        \[
            \dot\sigma_+(t_i) = \dot\sigma_-(t_i) - 2\langle \dot\sigma_-(t_i), n\rangle n,
        \]
        where 
        \( \dot\sigma_s(t_i) = \lim_{t \to t_i^s}\frac{1}{t}(\sigma(t) - \sigma(t_1))\), for \( s\in \{+ , -\} \),
        are the side limits of \(\sigma\) at \(t_i\), \(n\) is the normal to \( \partial\mathcal{D} \) at \(\sigma(t_i)\) and \( \langle \cdot, \cdot \rangle\) is the usual Euclidean dot product.
    \end{itemize}     
\end{defi}
It is clear that every closed orbit of the billiard flow is also a billiard cycle, as well as every closed geodesic of the length metric \( d_\mathcal{L} \).  

\begin{theorem}\label{thm:char_billiard_cycles_1}
    There is a surjection, induced by \( p \), between the set \( \mathcal{C}(K) \) and the collection of all billiard cycles.
\end{theorem}
\begin{proof}
    \par We begin by showing that, for any \( \alpha \in \mathcal{C}(K) \), the geodesic \( \gamma^\alpha_0 \) projects to a billiard cycle.
    We consider an arc-length parametrisation \( \gamma^\alpha_0\colon [0, l_\alpha] \to K_0 \).
    We need to construct a partition \( \{t_0 < \cdots < t_{k-1}\} \) of \( [0, l_\alpha] \) satisfying properties (i), (ii) and (iii).
    \par First, suppose \( \gamma^\alpha_0 \subset B(\gamma^\alpha_0) \).
    In this case, the curve \( \gamma^\alpha_0 \) coincides with one the obstacles \( \Gamma_i \), the trace of which is already a geodesic of \( d_\mathcal{L} \).
    We simply take \( t_0 = 0 \) and it follows that \( p(\gamma^\alpha_0) \) has the desired properties.
    \par Otherwise, the set \( \gamma^\alpha_0\setminus B(\gamma^\alpha_0) = K^\circ\cap\gamma^\alpha_0 \) is a non-empty open subspace of \( \gamma^\alpha_0\).
    Hence there is a finite collection of open subintervals \(I_j \subset [0, l_\alpha] \) such that the decomposition of  \( K^\circ\cap\gamma^\alpha_0 \) into connected components is
    \[
        K^\circ\cap\gamma^\alpha_0 = \coprod_{j=0}^{q}\gamma^\alpha_0(I_j)
    \]
    We assume our intervals are labelled in order, such that \( \sup I_j < \inf I_{j+1} \).
    Let \( t_j := \sup I_j \). 
    By construction \( \gamma^\alpha_0(t_j) \) is a collision point and, for any \( t \in I_j \), the triple \( (t_j-t; p(\gamma^\alpha_0(t))) \) is an element of \( A_0 \).
    It follows that each \( \gamma^\alpha_0\rvert_{I_j} \) projects to a straight line on \( \mathcal{D} \).
    Moreover, if the collision at \(t_j\) is non-degenerate, then \(\gamma^\alpha_0(t_j)\) is an isolated element of \(B(\gamma^\alpha_0)\), hence \(\inf I_{j+1} = \sup I_j = t_j\).
    \par Thus, at a non-degenerate collision \( \gamma^\alpha_0(t_j) \) the set
    \[
        p(\gamma^\alpha_0\rvert_{I_j})\cup p(\gamma^\alpha_0\rvert_{I_{j+1}})
    \]
    is the union of two geodesic segments in \( (\mathcal{D}, d_\mathcal{L}) \) with a common end point. If the collision at \(t_j\) is grazing, then \( p(\gamma^\alpha_0(t)\rvert_{I_j \cup I_{j+1}}) \) is actually a straight line. 
    On the other hand, if it is regular, then the geodesic \( \gamma^\alpha_0 \) is still conjugated to the geodesic flow after the collision, i.e., \(\inf I_{j+1} = t_j \) and \( (t_{j+1}-t; p(\gamma^\alpha_0(t))) \in A_0 \) for any \( t\in I_j \).
    In particular, \( \gamma^\alpha_0(t)\rvert_{I_j \cup I_{j+1}} \) projects to a billiard trajectory.
    In any case, we see that when the collision point at \(t_j\) is non-degenerate, the trajectory respect the desired specular reflection law at \( p(\gamma^\alpha_0(t)\rvert_{I_j \cup I_{j+1}}) \) at \(t_j\).
    \par If, on the other hand, the collision point at \(t_j\) is degenerate, then \(s_{j}:=\inf I_{j+1} > t_j \) and \(p(\gamma^\alpha_0\rvert_{[t_j, s_j]})\) is a segment of the boundary \( \partial \mathcal{D} \). 
    Now, we notice that in this case \(\gamma^\alpha_0(I_j \cup [t_j, s_j] \sup I_{j+1}\) is entirely contained in the closure of a connected component of \( K^\circ \), hence  \(p\) maps it into a geodesic of \( (\mathcal{D}, d_\mathcal{L}). \)   
    \par So we proceed as follows: the first element of our partition is \( t_0 \). 
    If \( t_1 \) is a regular collision point, we add it to the partition.
    If not, we look at the next regular collision point \( t_j \). 
    The segments between regular collisions points are unions of straight lines and boundary segments, and such unions are geodesics because they are images under \( p \) of segments of \( \gamma^\alpha_0 \). 
    As regular collision points obey the specular reflection law, the collection of all of them forms the desired partition.  
    This finishes the first part of the proof.
    \par Conversely, we show that every billiard cycle of \( \mathcal{D} \) is the image under \( p\) of a geodesic \( \gamma^\alpha_0 \).
    Let \( \sigma \) be a billiard cycle, and \( x_i = \sigma(t_i), \ i = 0, 1, \cdots, k-1, k \) its regular collision points.
    As before, we let \( t_k = t_0 \).
    We lift \( \sigma \) to a curve \( \sigma_\mathrm{up} \) in \( K_0 \) in the following way.
    For each \( x_i \in \partial\mathcal{D}\) there is a single point \( \widetilde x_i \in B(K) \) which projects to \( x_i \), since \( p\rvert_{B(K)} \) is a diffeomorphism.
    We lift the first geodesic segment \( \sigma\rvert_{[t_0, t_1]} \) to the sheet \( \mathcal{D}_\mathrm{up} \).
    This can be done in a unique way, since \( p\rvert_{\mathcal{D}_\mathrm{up}}: (\mathcal{D}_\mathrm{up}, d_0) \to (\mathcal{D}, d_\mathcal{L}) \) is an isometry.
    Now, in the same way, we lift the second geodesic segment \( \sigma\rvert_{[t_1, t_2]} \), this time to the sheet \( \mathcal{D}_\mathrm{down} \).
    As we showed in Theorem \ref{thm:K_0_as_gluing_metric_space}, if geodesic segments from different sheets are joined a collision point, then the curve obtained is minimising exactly when the collision is specular.
    Hence the lifting of \( \sigma\rvert_{[t_0, t_2]} \) thus defined is a geodesic segment of \( K_0 \).
    \par If we proceed this way, we get a closed curve \( \sigma_\mathrm{up} \) in \(K_0 \) which is everywhere minimising, except maybe at \( t_0 \).
    Indeed, if the lifts \( \sigma\rvert_{[t_0, t_1]} \) and \( \sigma\rvert_{[t_0, t_2]} \) are in the same sheet\footnote{
        this happens exactly when we have an even number \( k \) of regular collision times \( t_i \), or equivalently an odd number of collision points \( x_i \)
    },
    then \( \sigma\rvert_{[t_k, t_1]} \) is not minimising around \( t_0 \):
    convexity of the scatterers means that for any two points \( \sigma(t_0-\epsilon) \) and \( \sigma(t_0 + \epsilon) \) the line
    \[
        r(t) = \sigma(t_0-\epsilon) + t(\sigma(t_0+\epsilon) + \sigma(t_0-\epsilon)) 
    \]
    is entirely contained in the interior of \( \mathcal{D} \) and has Euclidean length smaller than the sum \( l_0(\sigma\rvert_{[t_0-\epsilon, t_0]}) + l_0(\sigma\rvert_{[t_0, t_0+\epsilon]}) \).
    To remedy this, we can consider the curve \( \sigma_\mathrm{down} \), which is the lift of \( \sigma \) to \( K_0 \) constructed as before, but starting at the sheet \( \mathcal{D}_\mathrm{down} \).
    Then the curve \( \sigma_\mathrm{up}\ast\sigma_\mathrm{down}^{-1} \), by construction, is a concatenation of geodesic segments, all of whose regular collisions are specular, and is therefore a geodesic of \( K_0 \).
    \par Finally, we let \( \alpha \) be the free homotopy class of \( \sigma_\mathrm{up}\ast\sigma_\mathrm{down}^{-1} \). 
    Since \( \sigma_\mathrm{up}\ast\sigma_\mathrm{down}^{-1} \) is a closed geodesic, it follows from Lemma \ref{lem:unique_minimising_geodesic} that \(\gamma^\alpha_0 = \sigma_\mathrm{up}\ast\sigma_\mathrm{down}^{-1} \).
    Thus \(p(\gamma^\alpha_0) = \sigma \), as we wanted.  
\end{proof}


\subsubsection{The enriched length functional.}
\par In this section we think of \( \mathcal{D} \) as periodic billiard table on \( \mathbb{R}^2 \), and consequently \( K \) is a periodic surface in \( \mathbb{R}^3 \). 
Our goal is to describe closed cycles as critical points of a length functional, much like in Appendix \ref{sec_periodicorbitsbilliard}.
To that end, we keep the same notation:
given a periodic sequence of scatterers \( \Gamma_{\rho_0}, \cdots , \Gamma_{\rho_{q-1}}, \  \rho_k \in \{ 1,\cdots, m \} \), for \( k \in \{ 0,\cdots, q-1 \} \). 
We let \( \rho = \rho(\sigma):=(\rho_k)_{k=0,\cdots,q-1} \). 
Let us consider the lift to \( \widetilde{\mathcal{D}} \) of a cycle \( \sigma \) starting in the cell \( (0,0) \), and denote by  \( \{(\tilde i_k,\tilde j_k)\}_{k=0,\cdots,q} \) the labels of the cells where the first \( (q+1) \) bounces happen, with \( (\tilde i_0,\tilde j_0)=(0,0) \).
For \( k\in \{1,\cdots,q\} \), set \( I_k:=(\tilde i_k,\tilde j_k)-(\tilde i_{k-1},\tilde j_{k-1}) \). 
We denote \( I = I(\sigma):=(I_k)_{k=1,\cdots,q} \).
The extended length functional \( \mathrm{EL} \) acts on a subset of \( \mathbb{T}_{\rho_0}^2 \times \mathbb{T}_{\rho_1}^2 \times \cdots \times \mathbb{T}_{\rho_{q-1}}^2 \).

\begin{defi}[Link sets] Given a pair of indices \( (i, j) \in \{1, \cdots, m\}^2 \), we define the \emph{link set} \( L_{ij} \) of the scatterers \( \Gamma_i, \ \Gamma_j \) by
\[
	L_{ij} := \{ y \in \Gamma_j  ; \ \exists \  x \in \Gamma_i \text{ such that } [x,y] \subset \mathrm{int}(\mathcal{D}) \}
\]
where \(  [x,y] \) is the line segment joining \( x \) to \( y \).
\end{defi}
Observe that these sets are closed intervals in the scatterers, the borders consisting of points whose trajectory (in a given direction), hits the next scatterer at (or after) a grazing collision.
\par We let
\[
D_i\colon \T^2_i \to \R, \ (e,s) = \min\{|e-s|, l_i-|e-s|\},
\]
that is, \( D_i(e,s)\) is the length of the smallest arc in \( \Gamma_i \) joining \(\gamma_i(e) \) to \( \gamma_i(s). \)
The function \( D_i \) is differentiable almost everywhere. 
Indeed, if we write 
\[
\begin{split}
    I_i(e) &:= [(e-l_i/2, e) \cup (e+l/2, \infty)] \cap [0,l_i]; \\
    J_i(e) &:= [(-\infty, e-l_i/2) \cup (e, e+l/2)] \cap [0,l_i].
\end{split}
\]
then we have \( I_i(e) \cup J_i(e) = [0, l_i]\setminus\{e, e \pm l_i/2\} \), and moreover:
\begin{equation}\label{eq:enriched_length_curved_part}
\begin{split}
    \partial_1D_i(e,s) &= \left\{
    \begin{array}{lcl} 
            1, &\text{ if } s \in I_i(e), \\
            -1, &\text{ if } s \in J_i(e), \\
            \text{undefined,} &\text{if } s \in \{e, e \pm l_i/2\};             
        \end{array}
        \right. \\
    \partial_2D_i(e,s) & = \left\{
    \begin{array}{lcl} 
            1, &\text{ if } s \in J_i(e), \\
            -1, &\text{ if } s \in I_i(e), \\
            \text{undefined,} &\text{if } s \in \{e, e \pm l_i/2\}.            
        \end{array}
        \right. 
\end{split}
\end{equation}

The enriched length functional measures the \(d_\mathcal{L}\) length of piecewise smooth curves consisting of concatenations of line and arc segments along the boundary.
\begin{defi}[Enriched length functional]
Let \( \mathbb{L}^I_\rho\) be the following product of link sets.
\[
   \mathbb{L}^I_\rho := L_{\rho_{q-1}\rho_{0}} \times L_{\rho_1\rho_0} \times L_{\rho_0\rho_1} \times L_{\rho_2\rho_1} \times \cdots \times L_{\rho_{q-2}\rho_{q-1}} \times L_{\rho_0\rho_{q-1}}.
\]
Let \( S =  (e_{0}, s_0, e_1, s_1, \cdots, e_{q-1}, s_{q-1}) \) be a point in \( \mathbb{L} \).
Then the enriched length functional associated to the sequence \( \rho \) is
\[
	\mathrm{EL}^{I}_\rho \colon \left\{
        \begin{array}{lcl} 
            \mathbb{L} \to \mathbb{R} \\
	       S \mapsto      {\displaystyle{\sum_{k=1}^{q}}} \tau_{I_k}(s_{k - 1},e_{k}) + D_{\rho_k}(e_{k}, s_{k}) 
            =
            L^I_\rho(s_0, e_1, \cdots, s_{q-1}, e_q) + {\displaystyle{\sum_{k=1}^{q}}} D_{\rho_k}(e_{k}, s_{k}).
        \end{array}
        \right.
\]
where we make the convention \( (e_q, s_q) = (e_0, s_0). \)
\end{defi}
Remark that if \(s =  (s_0, \cdots, s_{q-1}) \) is a periodic trajectory of the table, that is, a minimum of \( L^I_\rho, \),  then \(L^I_\rho(s) = \mathrm{EL}^I_\rho(S) \), where \(S =  (s_0, s_0, s_1, s_1, \cdots, s_{q-1}, s_{q_1}) \).
In particular, the point \(S\) is also a minimum of \( \mathrm{EL}^I_\rho \).
More generally, we have the following.
\begin{prop}
    The billiard cycles of \( \mathcal{D} \) are minima of the length functionals \( \mathrm{EL} \).
\end{prop}
\begin{proof}
Remark that \( \mathbb{L} \) is a convex subset of \( \mathbb{T}_{\rho_0}^2 \times \mathbb{T}_{\rho_1}^2 \times \cdots \times \mathbb{T}_{\rho_{q-1}}^2 \), which in turn implies that \( \mathrm{EL}^I_\rho \) is convex. 
Indeed, it is the restriction to a convex set of the sum of convex functions.
In particular, the Minimum Theorem applies, and a point \( S_0 \) is a minimum of \( \mathrm{EL}^I_\rho \) if and only if \( 0 \) belongs to the subdifferential \( \partial\mathrm{EL}^I_\rho(S_0) \).
Now, since \( \mathrm{EL}^I_\rho \) is a sum, it follows that
\[
    \partial L^I_\rho(S) + \partial D_{\rho_0}(S) + \cdots + \partial {\rho_{q-1}}(S) \subset \partial\mathrm{EL}^I_\rho(S),
\]
for every point \( S \).
\par Where a function is differentiable its subdifferential is a singleton consisting of the gradient, so it is enough to look at points where \( D_{\rho_i} \) has no derivative.
At these points, the partial derivatives exist, and therefore the subdifferential is obtained from the convex hull of sided limits, meaning that
\[
    \partial D_{\rho_i} (S) = [-1, 1]^2.
\]
Hence, \( 0 \in \partial D_{\rho_i}(S) + \cdots + \partial {\rho_{q-1}}(S) \) whenever \( S \) is a discontinuity point of \( D_{\rho_i} \).
\par Given a billiard cycle \( \sigma \), we may represent it as a point in the \( \mathbb{L}^I_\rho \) space by setting
\[
    S_\sigma = (e_0, s_0, \cdots, e_{q-1}, s_{q-1})
\]
where \( (e_i, s_i) \) are the collision points. 
For a non-degenerate collision, \( e_i = s_i \), while for degenerate collisions we think of \(e_i\) and \(s_i\) as the ``entry'' and ``exit'' points of the cycle in the scatterer \( \Gamma_{\rho_i} \).
\par Now, for a cycle \(\sigma\), either every collision is specular, and therefore \( S_\sigma \) is a critical point of \( L^I_\rho \) and a point of non-differentiability for every \(D_{\rho_i}\). 
Hence
\[
    \partial L^I_\rho(S_\sigma) = \{\nabla L^I_\rho(S_\sigma)\} = \{0\},
\]
and consequently \( 0 \in \partial \mathrm{EL}^I_\rho. \)
On the other hand, at degenerate collisions, the angle of entry (and exit) must be orthogonal to the normal direction, meaning the partial derivatives of \( E^I_\rho \) at that point will be in the set \( \{-1, 1\} \).
Comparing Equations \eqref{eq:first_gen_fct} and \eqref{eq:enriched_length_curved_part} above, we see that that \( 0 \in \partial \mathrm{EL}^I_\rho(S_\sigma) \) at such collisions as well.
\end{proof}

%%% Biblio
\begin{thebibliography}{10}

\bibitem{Abb_Mazz}
Alberto Abbondandolo and Marco Mazzucchelli.
\newblock Marked length spectrum of geodesic flows on the two-sphere, in preparation.

\bibitem{AbreschMeyer1997}
Uwe Abresch and Wolfgang Meyer.
\newblock Injectivity radius estimates and sphere theorems.
\newblock In {\em Comparison Geometry}, volume~30, pages 1--47. MSRI Publications, 1997.

\bibitem{andersson1977propagation}
Karl Andersson and Richard Melrose.
\newblock The propagation of singularities along gliding rays.
\newblock {\em Inventiones mathematicae}, 41(3):197--232, 1977.

\bibitem{baladi2018exponential}
Viviane Baladi, Mark Demers, and Carlangelo Liverani.
\newblock Exponential decay of correlations for finite horizon sinai billiard flows.
\newblock {\em Inventiones mathematicae}, 211(1):39--177, 2018.

\bibitem{ballmann1995lectures}
Walter Ballmann.
\newblock {\em Lectures on Spaces of Nonpositive Curvature}.
\newblock Oberwolfach Seminars. Birkh{\"a}user Basel, 1995.

\bibitem{bankovic_marked-length-spectral_2017}
Anja Bankovic and Christopher Leininger.
\newblock Marked-length-spectral rigidity for flat metrics.
\newblock {\em Transactions of the American Mathematical Society}, 370(3):1867--1884, October 2017.

\bibitem{bonahon_bouts_1986}
Francis Bonahon.
\newblock Bouts des {Variétés} {Hyperboliques} de {Dimension} 3.
\newblock {\em Annals of Mathematics}, 124(1):71--158, 1986.

\bibitem{bourdon_sur_1996}
Marc Bourdon.
\newblock Sur le birapport au bord des {CAT}(— 1)-espaces.
\newblock {\em Publications Mathématiques de l'Institut des Hautes Études Scientifiques}, 83(1):95--104, December 1996.

\bibitem{bridson2011metric}
Martin Bridson and Andr\'e H{\"a}fliger.
\newblock {\em Metric Spaces of Non-Positive Curvature}.
\newblock Grundlehren der mathematischen Wissenschaften. Springer Berlin Heidelberg, 2011.

\bibitem{bunimovich_variational_1995}
Leonid Bunimovich.
\newblock Variational principle for periodic trajectories of hyperbolic billiards.
\newblock {\em Chaos: An Interdisciplinary Journal of Nonlinear Science}, 5(2):349--355, June 1995.

\bibitem{burago2022course}
Dmitri Burago, Yuri Burago, and Sergei Ivanov.
\newblock {\em A Course in Metric Geometry}.
\newblock Graduate Studies in Mathematics. American Mathematical Society, 2022.

\bibitem{calderon2018hear}
Aaron Calderon, Solly Coles, Diana Davis, Justin Lanier, and Andre Oliveira.
\newblock How to hear the shape of a billiard table.
\newblock {\em arXiv preprint arXiv:1806.09644}, 2018.

\bibitem{chernov_chaotic_2006}
Nikolai Chernov and Roberto Markarian.
\newblock {\em Chaotic {Billiards}}.
\newblock American Mathematical Soc., 2006.

\bibitem{croke1990rigidity}
Christopher Croke.
\newblock Rigidity for surfaces of non-positive curvature.
\newblock {\em Commentarii Mathematici Helvetici}, 65(1):150--169, 1990.

\bibitem{de2023marked}
Jacopo De~Simoi, Vadim Kaloshin, and Martin Leguil.
\newblock Marked length spectral determination of analytic chaotic billiards with axial symmetries.
\newblock {\em Inventiones mathematicae}, 233(2):829--901, 2023.

\bibitem{desimoikal}
Jacopo De~Simoi, Vadim Kaloshin, Qiaoling Wei, and Hamed Hezari.
\newblock Dynamical spectral rigidity among $\mathbb{Z}_2$-symmetric strictly convex domains close to a circle (appendix b coauthored with h. hezari).
\newblock {\em Annals of Mathematics}, 186(1):277--314, 2017.

\bibitem{duchin2021you}
Moon Duchin, Viveka Erlandsson, Christopher~J Leininger, and Chandrika Sadanand.
\newblock You can hear the shape of a billiard table: Symbolic dynamics and rigidity for flat surfaces.
\newblock {\em Commentarii Mathematici Helvetici}, 96(3):421--463, 2021.

\bibitem{fisher2019hyperbolic}
Todd Fisher and Boris Hasselblatt.
\newblock {\em Hyperbolic Flows}.
\newblock Zurich lectures in advanced mathematics. European Mathematical Society, 2019.

\bibitem{furman_coarse-geometric_2002}
Alex Furman.
\newblock Coarse-{Geometric} {Perspective} on {Negatively} {Curved} {Manifolds} and {Groups}.
\newblock In Marc Burger and Alessandra Iozzi, editors, {\em Rigidity in {Dynamics} and {Geometry}: {Contributions} from the {Programme} {Ergodic} {Theory}, {Geometric} {Rigidity} and {Number} {Theory}, {Isaac} {Newton} {Institute} for the {Mathematical} {Sciences} {Cambridge}, {United} {Kingdom}, 5 {January} – 7 {July} 2000}, pages 149--166. Springer Berlin Heidelberg, Berlin, Heidelberg, 2002.

\bibitem{ghys_flots_1984}
Etienne Ghys.
\newblock Flots d'{Anosov} sur les 3-variétés fibrées en cercles.
\newblock {\em Ergodic Theory and Dynamical Systems}, 4(1):67--80, March 1984.

\bibitem{guillarmou_geodesic_2022}
Colin Guillarmou, Gerhard Knieper, and Thibault Lefeuvre.
\newblock Geodesic stretch, pressure metric and marked length spectrum rigidity.
\newblock {\em Ergodic Theory and Dynamical Systems}, 42(3):974--1022, March 2022.

\bibitem{guillarmou_marked_2019}
Colin Guillarmou and Thibault Lefeuvre.
\newblock The marked length spectrum of {Anosov} manifolds.
\newblock {\em Annals of Mathematics}, 190(1):321--344, July 2019.
\newblock Publisher: Department of Mathematics of Princeton University.

\bibitem{guillarmou2025marked}
Colin Guillarmou, Thibault Lefeuvre, and Gabriel~P Paternain.
\newblock Marked length spectrum rigidity for anosov surfaces.
\newblock {\em Duke Mathematical Journal}, 174(1):131--157, 2025.

\bibitem{hirsch2012differential}
Morris Hirsch.
\newblock {\em Differential Topology}.
\newblock Graduate Texts in Mathematics. Springer New York, 2012.

\bibitem{klingenberg_riemannian_1974}
Wilhelm Klingenberg.
\newblock Riemannian {Manifolds} {With} {Geodesic} {Flow} of {Anosov} {Type}.
\newblock {\em Annals of Mathematics}, 99(1):1--13, 1974.
\newblock Publisher: [Annals of Mathematics, Trustees of Princeton University on Behalf of the Annals of Mathematics, Mathematics Department, Princeton University].

\bibitem{knieper_new_2012}
Gerhard Knieper.
\newblock New results on noncompact harmonic manifolds.
\newblock {\em Commentarii Mathematici Helvetici}, 87(3):669--703, July 2012.

\bibitem{kourganoff_anosov_2016}
Micka{\"e}l Kourganoff.
\newblock Anosov {Geodesic} {Flows}, {Billiards} and {Linkages}.
\newblock {\em Communications in Mathematical Physics}, 344(3):831--856, June 2016.

\bibitem{lee2019introduction}
John Lee.
\newblock {\em Introduction to Riemannian Manifolds}.
\newblock Graduate Texts in Mathematics. Springer International Publishing, 2019.

\bibitem{otal_spectre_1990}
Jean-Pierre Otal.
\newblock Le {Spectre} {Marque} des {Longueurs} des {Surfaces} a {Courbure} {Negative}.
\newblock {\em Annals of Mathematics}, 131(1):151--162, 1990.

\bibitem{otal_sur_1992}
Jean-Pierre Otal.
\newblock Sur la géometrie symplectique de l'espace des géodésiques d'une variété à courbure négative.
\newblock {\em Revista Matemática Iberoamericana}, 8(3):441--456, December 1992.

\bibitem{walter2013ordinary}
Wolfgang Walter.
\newblock {\em Ordinary Differential Equations}.
\newblock Graduate Texts in Mathematics. Springer New York, 2013.

\bibitem{wilkinson_anosov}
Amie Wilkinson.
\newblock Lectures on marked length spectrum rigidity.
\newblock Lecture notes from the IAS/PCMI Summer School, 2013, 2013.

\bibitem{wojtkowski_two_1994}
Maciej Wojtkowski.
\newblock Two applications of {Jacobi} fields to the billiard ball problem.
\newblock {\em Journal of Differential Geometry}, 40(1):155--164, January 1994.
\newblock Publisher: Lehigh University.

\end{thebibliography}







\end{document}
